\documentclass[11pt]{article}
\usepackage{tcee}
\addbibresource{1D01.bib}
\tcetema[Guía del dictamen]{Primer ejercicio}{Guía para el dictamen de coyuntura económica}
\tceeguia
\tcefecha{23/10/2026}
\begin{document}
\tceetitulo

\section*{Introducción}

\begin{itemize}
\item \textbf{\emph{Para qué sirve}}:
  \begin{itemize}
  \item La segunda parte del primer ejercicio es un \concepto{dictamen de coyuntura}: con un cuadro macroeconómico y unos gráficos de un país, casi siempre anónimo, hay que contestar por escrito a 5 o 6 preguntas en 1 hora y 45 minutos \parencite{BOE2025ConvocatoriaTCEE}.
  \item El \concepto{análisis de coyuntura} estudia la situación actual de una economía, los factores que la han llevado hasta ahí y su evolución probable a corto plazo. No basta con describir qué sube y qué baja: hay que \emph{interpretar} la evolución de los agregados, proponer relaciones causales con prudencia y decir qué importa en ese país.
  \item Un análisis completo tiene 3 partes:
    \begin{enumerate}
    \item una idea de conjunto del estado de la economía, con sus fortalezas y debilidades;
    \item los factores que la han llevado a esa situación, para lo que hay que dominar las relaciones entre las macromagnitudes (y saber que no siempre se cumplen ni se ven en todos los países);
    \item un esbozo de la evolución futura y de las políticas que corregirían los desequilibrios.
    \end{enumerate}
  \item «Lo bueno, si breve, dos veces bueno», escribió Baltasar \autor{Gracián}. El dictamen no se valora al peso: el tribunal pide capacidad de análisis y de síntesis, y en 105 minutos solo cabe lo importante.
  \end{itemize}

\item \textbf{\emph{Cómo está organizada}}:
  \begin{itemize}
  \item El capítulo~\ref{sec:metodo} es el \emph{método}: qué es el ejercicio, qué valora el tribunal, cómo repartir el tiempo, cómo leer el cuadernillo, cómo redactar y qué errores evitar.
  \item Los capítulos siguientes siguen el orden en que suelen preguntarse las cosas: el perfil cíclico (capítulo~\ref{sec:perfil}), los mercados de factores, la competitividad, el sector exterior, las políticas y las recomendaciones.
  \item Cada bloque empieza con un recuadro \emph{Qué hay que contestar}, y cada indicador sigue el mismo esquema: \emph{qué mide}, \emph{cómo se calcula}, \emph{umbral o referencia}, \emph{cómo se lee y con qué se cruza} y el tema del temario donde se estudia a fondo.
  \item Los anexos recogen lo que se consulta y no se memoriza entero: la tabla de indicadores y umbrales, lo que se ha preguntado en los dictámenes oficiales desde 2011, expresiones útiles para redactar, la lectura de series, los cálculos tipo y la programación financiera del FMI.
  \end{itemize}

\item \textbf{\emph{Cómo usarla}}:
  \begin{itemize}
  \item La guía no desarrolla modelos: dice qué mirar, cómo leerlo y dónde está la teoría en el temario (\vertema{3.A.30} para la contabilidad nacional, \vertema{3.A.42} para el ciclo, etc.).
  \item Se aprende con dictámenes resueltos: primero el método y el bloque del perfil cíclico, después un dictamen oficial completo con el reloj delante, y luego los demás bloques a medida que salen en los exámenes.
  \item Las cifras de referencia (umbrales) son de 2 tipos y la guía los distingue siempre: los \emph{oficiales} (reglas fiscales, Basilea~III, cuadro de indicadores de la Comisión), con su fuente, y los \emph{orientativos} (reglas prácticas sin norma detrás), que sirven para ordenar el análisis y nunca para afirmar nada con rotundidad.
  \end{itemize}
\end{itemize}
\nocite{VichGadella2007Guia,GuiaDictamenJLCT,Sahuquillo2023,CuadradoRoura2015,FMI2013CaseStudy,FMI2022edXFPP,Afi2022Jornada}

% ==========================================================================
\section{El ejercicio y el método}\label{sec:metodo}

\begin{itemize}
\item Este capítulo da la conclusión por adelantado: el dictamen se gana \emph{antes} de escribir, en los 10-12 primeros minutos de lectura, y \emph{mientras} se escribe, con un reparto del tiempo decidido en casa y respetado en el examen.
\end{itemize}

\subsection{Qué es el ejercicio y cómo se califica}\label{sec:ejercicio}

\begin{itemize}
\item \emph{Qué dice la convocatoria}: la segunda parte del primer ejercicio «será de carácter práctico y consistirá en la elaboración por escrito de un dictamen sobre un escenario macroeconómico con las informaciones o datos que hayan sido suministrados por el Tribunal. La duración máxima de esta segunda parte será de una hora y cuarenta y cinco minutos» \parencite{BOE2025ConvocatoriaTCEE}. El texto es el mismo que el de la convocatoria anterior \parencite{BOE2024ConvocatoriaTCEE}.
\item \emph{Qué trae el cuadernillo}:
  \begin{itemize}
  \item un cuadro macroeconómico de unos 8-12 años (a veces con previsiones), con el PIB y sus componentes, el mercado de trabajo, los precios, el sector exterior, las cuentas públicas y los tipos de interés;
  \item gráficos (indicadores de alta frecuencia, tipo de cambio, balance del banco central, curva de \autor{Beveridge}, aranceles…);
  \item a veces, unas líneas de contexto: régimen cambiario, objetivo de inflación, pertenencia al euro o al G-20, un programa del FMI, un sector dominante;
  \item 5 o 6 preguntas con subpreguntas (desde 2023).
  \end{itemize}
\item \emph{Cómo se escribe}: a mano, en hojas autocopiativas, con bolígrafo azul o negro y sin calculadora. En 2026 solo se admitían las abreviaturas del enunciado y las convencionales del análisis económico (PIB, pp, pb, FMI); las demás se escriben completas.
\item \emph{La lectura}: quienes superan el test leen el dictamen en sesión pública y el tribunal puede hacer preguntas «durante un período máximo de quince minutos» \parencite{BOE2025ConvocatoriaTCEE}.
\item \emph{Calificación}: el dictamen se puntúa de 0 a 10 y hay que obtener al menos un 4; el primer ejercicio exige, además, 12 de los 30 puntos del conjunto (test, 0-20; dictamen, 0-10). Solo leen el dictamen quienes superan en el test la nota que fije el tribunal, que no puede ser inferior a 8 \parencite{BOE2025ConvocatoriaTCEE}.
\item \emph{El formato cambia}: en 15 años el tiempo ha pasado de 2 horas a 2 horas y 15 minutos y luego a 1 hora y 45 minutos, y el número de preguntas, de 7-10 a 5-6 (tabla~\ref{tab:formato}). La guía, por eso, da reglas que valen para cualquier formato y no un horario cerrado.
\end{itemize}

\begin{table}[H]
\centering\small
\caption{Formato del dictamen en las convocatorias de 2011 a 2026}\label{tab:formato}
\begin{tabular}{>{\raggedright\arraybackslash}p{2.6cm}>{\centering\arraybackslash}p{1.9cm}>{\centering\arraybackslash}p{1.6cm}>{\centering\arraybackslash}p{2.2cm}>{\raggedright\arraybackslash}p{5.2cm}}
\toprule
\rowcolor{tceeazulpalido}
\textbf{Examen} & \textbf{Tiempo} & \textbf{Preguntas} & \textbf{Puntos por pregunta} & \textbf{Rasgos} \\
\midrule
\rowcolor{tceesalvia}
2023, 2024, 2025, 2026 & 1 h 45 min & 5-6 & no & sin cálculos; subpreguntas; en 2026, abreviaturas solo del enunciado \\
2022 & 2 h 15 min & 7 & sí (1-2,5) & programa del FMI; solo PIB e IPC \\
2021 & 2 h 15 min & 6 & sí (1-2) & país del G-20 con flotación sucia \\
2020 & 2 h 15 min & 7 & no & economía abierta agroalimentaria \\
2019 & 2 h 15 min & 10 & sí (1 cada una) & cálculos de CLU y esfuerzo estructural \\
2018 & 2 h & 7 & sí (1-2) & proyecciones y brecha de producción \\
2016 (enero y noviembre) & 2 h 15 min & 8 & no & cálculos; vivienda, petróleo, PMI \\
2011, 2013, 2014 & 2 h & 3-6 & no & calculadora solo en 2013 \\
\bottomrule
\end{tabular}
\fuente{Elaboración propia a partir de los enunciados oficiales \parencite{ATCEE2026Examenes}.}
\end{table}

\begin{notaopositor}
El dictamen se escribe justo después del test, sin descanso y con cansancio; conviene entrenarlo también así, con un test y a continuación un dictamen completo con el reloj delante. Lo que no se ha decidido en casa (el orden de las preguntas, el reparto del tiempo, la estructura de cada bloque) no se decide bien en el examen.
\end{notaopositor}

\subsection{Qué valora el tribunal}\label{sec:que-valora}

\begin{itemize}
\item \emph{Lo dicen las instrucciones}. En 2023, 2024 y 2025, «la capacidad de síntesis y la capacidad de repartir adecuadamente el tiempo proporcionado entre las distintas preguntas»; en 2026, «fundamentalmente las capacidades analítica y de síntesis del aspirante, así como su claridad expositiva» \parencite{ATCEE2026Examenes}.
\item \emph{El objetivo} es demostrar que se domina el análisis básico de una economía: entender un cuadro macroeconómico, las relaciones entre las macromagnitudes, el efecto de las políticas y la acumulación de riesgos.
\item De ahí salen 4 capacidades:
  \begin{enumerate}
  \item \concepto{Capacidad de análisis}: interpretar los datos, no describirlos. Que el consumo crece un 3~\% lo ve cualquiera; lo que se valora es por qué (renta disponible, empleo, tipos, confianza) y qué implica. El examen es un puzle que siempre es parecido: conviene llevar memorizada la estructura de cada pregunta y saber dónde buscar las fortalezas y las debilidades.
  \item \concepto{Capacidad de síntesis}: lo importante primero, con una frase de diagnóstico al empezar cada respuesta (apartado~\ref{sec:redactar}), y una visión de conjunto de la serie, no año a año.
  \item \concepto{Comprensión lectora}: contestar a todo lo que se pregunta, cada subpregunta con su rótulo, y solo a eso.
  \item \concepto{Demostrar conocimiento}: usar el temario cuando viene a cuento (una teoría, una regla fiscal, un indicador) y explicar cómo se ha obtenido cada cálculo o por qué el indicador elegido es el relevante, sin convertir la respuesta en un tema.
  \end{enumerate}
\item Y 2 principios que ordenan todo el análisis:
  \begin{itemize}
  \item \emph{Prudencia}: se pide mucho para el tiempo y los datos disponibles. Si una relación no se aprecia, quizá no está en el cuadro. Cuando cuesta contestar (demasiada información o demasiado poca), un análisis básico y prudente, sin supuestos innecesarios, no se penaliza.
  \item \emph{Coherencia entre las respuestas}: lo que se dice del ciclo en la primera pregunta tiene que cuadrar con lo que se dice después de la inflación, la política fiscal o la monetaria.
  \end{itemize}
\item La solución no es única: se acepta cualquier respuesta plausible, razonada, no tajante y con sentido económico. Hay años en que cambian el tiempo o el tipo de preguntas, y se valora la flexibilidad para centrarse en los retos principales del país.
\end{itemize}

\subsection{El reparto del tiempo}\label{sec:tiempo}

\begin{itemize}
\item La idea: el reparto se decide antes del examen con una regla sencilla, que vale aunque cambien la duración o el número de preguntas, y en el examen solo se aplica.
\item \emph{La regla}, con $T$ minutos y $n$ preguntas:
  \begin{enumerate}
  \item \emph{Lectura y plan}: \tiempo{10}-\tiempo{12} al principio, sea cual sea la duración (en torno al 10~\% del tiempo). Se escribe solo en el cuadernillo (apartado~\ref{sec:leer}).
  \item \emph{Repaso final}: \tiempo{4}-\tiempo{8} (en torno al 5~\%), para releer, completar la frase de diagnóstico y cerrar lo que haya quedado abierto.
  \item \emph{Preguntas}: el resto, $W=T-\text{lectura}-\text{repaso}$.
    \begin{itemize}
    \item Si las preguntas no tienen puntuación, cada una recibe $t=W/(n+0{,}25)$ y el perfil cíclico, que es la más larga, $1{,}25\,t$.
    \item Si el tribunal da la puntuación (como en 2018, 2019, 2021 y 2022), el tiempo se reparte en proporción a los puntos: $W$ entre los puntos totales da los minutos por punto.
    \end{itemize}
  \end{enumerate}
\item \emph{Con el formato vigente} (1 hora y 45 minutos, 5 o 6 preguntas sin puntuación), la regla da \tiempo{12} de lectura y, con 5 preguntas, \tiempo{21} para el perfil cíclico y \tiempo{17} para cada una de las demás; con 6, \tiempo{18} y \tiempo{14} (tabla~\ref{tab:reparto}).
\end{itemize}

\begin{table}[H]
\centering\small
\caption{Reparto del tiempo según la regla, por convocatoria}\label{tab:reparto}
\begin{tabular}{>{\raggedright\arraybackslash}p{2.3cm}>{\centering\arraybackslash}p{3.3cm}>{\centering\arraybackslash}p{1.4cm}>{\centering\arraybackslash}p{2.3cm}>{\centering\arraybackslash}p{2.4cm}>{\centering\arraybackslash}p{1.4cm}}
\toprule
\rowcolor{tceeazulpalido}
\textbf{Examen} & \textbf{Formato} & \textbf{Lectura} & \textbf{Perfil cíclico} & \textbf{Cada una de las demás} & \textbf{Repaso} \\
\midrule
\rowcolor{tceesalvia}
2025, 2026 & 105 min, 5 preg. & \tiempo{12} & \tiempo{21} & \tiempo{17} & \tiempo{4} \\
\rowcolor{tceesalvia}
2023, 2024 & 105 min, 6 preg. & \tiempo{12} & \tiempo{18} & \tiempo{14} & \tiempo{5} \\
2022 & 135 min, 7 preg., 10 puntos & \tiempo{13} & \tiempo{17,5} (1,5 puntos) & \tiempo{11,5} por punto & \tiempo{7} \\
2021 & 135 min, 6 preg., 10 puntos & \tiempo{13} & \tiempo{17,5} (1,5 puntos) & \tiempo{11,5} por punto & \tiempo{7} \\
2020 & 135 min, 7 preg. & \tiempo{13} & \tiempo{19} & \tiempo{16} & \tiempo{7} \\
2019 & 135 min, 10 preg., 1 punto cada una & \tiempo{13} & \tiempo{11,5} & \tiempo{11,5} & \tiempo{7} \\
2018 & 120 min, 7 preg., 10 puntos & \tiempo{12} & \tiempo{20} (2 puntos) & \tiempo{10} por punto & \tiempo{8} \\
2016 (enero y noviembre) & 135 min, 8 preg. & \tiempo{13} & \tiempo{17} & \tiempo{14} & \tiempo{7} \\
\bottomrule
\end{tabular}
\fuente{Elaboración propia con la regla del texto y los formatos de la tabla~\ref{tab:formato}.}
\end{table}

\begin{itemize}
\item \emph{Ninguna pregunta en blanco}: si se acaba el tiempo de una, se cierra con una frase de conclusión y se pasa a la siguiente; al final, si sobra, se vuelve. Un esquema en viñetas es mejor que nada.
\item \emph{El perfil cíclico es el que más se alarga}: cayó en 13 de los 14 dictámenes de 2011 a 2026 (no en 2025) y casi siempre es la primera pregunta \parencite{ATCEE2026Examenes}. Ordena todo lo demás, pero lo que se diga de la inflación, el paro o el sector exterior va en su pregunta.
\item \emph{Antes de cada pregunta}, si el tiempo lo permite, 1 o 2 minutos para pensar la frase que anticipa el diagnóstico de ese bloque.
\end{itemize}

\subsection{Cómo leer el cuadernillo}\label{sec:leer}

\begin{itemize}
\item La idea: en los 10-12 primeros minutos no se escribe nada en las hojas de respuesta. Se trabaja sobre el cuadernillo, en el que sí se puede anotar, y se sale de ahí con el dictamen pensado.
\item Los pasos, por este orden:
  \begin{enumerate}
  \item \emph{Leer las instrucciones y todas las preguntas enteras}, subrayando los periodos («entre 2019 y 2022») y lo que pide cada subpregunta. Si las preguntas o las tablas vienen grapadas, se separan para tenerlas a la vista.
  \item \emph{Mirar las unidades} de cada fila: tasa de variación, nivel, porcentaje del PIB, puntos porcentuales, moneda nacional por dólar o al revés. En 2026 el tribunal tuvo que rectificar las unidades del tipo de cambio y de otros cuadros del enunciado \parencite{ATCEE2026Examenes}.
  \item \emph{Dividir el ciclo}: con el crecimiento del PIB real y, si la hay, la brecha de producción, trazar líneas verticales en el cuadro entre las fases (apartado~\ref{sec:etapas}) y ponerles nombre.
  \item \emph{Marcar en verde las fortalezas y en rojo las debilidades}: datos anómalos, máximos y mínimos, cambios de tendencia. Un indicador poco habitual suele estar ahí para que se use.
  \item \emph{Emparejar cada gráfico, tabla y fila con su pregunta}: escribir al lado «P2», «P4». Ningún gráfico está de adorno (días de huelga, uso de la capacidad productiva, precios de los activos…).
  \item \emph{Hacer al principio los cálculos que se pidan} (y los que se vayan a usar), por 2 motivos: se está más fresco y se evita contradecirse en la primera respuesta.\footnote{En las preguntas de cálculo hay que escribir el procedimiento, no solo el resultado. Desde 2022 no se han pedido cálculos, pero el análisis los usa (contribuciones, productividad, CLU, brecha); las fórmulas y ejemplos están en el anexo~\ref{anexo:calculos}.}
  \item \emph{Anotar las 2 o 3 ideas del temario} que encajan (una regla fiscal, el traspaso del tipo de cambio, la curva de \autor{Beveridge}) y empezar la lista de fortalezas y debilidades (apartado~\ref{sec:matriz}).
  \end{enumerate}
\end{itemize}

\begin{table}[H]
\centering\footnotesize
\caption{Un cuadro macroeconómico leído: el dictamen de 2026}\label{tab:cuadernillo}
\setlength{\tabcolsep}{3pt}
\begin{tabular}{>{\raggedright\arraybackslash}p{3.5cm}*{11}{>{\centering\arraybackslash}p{0.81cm}}}
\toprule
\rowcolor{tceeazulpalido}
 & \multicolumn{5}{c}{\textbf{Crecimiento débil o moderado}} & \multicolumn{2}{c}{\textbf{Recesión}} & \textbf{Recup.} & \multicolumn{3}{c}{\textbf{Desaceleración}} \\
\rowcolor{tceeazulpalido}
\textbf{Indicador} & \textbf{2014} & \textbf{2015} & \textbf{2016} & \textbf{2017} & \textbf{2018} & \textbf{2019} & \textbf{2020} & \textbf{2021} & \textbf{2022} & \textbf{2023} & \textbf{2024} \\
\midrule
PIB real (\%) & 0,9 & 1,8 & 0,7 & 1,6 & 0,8 & $-$0,3 & \cellcolor{tceesalmon}$-$4,3 & 3,6 & 1,3 & 0,7 & \cellcolor{tceesalmon}$-$0,2 \\
Contrib. demanda interna (pp) & 0,9 & 1,3 & 0,2 & 1,1 & 0,8 & 0,1 & $-$3,5 & 2,5 & 1,8 & 0 & $-$0,2 \\
Contrib. demanda externa (pp) & 0 & 0,5 & 0,5 & 0,6 & 0 & $-$0,5 & $-$0,8 & 1,1 & $-$0,5 & 0,7 & 0 \\
RNB real (\%) & 0,9 & \cellcolor{tceesalvia}3,3 & 1,1 & 1,2 & 0,2 & $-$0,1 & $-$3,7 & 3,3 & 0,7 & \cellcolor{tceesalvia}1,5 & 0,6 \\
Tasa de paro (\%) & 3,6 & 3,4 & 3,1 & 2,8 & \cellcolor{tceesalvia}2,4 & 2,4 & 2,8 & 2,8 & 2,6 & 2,6 & \cellcolor{tceesalvia}2,5 \\
Tasa de dependencia de los mayores (\%) & 42,3 & 43,8 & 45,1 & 46,1 & 46,9 & 47,5 & 48,0 & 48,6 & 48,8 & 49,0 & \cellcolor{tceesalmon}49,2 \\
IPC (\%) & 2,7 & 0,8 & $-$0,1 & 0,5 & 1,0 & 0,5 & 0,0 & $-$0,2 & 2,5 & \cellcolor{tceesalmon}3,2 & 2,7 \\
Tipo de cambio efectivo real (2020 $=$ 100) & 96,0 & 91,6 & 102,8 & 97,0 & 95,9 & 99,1 & 100,0 & 91,3 & 78,1 & 71,6 & \cellcolor{tceesalvia}66,9 \\
Saldo por cuenta corriente (\% PIB) & 0,7 & 3,0 & 3,8 & 4,0 & 3,4 & 3,4 & 2,9 & 3,7 & 2,0 & 3,6 & \cellcolor{tceesalvia}4,5 \\
Saldo presupuestario (\% PIB) & \cellcolor{tceesalmon}$-$5,6 & $-$3,6 & $-$3,6 & $-$3,0 & $-$2,4 & $-$3,0 & \cellcolor{tceesalmon}$-$9,0 & $-$6,3 & $-$4,2 & $-$2,3 & $-$1,7 \\
\bottomrule
\end{tabular}
\fuente{Selección de filas del cuadernillo del dictamen de 2026 (versión rectificada) \parencite{ATCEE2026Examenes}; las fases y los colores son de la lectura: en salvia, fortalezas; en salmón, debilidades.}
\end{table}

\begin{itemize}
\item \emph{Qué hay que ver en la tabla~\ref{tab:cuadernillo}}: el cuadro no da la brecha de producción, así que las fases salen solo del crecimiento del PIB.
  \begin{itemize}
  \item El crecimiento es bajo en toda la serie (nunca supera el 1,8~\%) con un paro muy bajo (del 3,6 al 2,5~\%): el potencial parece bajo, y la población que cae y la tasa de dependencia al alza apuntan a la demografía (pregunta 3).
  \item La RNB crece más que el PIB en varios años y el superávit corriente lo explican las rentas primarias: un acreedor neto frente al exterior (preguntas 1 y 2).
  \item El tipo de cambio real se deprecia en torno a un tercio desde 2020 mientras el IPC sube solo al 2,5-3,2~\%: indicio de un traspaso bajo (pregunta 2).
  \item El déficit público es persistente, con un salto en 2020, aunque se reduce desde entonces (pregunta 4).
  \end{itemize}
\item Ni hace falta ni conviene adivinar el país: basta con lo que dice el cuadernillo.
\end{itemize}

\subsection{Cómo redactar}\label{sec:redactar}

\begin{itemize}
\item \emph{Frase de diagnóstico}: cada respuesta empieza con 1 o 2 frases que dan la conclusión, como un resumen ejecutivo, y el resto la justifica con datos. Así abren sus informes el FMI en las consultas del artículo~IV y el Banco de España en su informe anual \parencite{FMI2026ArticuloIV,BdE2026InformeAnual}. La frase de diagnóstico general de la economía se escribe al final, en un hueco que se deja al principio de la primera respuesta (apartado~\ref{sec:etapas}).
\item \emph{Analizar, no describir}:
  \begin{itemize}
  \item Descriptivo (a evitar): «la contribución de la demanda interna es positiva en 2014, negativa en 2015 y vuelve a ser positiva en 2016».
  \item Analítico: «la demanda interna contribuye con fuerza al crecimiento en las fases de bonanza y se desploma en las recesiones, cuando la demanda externa amortigua en parte la caída: la interna es procíclica y la externa, contracíclica».
  \item La descripción de los datos es tarea de la estadística; al economista le toca explicar lo que hay detrás.
  \end{itemize}
\item \emph{Estructura visible}: rótulos («a) Perfil cíclico», «b) Determinantes de la inversión», «Pasando a la segunda cuestión…»), párrafos cortos, una idea por párrafo, interlineado amplio y márgenes grandes. Cada pregunta empieza en hoja nueva (si una se alarga, su continuación va en una hoja «1 bis», no en la de la siguiente).
\item \emph{Cifras justas}: 1 o 2 datos por idea, con su año y su unidad («la contribución de la demanda externa pasa de $-1{,}2$~pp a $+0{,}8$~pp»). No se copian filas enteras.
\item \emph{Condicionales y matices}: con pocos datos y poco tiempo solo se ven correlaciones, no causalidad (\vertema{5.B.4}). «Podría deberse a…», «parece reflejar…», «sería coherente con…». «\emph{An economist is an expert who will know tomorrow why the things he predicted yesterday didn't happen today}», reza una frase atribuida a Laurence J. \autor{Peter} \parencite{Peter1977}. Tampoco hay que ir al extremo contrario: cada respuesta tiene que llegar a una conclusión.
\item \emph{Sin solapes}: cada cosa en su pregunta, tratadas como estancas. Si algo ya se ha explicado, basta una remisión («como se ha visto al analizar la demanda…»).
\item \emph{Conocimiento a cuento}: una línea de teoría que explique un dato vale más que un párrafo de teoría sin datos. Se evitan las frases hechas sin contenido («la brecha de producción es el mejor indicador de la posición cíclica»).
\item \emph{Recomendaciones concretas}: qué medida, para qué desequilibrio y con qué riesgo («una consolidación gradual, centrada en el gasto corriente, para no frenar la recuperación»).
\end{itemize}

\subsection{Errores típicos y la lectura}\label{sec:errores}

\begin{itemize}
\item \emph{Errores que más se repiten}:
  \begin{itemize}
  \item describir sin analizar («en 2015 sube, en 2016 baja, en 2017 sube»);
  \item confundir porcentajes con puntos porcentuales, niveles con variaciones o tasas nominales con reales;
  \item hablar de «recuperación» cuando el PIB cae menos que el año anterior (pasar de $-7$~\% a $-1$~\% sigue siendo contracción);
  \item leer al revés el tipo de cambio (directo o indirecto) o la posición de inversión internacional;
  \item juzgar el tono de la política monetaria solo por si el tipo nominal sube o baja, o el de la fiscal por el saldo total y no por el estructural primario;
  \item dedicar demasiado tiempo al perfil cíclico y dejar a medias las últimas preguntas;
  \item no contestar a una subpregunta, o contestar a lo que no se pregunta.
  \end{itemize}
\item \emph{No adivinar el país}: aunque el país suele ser real, no es lo que se pide y, si se falla, lleva a errores graves (atribuirle un régimen cambiario, unas instituciones o una historia que no tiene). Se analiza en abstracto («esta economía») y solo con la información del cuadernillo, salvo que el enunciado diga otra cosa.
\item \emph{La lectura}: lo escrito no se puede cambiar. Se lee tal cual, con buena dicción, a un ritmo pausado y marcando los rótulos para que el tribunal sepa en qué pregunta está; por eso la redacción tiene que poder seguirse de oído. Si el tribunal pregunta (hasta 15 minutos), suele pedir que se aclare o justifique algo del texto o que se complete una respuesta corta: se contesta con calma, con la misma prudencia, y si se detecta un error propio se reconoce.
\end{itemize}

\begin{notaopositor}
Formato práctico: no usar corrector (lo que esté mal se tacha con una raya); cuidado con el papel autocopiativo, que calca lo que se apoye encima (separar las hojas ya escritas); numerar las hojas.
\end{notaopositor}

\subsection{La matriz de fortalezas y debilidades}\label{sec:matriz}

\begin{itemize}
\item La idea: durante todo el examen se va rellenando, en una hoja aparte o en el margen del cuadernillo, una lista de fortalezas (en verde) y debilidades (en rojo) por bloque, con el indicador que las muestra y su importancia.
\item Sirve para 3 cosas:
  \begin{enumerate}
  \item anticipar el diagnóstico en la primera pregunta (apartado~\ref{sec:anticipar});
  \item no contradecirse entre preguntas;
  \item contestar la pregunta final de desequilibrios y recomendaciones, que cayó en 12 de los 14 dictámenes, casi siempre en último lugar \parencite{ATCEE2026Examenes}.
  \end{enumerate}
\end{itemize}

\begin{table}[H]
\centering\small
\caption{Matriz de fortalezas y debilidades (plantilla con ejemplos)}\label{tab:matriz}
\begin{tabular}{>{\raggedright\arraybackslash}p{2.4cm}>{\raggedright\arraybackslash}p{4.6cm}>{\raggedright\arraybackslash}p{4.6cm}>{\centering\arraybackslash}p{1.8cm}}
\toprule
\rowcolor{tceeazulpalido}
\textbf{Bloque} & \textbf{\textcolor{tceeverde}{Fortalezas} (indicador)} & \textbf{\textcolor{tceerojo}{Debilidades} (indicador)} & \textbf{Peso} \\
\midrule
Perfil cíclico & crecimiento sostenido y equilibrado (contribuciones de demanda interna y externa positivas) & débil crecimiento o recesión; consumo anémico; dependencia de un sector & *** \\
Mercado de trabajo & creación de empleo; paro a la baja & paro alto y persistente; CLU que crecen más que en los socios & *** \\
Sistema financiero & solvencia y liquidez holgadas & endeudamiento privado alto; morosidad; racionamiento del crédito & ** \\
Precios & inflación cerca del objetivo & inflación alta y persistente; sobrecalentamiento & ** \\
Sector exterior & superávit corriente; PIIN mejorando & déficit corriente persistente; deuda externa alta; pérdida de competitividad & *** \\
Cuentas públicas & saldo primario estructural positivo; deuda baja & déficit y deuda altos; trayectoria insostenible & *** \\
\bottomrule
\end{tabular}
\fuente{Elaboración propia. La recomendación tipo para cada debilidad está en el apartado~\ref{sec:por-desequilibrio}.}
\end{table}

\begin{itemize}
\item \emph{Dónde está la teoría}: cada indicador remite al tema donde se estudia a fondo, primero el del tercer ejercicio y después el del cuarto; la tabla~\ref{tab:donde-se-estudia}, en los anexos, los reúne por bloque.
\item \emph{Recapitulación}: el método se resume en 4 verbos: \emph{leer} (10-12 minutos con el cuadernillo), \emph{repartir} (la regla del tiempo), \emph{diagnosticar} (una frase al empezar cada respuesta) y \emph{anotar} (la matriz de fortalezas y debilidades). Con él, pasamos a la primera pregunta de casi todos los dictámenes: el perfil cíclico.
\end{itemize}

% ==========================================================================
\section{Perfil cíclico de la economía}\label{sec:perfil}

\begin{itemize}
\item La conclusión por adelantado: la pregunta del perfil cíclico se contesta en 3 pasos, siempre los mismos:
  \begin{enumerate}
  \item dividir la serie en fases con el crecimiento del PIB real y la brecha de producción;
  \item explicar qué tira del crecimiento, por el lado de la demanda (interna o externa) y, si hay datos, por el de la oferta (empleo o productividad);
  \item anticipar el diagnóstico, con los desequilibrios y las fortalezas.
  \end{enumerate}
\end{itemize}

\begin{ideaclave}
\textbf{Qué hay que contestar}: ¿en qué fase está la economía y desde cuándo? ¿Qué tira del crecimiento: la demanda interna o la externa; el empleo o la productividad? ¿Crece por encima o por debajo de su potencial?
\end{ideaclave}

\begin{notaopositor}
Conviene hacer esta pregunta en primer lugar, porque ordena todo lo demás, pero sin pasarse del tiempo previsto (apartado~\ref{sec:tiempo}): en la coyuntura nunca sobra tiempo, y hay cosas mejores que hacer que analizar el último detalle de esta guía.
\end{notaopositor}

\subsection{Diagnóstico y etapas del ciclo}\label{sec:etapas}

\paragraph{La frase ejecutiva}
\begin{itemize}
\item El FMI, en los informes de las consultas del artículo~IV, y el Banco de España, en su informe anual, empiezan con una frase que resume el estado de la economía \parencite{FMI2026ArticuloIV,BdE2026InformeAnual}. Es una buena manera de empezar el dictamen.
\item Se limita a la situación actual (si hay previsiones, se mencionan sin más) y se escribe al final del examen, en un hueco que se deja antes de la primera respuesta, cuando ya se conoce bien la economía.
\item Ejemplo: el FMI abría así su evaluación de España en 2026: la economía española «siguió creciendo más que la zona del euro pese al impacto adverso de la guerra en Oriente Medio» \parencite{FMI2026ArticuloIV}.\footnote{Traducción propia de la nota de prensa del FMI. Una frase de diagnóstico completa, en el estilo de la del FMI de 2023 sobre España, encadena 3 ideas: dónde está la economía y qué la ha sostenido, qué la está frenando y qué riesgos tiene \parencite{FMI2023ArticuloIV2022}.}
\end{itemize}

\paragraph{Las fases del ciclo}
\begin{itemize}
\item El \concepto{ciclo económico} es la oscilación de la actividad alrededor de su tendencia: de un pico a un valle, la contracción; de un valle a un pico, la expansión \vertema{3.A.42}.
\item En el dictamen se usan 4 fases, que se reconocen por el crecimiento del PIB real ($\gamma_Y$) frente al potencial ($\gamma_{Y^p}$) y por el signo de la brecha de producción ($\mathrm{OG}$, apartado~\ref{sec:medidores}):
  \begin{itemize}
  \item \concepto{recuperación}: el PIB crece y más que el potencial, pero la brecha aún es negativa (se está cerrando);
  \item \concepto{expansión}: la brecha es positiva y se sigue abriendo;
  \item \concepto{desaceleración}: el PIB aún crece, pero menos que el potencial, y la brecha se reduce;
  \item \concepto{recesión}: el PIB cae y la brecha se vuelve negativa (gráfico~\ref{fig:fases-ciclo}).
  \end{itemize}
\item La forma de la recesión se describe con una letra: en \emph{V} (caída y recuperación rápidas, como la de 2020), en \emph{U} (estancamiento en el fondo), en \emph{W} (recaída tras una primera recuperación) o en \emph{L} (sin recuperación del nivel anterior).
\item Recesión es, como regla práctica, una caída del PIB real durante 2 trimestres seguidos; depresión, una caída de más del 10~\% o una recesión de más de 3 años.\footnote{La regla de los 2 trimestres consecutivos es una convención periodística popularizada por \textcite{Shiskin1974}. Los comités que fechan los ciclos (el NBER en EE.~UU. y el del CEPR en la zona del euro) usan un criterio más amplio: un descenso significativo y generalizado de la actividad durante más de unos meses, mirando también el empleo y la producción industrial \parencite{NBER2023Fechado}. Para la depresión no hay definición oficial: la regla del texto es la que recoge \textcite{Economist2008Depression}.}
\end{itemize}

\begin{figure}[H]
\centering
\caption{Las 4 fases del ciclo según la brecha de producción}\label{fig:fases-ciclo}
\grafico[0.75\textwidth]{fases-ciclo-brecha}
\fuente{Elaboración propia. Qué hay que ver: arriba, el PIB crece más que su potencial y la brecha avanza hacia la derecha (recuperación y expansión); abajo, crece menos o cae y la brecha retrocede (desaceleración y recesión). El paso por cero separa cada pareja de fases.}
\end{figure}

\begin{itemize}
\item \emph{Cómo se escribe}: «el cuadro recoge 2 etapas: una expansión entre 2014 y 2017, en la que la economía crece por encima de su potencial y la brecha se hace positiva, y una desaceleración en 2018-2019, en la que sigue creciendo, pero menos, y la brecha se cierra; en 2020 la recesión rompe el ciclo». Si el cuadro recoge más de un ciclo (por ejemplo, la crisis financiera y la de la COVID-19), se dice al principio.
\item \emph{Visión de conjunto}: lo atípico se señala como excepción («crecimiento sostenido entre 2014 y 2018, salvo una ligera desaceleración en 2016»), sin contar la serie año a año.
\end{itemize}

\subsection{Medidores}\label{sec:medidores}

\begin{itemize}
\item La idea: para dividir la serie se usa el crecimiento del PIB real, que el cuadro da siempre, y se completa con la brecha de producción, que no siempre está. Se elige el que dé fases más claras.
\item No siempre coinciden: un país puede crecer al 2~\% con la brecha abriéndose hacia abajo si su potencial crece al 3~\%. Si discrepan, se dice y se explica.
\end{itemize}

\paragraph{Crecimiento del PIB real}
\begin{itemize}
\item \emph{Qué mide}: el ritmo de avance de la producción, sin el efecto de los precios.
\item \emph{Cómo se calcula}: $\gamma_{Y,t}=\frac{Y_t-Y_{t-1}}{Y_{t-1}}\times100$, con $Y$ el PIB en volumen.
\item \emph{Umbral o referencia} (orientativo, sin norma detrás): en una economía avanzada, negativo es contracción; entre el 0 y el 1~\%, estancamiento; entre el 1 y el 2~\%, crecimiento moderado, cerca del potencial de la mayoría de las avanzadas; entre el 2 y el 3~\%, crecimiento por encima del potencial; desde el 3~\%, crecimiento sólido. Las economías avanzadas crecieron un 2,0~\% en 2025 y las emergentes y en desarrollo, un 4,5~\% \parencite{FMI2026WEO}: en un emergente, un 3~\% puede ser débil.
\item \emph{Cómo se lee}:
  \begin{itemize}
  \item el \emph{signo} y el \emph{nivel} dicen si crece y cuánto;
  \item la \emph{variación} (la «segunda derivada») dice si se acelera o se desacelera;
  \item en los países en desarrollo las tasas son más altas y más volátiles, porque tienen menos estabilizadores automáticos.
  \end{itemize}
\item \emph{Con qué se cruza}: con la brecha, el empleo y la inflación.
\item \emph{Tema}: \vertema{3.A.30} \vertema{4.A.1}.
\end{itemize}

\paragraph{Brecha de producción y PIB potencial}
\begin{itemize}
\item \emph{Qué mide}: la \concepto{brecha de producción} (\emph{output gap}) compara el PIB observado con el \concepto{PIB potencial}, el que se puede producir con los factores empleados con su intensidad normal y sin tensiones inflacionistas.\footnote{En una economía de mercado siempre hay recursos ociosos, porque cambian continuamente de actividad: hay un nivel normal de inactividad, y el producto que le corresponde (con una inflación estable) es el potencial. Por eso se le llama también producto no aceleracionista de la inflación, asociado a la NAIRU \vertema{3.A.28}. La producción puede estar por debajo del potencial cuando hay mucho paro de trabajo y de capital, y por encima durante un tiempo limitado cuando la demanda es muy fuerte y los factores se usan más de lo normal.}
\item \emph{Cómo se calcula}: en porcentaje del potencial,
\[
\mathrm{OG}_t\equiv\frac{Y_t-Y^p_t}{Y^p_t}\times100 .
\]
El potencial no se observa: se estima con una \emph{función de producción} (potencial del trabajo con la NAWRU, del capital y de la productividad total de los factores), como hacen la Comisión Europea, la OCDE y el FMI \parencite{Havik2014Potencial} \vertema{4.B.4}, o con un \emph{filtro estadístico}, como el de \autor{Hodrick} y \autor{Prescott}, que extrae la tendencia de la serie \parencite{HodrickPrescott1997} \vertema{3.A.42}. El filtro pide menos datos, pero su último tramo se revisa mucho.
\item \emph{Umbral o referencia}: 0, la producción en su potencial.
\item \emph{Cómo se lee}:
  \begin{itemize}
  \item \emph{Negativa}: producción por debajo del potencial, recursos infrautilizados (paro, capacidad ociosa) y presiones a la baja sobre la inflación.
  \item \emph{Positiva}: la demanda empuja por encima del potencial; aparecen cuellos de botella y riesgo de sobrecalentamiento e inflación.
  \item \emph{Cercana a 0}: paro y uso del capital en sus niveles normales e inflación estable. Es lo deseable.
  \item Su \emph{variación}, medida en puntos, dice si el PIB crece más o menos que el potencial: $\Delta\mathrm{OG}_t\approx\gamma_{Y,t}-\gamma_{Y^p,t}$. «La brecha se cierra» o «se abre» (anexo~\ref{anexo:expresiones}).
  \end{itemize}
\item \emph{Con qué se cruza}: si el cuadro da la brecha y el crecimiento, se deduce el crecimiento del potencial,\footnote{Los choques de oferta desplazan el PIB potencial y, con él, la brecha: el potencial crece con más capital (inversión en bienes de equipo), más trabajo (inmigración) o más productividad.}
\[
\gamma_{Y^p,t}\approx\gamma_{Y,t}-(\mathrm{OG}_t-\mathrm{OG}_{t-1}),
\]
que se pregunta a veces (2018 y 2021). Si la economía crece con fuerza y la brecha no se cierra, el potencial crece mucho: no es un cuadro «débil». Al revés, cerrar la brecha creciendo poco no es buena señal: el potencial crece poco.
\item \emph{Cuidado}: la brecha y el potencial se estiman con el crecimiento pasado y se revisan mucho; sirven poco en economías con cambios estructurales rápidos, y no conviene dividir el ciclo solo con ellos.
\item \emph{Tema}: \vertema{3.A.42} \vertema{3.A.45} \vertema{3.A.28} \vertema{4.B.4}.
\end{itemize}

\paragraph{PIB frente a renta nacional bruta}
\begin{itemize}
\item \emph{Qué mide}: el PIB es lo que se produce \emph{dentro} del país; la \concepto{renta nacional bruta} (RNB), lo que ganan sus \emph{residentes}, dentro o fuera.
\item \emph{Cómo se calcula}: la RNB es el PIB más las rentas primarias recibidas del exterior menos las pagadas al exterior. Las rentas primarias son las del trabajo (trabajadores transfronterizos), las de la inversión (beneficios, dividendos e intereses de la inversión directa, la de cartera y los préstamos) y algunos impuestos y subvenciones con la UE. Si se suman las transferencias corrientes netas (remesas, cooperación), se obtiene la renta nacional disponible \parencite{SEC2010}.
\item \emph{Umbral o referencia}: no hay; importa el signo y si la diferencia es grande y persistente.
\item \emph{Cómo se lee}:
  \begin{itemize}
  \item $\mathrm{RNB}<\mathrm{PIB}$: el país paga al exterior más rentas de las que recibe. Es típico de un deudor neto (intereses de la deuda externa) o de un país con mucha inversión extranjera directa, cuyas multinacionales repatrían beneficios.
  \item $\mathrm{RNB}>\mathrm{PIB}$: acreedor neto (rentas de sus inversiones fuera) o país con muchos trabajadores en el extranjero.
  \item Un salto puntual suele deberse a multinacionales que trasladan activos (patentes) o beneficios: el PIB sube sin que suba la renta de los residentes, como en Irlanda en 2015.
  \end{itemize}
\item \emph{Con qué se cruza}: con la cuenta de rentas primarias de la balanza de pagos, la posición de inversión internacional neta y la inversión directa (capítulo~\ref{sec:exterior}). Se preguntó en 2023 y 2026.
\item \emph{Tema}: \vertema{3.A.30} \vertema{3.B.11} \vertema{4.A.21}.
\end{itemize}

\paragraph{Indicadores adelantados y de alta frecuencia}
\begin{itemize}
\item \emph{Qué miden}: el cuadro macroeconómico es anual o trimestral y llega con retraso; estos indicadores son mensuales (o diarios) y algunos se anticipan al PIB. Dicen cómo está la economía \emph{ahora} y hacia dónde va (2016, 2019 y 2024).
\item \emph{Los más habituales}:
  \begin{itemize}
  \item \concepto{Índice de gestores de compras} (PMI) de manufacturas y de servicios: un índice de difusión (cuántas empresas dicen que la producción, los pedidos o el empleo suben, bajan o siguen igual que el mes anterior). Por encima de 50, la actividad \emph{crece} respecto al mes anterior (no que esté alta); por debajo, se contrae \parencite{SPGlobalPMI}.
  \item \concepto{Indicador de sentimiento económico} de la Comisión Europea (ESI) y los de confianza de consumidores y empresarios: su media de largo plazo es 100 \parencite{CE2025Encuestas}.
  \item \emph{Pendiente de la curva de tipos} (tipo a 10 años menos tipo a corto): si se invierte, suele anticipar una recesión \parencite{EstrellaMishkin1998}.\footnote{La \emph{curva de tipos} ordena los rendimientos de la deuda pública por plazo. Si el cuadro da 2 tipos (corto y largo), la «curva» es su diferencia. \emph{Creciente}: lo normal, por la prima de liquidez; si se empina, los inversores esperan subidas de tipos, a menudo porque esperan crecimiento. \emph{Plana}: se esperan bajadas que compensan la prima. \emph{Invertida}: los inversores esperan que el banco central baje los tipos para combatir una recesión y compran a largo plazo, lo que baja su rentabilidad \vertema{3.B.24}.}
  \item Otros adelantados: bolsa, crédito nuevo, pedidos industriales, permisos de construcción.
  \item \emph{Coincidentes}: producción industrial, afiliación, ventas minoristas, consumo de electricidad.
  \item \emph{Retrasados}: tasa de paro, inflación subyacente, tipos de interés de los préstamos \parencite{ConferenceBoard2026LEI}.
  \end{itemize}
\item \emph{Umbral o referencia}: 50 en el PMI; 100 en el ESI; pendiente negativa en la curva de tipos.
\item \emph{Cómo se lee}: importa la dirección y que varios coincidan. Un PMI que baja de 55 a 51 sigue indicando crecimiento, pero más lento. Si los indicadores mensuales caen mientras el PIB anual aún crece, la economía se está desacelerando.
\item \emph{Con qué se cruza}: con el perfil cíclico («los indicadores de alta frecuencia apuntan a una desaceleración a finales de año») y con las previsiones del cuadro. La clasificación de las variables por su adelanto o retraso está en el anexo~\ref{anexo:series}.
\item \emph{Tema}: \vertema{3.A.42} \vertema{4.A.1}.
\end{itemize}

\subsection{Demanda}\label{sec:demanda}

\begin{itemize}
\item La idea: el lado de la demanda explica la evolución del PIB por el enfoque del gasto: quién compra lo que se produce. Es lo primero que se mira, porque el cuadro siempre trae sus componentes.
\item Con $C$ el consumo privado, $I$ la inversión (formación bruta de capital), $G$ el consumo de las administraciones públicas, $X$ las exportaciones y $M$ las importaciones \vertema{3.A.30},
\[
Y_t=\underbrace{C_t+I_t+G_t}_{\text{demanda interna}}+\underbrace{X_t-M_t}_{\text{demanda externa neta}} .
\]
En la contabilidad nacional, $G$ es solo el consumo público; la inversión pública está dentro de $I$.
\end{itemize}

\begin{ideaclave}
\textbf{Qué hay que contestar}: ¿cuál es el motor del crecimiento, la demanda interna o la externa? ¿Qué componentes explican cada fase y por qué?
\end{ideaclave}

\paragraph{Contribuciones al crecimiento: el motor}
\begin{itemize}
\item \emph{Qué mide}: cuántos puntos del crecimiento del PIB aporta cada componente.
\item \emph{Cómo se calcula}: restando a la identidad de la demanda la del año anterior y dividiendo por $Y_{t-1}$,
\[
\gamma_{Y,t}=\underbrace{\frac{C_{t-1}}{Y_{t-1}}\gamma_{C,t}+\frac{I_{t-1}}{Y_{t-1}}\gamma_{I,t}+\frac{G_{t-1}}{Y_{t-1}}\gamma_{G,t}}_{\text{contribución de la demanda interna}}+\underbrace{\frac{X_{t-1}}{Y_{t-1}}\gamma_{X,t}-\frac{M_{t-1}}{Y_{t-1}}\gamma_{M,t}}_{\text{contribución de la demanda externa}} .
\]
Cada contribución es el peso del componente en el año anterior por su crecimiento: depende de los 2. Las importaciones restan.
\item \emph{Unidades}: las contribuciones se miden en \emph{puntos porcentuales}. Si el PIB crece un 3~\% y la demanda interna contribuye con 2~pp, no se puede decir que la demanda interna haya crecido un 2~\%: de los 3 puntos de crecimiento, 2 se deben a la demanda interna y 1 a la externa.
\item \emph{Cuidado}: con volúmenes encadenados la suma no es exacta, y las existencias también contribuyen (si el cuadro no las da, se suponen nulas). Si el cuadernillo da las contribuciones, se usan esas \parencite{Eurostat2026nama10gdp}.
\item \emph{Cómo se lee}:
  \begin{itemize}
  \item contribuciones positivas y parecidas de las 2: \emph{crecimiento equilibrado};
  \item en las crisis de economías con fuerte demanda interna, la interna pasa a restar (es \emph{procíclica}: caen el consumo y la inversión) y la externa a sumar (tiene un componente \emph{contracíclico}), porque las importaciones caen con la renta, el resto del mundo puede haber sufrido menos y las empresas buscan fuera la demanda que les falta dentro;
  \item si el sector exterior «toma el relevo» como motor en una crisis, hay que ver si suben las exportaciones o solo caen las importaciones: lo primero es una fortaleza; lo segundo, el reflejo de la caída de la renta;
  \item ante un cambio de patrón, ¿es coyuntural (un auge del precio del petróleo en un exportador) o estructural?
  \end{itemize}
\item \emph{Tema}: \vertema{3.A.30} \vertema{4.A.29}.
\end{itemize}

\begin{figure}[H]
\centering
\caption{Contribuciones al crecimiento del PIB real de la zona del euro}\label{fig:contribuciones-pib-zona-euro}
\grafico[0.85\textwidth]{contribuciones-pib-zona-euro}
\fuente{\textcite{Eurostat2026nama10gdp}; zona del euro de 21 países, 2012-2025. Qué hay que ver: la demanda interna explica casi todo el crecimiento salvo en las crisis; en 2020 se hunde el consumo privado y en 2023 restan las existencias; la demanda externa neta suma en 2012-2013 y en 2021 y resta en 2025 ($-0{,}7$~pp).}
\end{figure}

\paragraph{Consumo privado}
\begin{itemize}
\item \emph{Qué incluye}: el gasto de los hogares (y de las instituciones sin fines de lucro a su servicio) en bienes y servicios: alimentación, ropa, alquiler, bienes duraderos. No incluye la compra de vivienda, que es inversión.
\item \emph{Hechos estilizados} (comprobar si se cumplen en el cuadro):
  \begin{itemize}
  \item \emph{Peso}: el componente más cuantioso, algo más de la mitad del PIB en la UE (53~\% en 2025) y en torno a 2 tercios en EE.~UU. \parencite{Eurostat2026nama10gdp,BM2026WDI}.
  \item \emph{Volatilidad}: el menos volátil de los componentes privados, por su peso y porque los hogares prefieren suavizar el consumo según su renta permanente \vertema{3.A.33}. Si en el cuadro es más volátil que el PIB, puede haber sobrecalentamiento e inflación de demanda.
  \item \emph{Comovimiento}: procíclico y contemporáneo con el PIB; los bienes duraderos, más volátiles \parencite{StockWatson1999Ciclo}.
  \end{itemize}
\item \emph{Determinantes}, con el signo de su relación con el consumo:
  \begin{itemize}
  \item \emph{Renta disponible real} ($+$), sin duda el principal. A su vez depende del empleo y del salario real (que cae con la inflación), de las rentas del capital, de los impuestos y de las transferencias.\footnote{A veces el cuadro da directamente la renta disponible. La propensión a consumir la renta que procede de más empleo es mayor que la que procede de salarios reales más altos, porque el empleo señala mejor un aumento de la renta permanente.}
  \item \emph{Riqueza} ($+$): precio de la vivienda en propiedad y de los activos financieros (efecto riqueza; discutible en los hogares de renta baja, que tienen pocos activos).
  \item \emph{Tipo de interés real a corto plazo} ($-$): es el coste de endeudarse y lo que rinde el ahorro (los nominales se corrigen con el IPC).
  \item \emph{Inflación} ($-$ la actual; ambiguo la esperada): la actual reduce la renta real; la esperada también, pero adelanta compras.\footnote{El precio del petróleo pesa mucho en la renta real de los hogares, porque la demanda de energía es poco elástica.}
  \item \emph{Crédito} ($+$) y \emph{endeudamiento de los hogares} ($-$).
  \item \emph{Incertidumbre y paro esperado} ($-$): ahorro precautorio \vertema{3.A.10}.
  \item \emph{Demografía}: crece con la población; con la estructura de edades el efecto es ambiguo. Según la hipótesis del ciclo vital, ahorran sobre todo las personas en edad de trabajar \parencite{ModiglianiBrumberg1954}: más población dependiente (joven o mayor) supone más consumo en proporción a la renta \parencite{Leff1969Dependencia}.
  \item \emph{Expectativas} ($+$): confianza de los consumidores, PMI, pendiente de la curva de tipos, crecimiento mundial. El índice de miseria (inflación más paro) resume 2 frenos.
  \item \emph{Políticas futuras}: si los hogares esperan más impuestos para pagar la deuda que se emite hoy, ahorran más y el estímulo fiscal pierde fuerza (\concepto{equivalencia ricardiana}, \textcite{Barro1974}, \vertema{3.A.39}); una bajada de tipos anunciada adelanta consumo.
  \end{itemize}
\item \emph{Otro enfoque}: si los hogares siguen la hipótesis de la renta permanente \parencite{Friedman1957} y tienen expectativas racionales, solo cambian su consumo cuando les llegan noticias que revisan esa renta. El consumo sigue entonces un \emph{paseo aleatorio}, y el mejor pronóstico del consumo futuro es el presente \parencite{Hall1978}.\footnote{Una condición para que se cumpla es que los hogares puedan prestar y pedir prestado, un supuesto razonable solo con mercados financieros desarrollados.}
\item \emph{Cómo se escribe}: se vinculan los determinantes con los años: «entre los determinantes habituales del consumo, la renta disponible y el empleo evolucionan de modo que podrían explicar su caída en 2020 y su recuperación posterior».
\end{itemize}

\begin{figure}[H]
\centering
\caption{Consumo, renta disponible y tasa de ahorro de los hogares de la zona del euro}\label{fig:consumo-renta-ahorro}
\grafico[0.85\textwidth]{consumo-renta-ahorro-zona-euro}
\fuente{\textcite{Eurostat2026nama10gdp}; elaboración propia de la renta disponible real. Qué hay que ver: en 2020 la renta disponible apenas cae ($-0{,}3$~\%) y el consumo se desploma ($-7{,}7$~\%): la tasa de ahorro sube al 19~\% (ahorro forzoso y por precaución) y se gasta en 2021-2022. Desde 2023 el ahorro sigue alto (en torno al 15~\%) y el consumo crece menos que la renta.}
\end{figure}

\paragraph{Inversión}
\begin{itemize}
\item \emph{Qué incluye}: 3 partidas \parencite{SEC2010}:
  \begin{itemize}
  \item la \concepto{formación bruta de capital fijo} (FBCF): adquisiciones menos cesiones de activos fijos, los que se usan para producir más de un año (vivienda, otras construcciones, maquinaria y equipo, armamento, recursos biológicos, I+D, programas informáticos y otros derechos de propiedad intelectual).\footnote{La inversión residencial es el único componente de la inversión que hacen los hogares, y en algunos países el más importante. Depende del crédito, los tipos, la regulación y los impuestos, la renta de los hogares y la demanda de vivienda frente al parque disponible \vertema{4.A.11}.} Incluye la reposición del capital que se desgasta (consumo de capital fijo);
  \item la \concepto{variación de existencias}: entradas menos salidas de existencias;
  \item las adquisiciones menos cesiones de \emph{objetos valiosos} (bienes que se mantienen como depósito de valor).
  \end{itemize}
\item \emph{Por qué importa}: es demanda a corto plazo y oferta a largo, porque acumula capital y eleva el potencial \vertema{3.A.45}.
\item \emph{Hechos estilizados} (comprobar si se cumplen):
  \begin{itemize}
  \item \emph{Peso}: en torno a una quinta parte del PIB; la FBCF fue el 21,5~\% del PIB en la UE en 2025 \parencite{Eurostat2026nama10gdp}.
  \item \emph{Volatilidad}: el componente más volátil, varias veces más que el PIB, por su dependencia de las expectativas (los \emph{animal spirits} de \textcite{Keynes1936}).
  \item \emph{Comovimiento}: procíclica; la inversión en equipo, contemporánea o algo retrasada, y la residencial, adelantada \parencite{StockWatson1999Ciclo}.
  \item \emph{Existencias}: pesan poco, pero son muy volátiles y \emph{procíclicas}: en las recesiones las empresas dejan de reponerlas y eso explica buena parte de la caída del PIB \parencite{BlinderMaccini1991Existencias}. A muy corto plazo amortiguan la demanda inesperada (si la demanda sube por sorpresa, primero caen las existencias), pero en el ciclo amplifican; una acumulación excesiva puede anunciar la fase bajista.
  \end{itemize}
\item \emph{Determinantes}:
  \begin{itemize}
  \item \emph{Positivos}: crecimiento esperado de la demanda (interna y externa); \emph{brecha de producción y grado de utilización de la capacidad}, porque con la brecha positiva las empresas necesitan ampliar capacidad (acelerador); rentabilidad empresarial (excedente bruto de explotación o, si falta, el margen, apartado~\ref{sec:trabajo}); incentivos fiscales; crédito; expectativas (confianza empresarial, bolsa, curva de tipos).
  \item \emph{Negativos}: tipo de interés real a largo plazo;\footnote{El tipo real a largo plazo pide expectativas de inflación a largo plazo, que el cuadro casi nunca da. Por orden de preferencia: deflactar el tipo nominal a largo con la inflación del periodo (mejor con el deflactor del PIB que con el IPC, porque es más relevante para las empresas); usar el tipo real a corto plazo; usar el nominal a largo, que ya incluye una prima por la inflación. En todos los casos hay que decir qué se ha hecho y sus limitaciones.} costes de construcción y precio real del equipo; incertidumbre económica y política.\footnote{En los países en desarrollo, la falta de mercados financieros, el papel del Estado en la formación de capital y las restricciones cambiarias debilitan los modelos estándar de inversión: la disponibilidad de financiación limita más que su coste, y el crédito bancario, la entrada de capital extranjero y los beneficios retenidos son los determinantes principales.}
  \item \emph{Expulsión} (\emph{crowding-out}): con déficit público alto y tipos al alza, el endeudamiento público puede detraer financiación a la inversión privada.
  \item Si el cuadro no da indicadores de expectativas, se pueden mencionar como posible factor, con condicionales («es posible que el crecimiento del resto del mundo esté generando unas expectativas que expliquen la inversión»).
  \end{itemize}
\item \emph{Cómo se escribe}: «la inversión es, como es habitual, el componente más volátil; sus caídas coinciden con las del crédito al sector privado y sus aumentos, con tipos de interés reales negativos». Se preguntaron sus determinantes en 2023 y 2026.
\item \emph{Tema}: \vertema{3.A.34} \vertema{4.A.29}.
\end{itemize}

\paragraph{Consumo e inversión públicos}
\begin{itemize}
\item \emph{Qué incluyen}: el consumo público son sobre todo los gastos de personal y las compras de bienes y servicios de las administraciones; la inversión pública, las infraestructuras y el equipamiento (dentro de la FBCF).
\item \emph{Cómo se leen}: responden sobre todo a decisiones políticas. En las economías avanzadas el consumo público es acíclico o algo contracíclico (crece en las recesiones, como en 2020); en muchos emergentes es procíclico, porque el Estado gasta cuando entra financiación y recorta cuando se corta \parencite{Kaminsky2004Rains}. Si sube con fuerza en una recesión, lo probable es una respuesta discrecional, que se analiza en la pregunta de política fiscal (apartado~\ref{sec:fiscal}).
\item La inversión pública puede arrastrar a la privada (infraestructuras que la hacen rentable) o expulsarla: por superposición (expulsión real) o porque el endeudamiento público detrae financiación (expulsión financiera).
\end{itemize}

\paragraph{Demanda externa}
\begin{itemize}
\item En el perfil cíclico se comenta en 1 o 2 frases: su contribución y si la explican las exportaciones o las importaciones. Si faltan los datos de bienes y servicios, se usa la balanza comercial, diciéndolo. El análisis de sus determinantes va en la pregunta del sector exterior (capítulo~\ref{sec:exterior}).
\item \emph{Exportaciones}: dependen positivamente de la competitividad-precio (el tipo de cambio real, $q=EP^*/P$, cuya subida es una depreciación real), de la competitividad no precio (calidad, marca, diferenciación) y de la renta del resto del mundo ($Y^*$), y negativamente de la demanda interna: cuando esta cae, las empresas buscan fuera los mercados que les faltan dentro, y cuando repunta, se repliegan al mercado nacional.\footnote{Una depreciación puede tener un efecto ambiguo sobre la inversión: la anima en las empresas exportadoras, pero encarece los bienes intermedios importados (sobre todo el petróleo) y reduce márgenes. Con la condición de \autor{Marshall}-\autor{Lerner}, una depreciación aumenta la contribución de la demanda externa al crecimiento (apartado~\ref{sec:tipo-cambio}).}
\item \emph{Importaciones}: dependen positivamente de la renta nacional (son procíclicas, por su alta elasticidad-renta) y negativamente de la competitividad-precio y no precio de la producción nacional; y de rasgos estructurales, como la dependencia energética o tecnológica.
\end{itemize}

\begin{figure}[H]
\centering
\caption{Crecimiento del PIB y del volumen de comercio mundiales}\label{fig:pib-comercio-mundial}
\grafico[0.85\textwidth]{pib-comercio-mundial}
\fuente{\textcite{FMI2026WEO}; 2026 y 2027, previsiones. Qué hay que ver: el comercio es más volátil que el PIB (cae un 10~\% en 2009 y un 8~\% en 2020); antes de 2012 crecía bastante más que el PIB y desde entonces lo hace a un ritmo parecido.}
\end{figure}

\begin{notaopositor}
En el perfil cíclico, la demanda externa se comenta en 1 o 2 frases; sus determinantes, la competitividad y la balanza de pagos van en la pregunta del sector exterior.
\end{notaopositor}

\subsection{Oferta}\label{sec:oferta}

\begin{itemize}
\item La idea: el lado de la oferta explica el PIB por el enfoque de la producción. A corto plazo la economía depende de la demanda; a largo plazo, de la oferta. Se pregunta menos (2013, 2014 y 2024), pero siempre da 2 ideas: si se crece con productividad o con más factores, y qué sectores tiran.
\end{itemize}

\paragraph{Motor de crecimiento: productividad o acumulación de factores}
\begin{itemize}
\item \emph{Qué mide}: si el crecimiento se basa en mejoras de la productividad o en más trabajo y capital.
\item \emph{Cómo se calcula}: con una función de producción $Y=F(A,K,L)$, la contabilidad del crecimiento descompone
\[
\gamma_Y=\gamma_A+\alpha\,\gamma_K+(1-\alpha)\,\gamma_L,
\]
con $\alpha$ la participación del capital en la renta y $\gamma_A$ el crecimiento de la productividad total de los factores \vertema{3.A.45}. Como el cuadro casi nunca da el capital ni la productividad total, se usa la \concepto{productividad aparente del trabajo} (PAT), $\mathrm{PAT}=Y/L$:
\[
\gamma_{\mathrm{PAT}}=\gamma_Y-\gamma_L\quad\Longleftrightarrow\quad\gamma_Y=\gamma_{\mathrm{PAT}}+\gamma_L .
\]
\item \emph{Cómo se lee}:
  \begin{itemize}
  \item crecimiento con empleo y PAT estancada: patrón extensivo (más factores);
  \item crecimiento con PAT al alza: patrón intensivo, más sostenible;
  \item cuidado: si el empleo cae más que el PIB, la PAT «sube» de forma espuria (se destruye el empleo menos productivo), y la causalidad entre productividad y ciclo es difícil de establecer: mejor no concluir que la expansión se debe a un choque de productividad.
  \end{itemize}
\item \emph{Con qué se cruza}: con el empleo, las horas y los CLU (apartado~\ref{sec:trabajo}) y con la brecha: si la economía crece y la brecha no se cierra, el potencial crece (apartado~\ref{sec:medidores}).
\item \emph{Tema}: \vertema{3.A.45} \vertema{3.A.43} \vertema{4.A.29}.
\end{itemize}

\paragraph{Sectores}
\begin{itemize}
\item \emph{Qué se mira}: la evolución del valor añadido bruto (VAB) de cada sector y su peso, con 4 preguntas: ¿cambia la estructura productiva (industrialización, terciarización)?, ¿qué factores usa con intensidad cada sector (trabajo, capital, I+D, energía)?, ¿tira del resto (eslabonamientos)?, ¿genera externalidades (como la industria aeronáutica)?
\item \emph{Rasgos de cada sector}:
  \begin{itemize}
  \item \emph{Primario}: en general acíclico; intensivo en trabajo; baja elasticidad-renta; protección y subvenciones; más peso cuanto menos desarrollado el país \vertema{4.A.4}.
  \item \emph{Energía y materias primas}: precios fijados en los mercados mundiales (inflación importada en un importador), subvenciones o impuestos, efecto sobre la balanza comercial, capacidad de refino (a menudo con inversión extranjera) \vertema{4.A.5}.
  \item \emph{Construcción}: grandes eslabonamientos hacia atrás, intensiva en trabajo poco cualificado, maduración larga, oferta fija a corto plazo, demanda condicionada por la demografía y por su carácter de inversión; riesgo de burbuja, fuerte efecto riqueza y simbiosis con el sector financiero (círculos virtuosos o viciosos con el crédito) \vertema{4.A.11}.
  \item \emph{Industria}: menos intensiva en trabajo; ver su relación con las exportaciones y, en un país en desarrollo, si hay una estrategia de industrialización orientada a la exportación \vertema{4.A.9}.
  \item \emph{Servicios}: muy heterogéneos, de alto y bajo valor añadido y a menudo no comercializables; el turismo depende de la movilidad y suele concentrar empleo joven, temporal y poco cualificado \vertema{4.A.13}.
  \end{itemize}
\item \emph{Tamaño de las empresas}: un tejido muy atomizado hace al país más vulnerable a las crisis y limita la innovación, la internacionalización y la financiación, y con ello la productividad.
\end{itemize}

\paragraph{Un sector dominante: petróleo y materias primas}
\begin{itemize}
\item \emph{Qué hay que mirar} (2016 y 2025): el peso del sector en el VAB, en las exportaciones y en los ingresos públicos, y el precio internacional, en dólares.
\item \emph{Canales de un cambio del precio del petróleo}:
  \begin{itemize}
  \item \emph{PIB y precios}: en un exportador, una subida eleva la renta y la inversión del sector y el deflactor del PIB más que el IPC (mejora la relación real de intercambio); en un importador, sube el IPC (inflación importada) y baja la renta disponible.
  \item \emph{Sector exterior}: la balanza de bienes y la cuenta corriente siguen al precio; con tipo flexible, la moneda del exportador se aprecia.
  \item \emph{Cuentas públicas}: regalías e impuestos sobre el sector (exportador) o subvenciones a los combustibles; los ingresos se vuelven volátiles y procíclicos.
  \end{itemize}
\item \emph{Riesgos}: la \concepto{enfermedad holandesa}, en la que el auge del sector aprecia el tipo de cambio real y expulsa a otros sectores exportadores \parencite{CordenNeary1982}, y una política fiscal procíclica, que gasta en el auge y recorta en la caída \parencite{Kaminsky2004Rains}.
\item \emph{Recomendaciones}: un fondo soberano o una regla fiscal estructural que ahorre los ingresos extraordinarios, diversificar la producción y las exportaciones, y retirar poco a poco las subvenciones a los combustibles con ayudas focalizadas.
\item \emph{Tema}: \vertema{3.B.28} \vertema{3.B.32} \vertema{4.A.5}.
\end{itemize}

\subsection{Anticipar el diagnóstico}\label{sec:anticipar}

\begin{itemize}
\item La idea: si está claro y da tiempo, la respuesta del perfil cíclico se cierra con un diagnóstico general que anuncia lo que se desarrollará después. Se deja hueco y se escribe al final.
\item \emph{Cómo se escribe}: «como se verá, parece tratarse de una economía avanzada cuyo crecimiento se ha apoyado en la demanda interna, pero que presenta desequilibrios como un déficit exterior persistente y una deuda pública elevada, de modo que su cuadro macroeconómico tiene una salud mejorable».
\end{itemize}

\begin{notaopositor}
Los problemas estructurales son más difíciles de ver en los datos que los coyunturales. Por eso se formulan con condicionales y apoyados en 2 o 3 datos (crecimiento potencial bajo, productividad estancada, paro persistente).
\end{notaopositor}

\paragraph{Economía avanzada o emergente: qué cambia en el análisis}
\begin{itemize}
\item Sin intentar adivinar el país, hay rasgos que cambian la lectura: crecimiento potencial y volatilidad, inflación y traspaso del tipo de cambio, peso de las materias primas, política fiscal procíclica, deuda en divisas, riesgo de frenazo súbito de capitales y sistema financiero poco profundo (tabla~\ref{tab:avanzada-emergente}, en los anexos). Como referencia, en 2025 las economías avanzadas crecieron un 2,0~\% con una inflación del 2,5~\% \parencite{FMI2026WEO}; la renta per cápita era de unos 90.000 dólares en EE.~UU., 61.000 en Alemania y 39.000 en España \parencite{FMI2026WEO}; y el sector primario pesaba un 1,8~\% del VAB de la UE, frente a más de una cuarta parte del PIB en los países de renta baja \parencite{Eurostat2026nama10a10,BM2026WDI}.
\item Ojo con los países exportadores de petróleo, con rentas per cápita altas y rasgos de emergente.
\end{itemize}

\paragraph{Desequilibrios y fortalezas típicos}
\begin{itemize}
\item Los desequilibrios que más se repiten en los dictámenes, con el indicador que los delata:
  \begin{itemize}
  \item \emph{Débil crecimiento o recesión}: PIB real que crece poco o cae; crecimiento potencial bajo.
  \item \emph{Demanda interna anémica}: consumo privado que apenas crece.
  \item \emph{Inflación y sobrecalentamiento}: brecha positiva persistente, inflación alta, consumo que crece mucho, caída brusca de existencias y salto de la penetración de importaciones.
  \item \emph{Mercado de trabajo deficiente}: paro alto y persistente, CLU crecientes.
  \item \emph{Déficit y deuda públicos}: altos, persistentes o con trayectoria insostenible.
  \item \emph{Déficit por cuenta corriente y deuda externa}: altos y persistentes, con pérdida de competitividad, exceso de absorción, ahorro escaso o déficit público.
  \item \emph{Riesgos en el sistema financiero}: prima de riesgo alta, endeudamiento privado y público alto, morosidad, racionamiento del crédito.
  \item \emph{Estructura sectorial desequilibrada}: dependencia de un sector (inmobiliario, energético, materias primas).
  \end{itemize}
\item Las fortalezas son su reverso (tabla~\ref{tab:matriz}). Cada uno se desarrolla en su bloque, y la recomendación tipo para cada desequilibrio está en el apartado~\ref{sec:por-desequilibrio}.
\end{itemize}

\begin{itemize}
\item \emph{Recapitulación}: el perfil cíclico se resuelve con 3 piezas que se repiten en todos los dictámenes: las fases (con el crecimiento y la brecha), el motor (contribuciones por el lado de la demanda y, si hay datos, empleo frente a productividad por el de la oferta) y un diagnóstico anticipado. Una vez situada la economía en el ciclo, pasamos a los mercados de factores, donde el crecimiento se traduce en empleo, salarios y crédito.
\end{itemize}

% ==========================================================================
\section{Mercados de factores}\label{sec:factores}

\begin{itemize}
\item La conclusión por adelantado: en los mercados de factores se ve si el crecimiento es sostenible. En el de trabajo, si crea empleo, qué tipo de paro hay y si los salarios siguen a la productividad; en el de capitales, si el crédito, la deuda, la banca y los precios de los activos acumulan riesgos.
\item Desde el punto de vista macroeconómico, el trabajo es, junto con el capital, uno de los 2 factores de producción (la tierra se suele incluir en el capital).
\end{itemize}

\subsection{Mercado de trabajo}\label{sec:trabajo}

\begin{ideaclave}
\textbf{Qué hay que contestar}: ¿cuánto empleo crea el crecimiento? ¿Qué tipo de paro hay? ¿Los salarios siguen a la productividad? ¿Es flexible y competitivo el mercado de trabajo?
\end{ideaclave}

\begin{notaopositor}
En el mercado de trabajo suelen pedirse cosas muy concretas (definir la tasa de actividad, la curva de \autor{Beveridge}, un cálculo de CLU). Conviene contestarlas al principio de la respuesta, con un rótulo, para que el tribunal las encuentre al oírlas en la lectura.
\end{notaopositor}

\paragraph{Indicadores de cantidad}
\begin{itemize}
\item \emph{Cómo se calculan}: con la notación de la Encuesta de Población Activa, $L$ ocupados, $U$ parados, $A=L+U$ activos y $P_{16}$ población de 16 y más años \vertema{4.A.1},
\[
\text{tasa de actividad}=\frac{A}{P_{16}},\qquad\text{tasa de empleo}=\frac{L}{P_{16}},\qquad u=\frac{U}{A} .
\]
\item \emph{Cuidado con la edad}: el INE calcula las tasas sobre la población de 16 y más años (incluye a los jubilados), y Eurostat y la OCDE, sobre la de 15 a 64 años. Con la misma definición, la tasa de actividad de España (74,9~\% en 2025) ya no está por debajo de la media de la OCDE (73,9~\% en 2024); sobre la población de 16 y más años es del 58~\%, por el peso de los mayores \parencite{OCDE2026LFS,Eurostat2026lfsa,INE2025EPA}. Solo se comparan tasas con la misma definición.
\item \emph{Otras medidas}: la tasa bruta de actividad (activos sobre población total) mezcla la participación con la estructura de edades; el factor demográfico (población en edad de trabajar sobre la total) encaja en la descomposición $Y/N=(Y/L)(L/A)(A/P_{16})(P_{16}/N)$ \vertema{4.A.30}.
\item \emph{Umbral o referencia}: el cuadro de indicadores del procedimiento de desequilibrios macroeconómicos alerta con un paro superior al 10~\% y con una caída de la tasa de actividad de más de 0,2 puntos en 3 años \parencite{CE2025AMR2026}. Como regla práctica (orientativa), un paro del 3-6~\% suele considerarse pleno empleo en las economías avanzadas.\footnote{Pleno empleo no es paro nulo: siempre hay quien está cambiando de trabajo. En los países menos desarrollados un paro bajo puede esconder subempleo (pocas horas, informalidad, rentas muy bajas).}
\item \emph{Cómo se lee}: la tasa de actividad depende de la demografía, de la incorporación de la mujer, de los años de estudio y de los incentivos (jubilación, prestaciones, fiscalidad). El empleo y el paro pueden subir a la vez si llegan más activos (inmigración) de los que encuentran trabajo, y bajar a la vez si los parados emigran. En 2020 el empleo y el paro apenas reflejaron la caída del PIB por los esquemas de mantenimiento del empleo, que además no cuentan a sus trabajadores como parados.
\item \emph{Tema}: \vertema{3.A.25} \vertema{4.A.16}.
\end{itemize}

\paragraph{Productividad}
\begin{itemize}
\item \emph{Qué mide}: la \concepto{productividad aparente del trabajo}, $\mathrm{PAT}=Y/L$, es lo que produce cada ocupado (o cada hora, si el cuadro da horas); es una aproximación a la productividad marginal, que no se observa.
\item \emph{Cómo se calcula}: $\gamma_{\mathrm{PAT}}=\gamma_Y-\gamma_L$ (apartado~\ref{sec:oferta}). Con horas trabajadas, el PIB por persona en edad de trabajar se descompone en PIB por hora y horas por persona.
\item \emph{Cómo se lee}: crece por \emph{intensificación del capital} (más capital por trabajador, gracias a la inversión neta) o por \emph{progreso técnico} (sube la productividad total de los factores). Si cae el empleo menos productivo, la PAT sube sin que nadie trabaje mejor. A largo plazo, la productividad es la que eleva la renta per cápita y los salarios reales.
\item \emph{Tema}: \vertema{3.A.45} \vertema{4.A.29}.
\end{itemize}

\paragraph{Salarios, costes laborales unitarios y margen}
\begin{itemize}
\item \emph{Qué mide}: el \concepto{coste laboral unitario} (CLU) es lo que cuesta en salarios producir una unidad: la remuneración por asalariado ($W$, nominal) dividida por la productividad. Es la medida de competitividad en costes que más se usa.
\item \emph{Cómo se calcula}:
\[
\mathrm{CLU}=\frac{WL}{Y}=\frac{W}{\mathrm{PAT}}\quad\Longrightarrow\quad\gamma_{\mathrm{CLU}}=\gamma_W-\gamma_{\mathrm{PAT}} .
\]
Como $\mathrm{CLU}=P\cdot\frac{WL}{PY}$, el \concepto{CLU real} ($\mathrm{CLU}/P$, con $P$ el deflactor del PIB) es la participación de la remuneración de los asalariados en el PIB.\footnote{Hay 2 salarios reales: el \emph{de consumo}, $W/\mathrm{IPC}$, que mide el poder adquisitivo del trabajador, y el \emph{de producto}, $W/P$, que mide el coste para la empresa. El CLU real usa el segundo.}
\item \emph{Margen}: si el precio es un margen $\mu$ sobre el coste laboral unitario, $P=(1+\mu)\,\mathrm{CLU}$, y el margen evoluciona como
\[
\gamma_{1+\mu}\approx\gamma_P-\gamma_{\mathrm{CLU}} .
\]
Los márgenes suben exactamente cuando caen los CLU reales. Es una aproximación: supone constantes los impuestos indirectos netos y el peso de los consumos intermedios importados. El excedente bruto de explotación es otra cosa, un nivel de rentas, y su crecimiento no es el del margen.
\item \emph{Umbral o referencia}: el cuadro de indicadores del procedimiento de desequilibrios alerta si los CLU nominales por hora crecen más del 9~\% en 3 años en la zona del euro (12~\% fuera) \parencite{CE2025AMR2026}. Como referencia, que los salarios nominales crezcan como la productividad más el objetivo de inflación mantiene estables los CLU reales.\footnote{Es la «regla de oro» salarial de \textcite{Watt2007WageSetting}: en la zona del euro, unos CLU nominales que crecen en torno al 2~\% no son una pérdida de competitividad frente al resto del euro. Varios estudios discuten, además, que los CLU agregados midan bien la competitividad exterior: \textcite{Kaldor1978} observó que los países cuyos CLU relativos más crecieron entre 1963 y 1975 fueron los que más cuota de exportación ganaron (\concepto{paradoja de \autor{Kaldor}}), y \textcite{FelipeKumar2011} no encuentran en la zona del euro relación entre el crecimiento de los CLU y el del PIB. Una explicación sencilla de los CLU está en \textcite{UCM2015CLU}.}
\item \emph{Cómo se lee}:
  \begin{itemize}
  \item CLU nominales que crecen \emph{más que la inflación} (CLU reales al alza): los salarios crecen más que la productividad o la productividad cae sin reflejo en los salarios. Anticipan presiones inflacionistas, menos margen y beneficio y, a medio plazo, sustitución de trabajo por capital y menos empleo.
  \item CLU nominales que crecen \emph{más que en los socios}: pérdida de competitividad-precio (apartado~\ref{sec:tipo-cambio}).
  \item CLU que crecen poco con paro alto: ajuste por salarios (devaluación interna).
  \end{itemize}
\item \emph{Tema}: \vertema{3.A.27} \vertema{3.A.28} \vertema{4.A.22}.
\end{itemize}

\paragraph{Demanda y oferta de trabajo}
\begin{itemize}
\item \emph{Demanda de trabajo} (la creación de empleo): depende negativamente de los CLU \emph{reales} (si los nominales suben al ritmo de los precios, el trabajo no se encarece para la empresa) y positivamente de los márgenes y del excedente empresarial, de la demanda interna y de las exportaciones; los factores institucionales (rigideces, problemas de emparejamiento) la condicionan.
\item \emph{Oferta de trabajo} (los activos): depende del salario real y de la demografía y la inmigración, y la reducen las prestaciones mal diseñadas.
\item \emph{Cómo se escribe}: «el crecimiento del PIB parece correlacionar bien con la creación de empleo; el persistente aumento de los CLU y la caída de los márgenes podrían explicar las tímidas tasas de creación de empleo».
\end{itemize}

\paragraph{Qué tipo de paro hay}
\begin{itemize}
\item La receta depende del diagnóstico, así que hay que identificar el tipo de paro:
  \begin{itemize}
  \item \concepto{clásico}: salarios reales demasiado altos para la productividad (rigideces: salario mínimo, cuña fiscal, poder de negociación); se delata por CLU que crecen mucho y de forma sistemática, o por salarios reales que suben mientras sube el paro;
  \item \concepto{keynesiano}: demanda insuficiente; consumo y demanda interna anémicos;
  \item \concepto{friccional} o de desajuste: búsqueda y emparejamiento; vacantes altas con paro alto (curva de \autor{Beveridge} alejada del origen);
  \item \concepto{de larga duración}: se retroalimenta (pérdida de cualificación, desánimo, estigma) y es causa de histéresis.
  \end{itemize}
\item \concepto{NAIRU}: la tasa de paro compatible con una inflación estable. Si el paro está por debajo, la inflación tiende a acelerarse; si está por encima, a frenarse. No es un techo al empleo: mide a partir de qué paro aparecen tensiones salariales \vertema{3.A.28}. En el examen se aproxima con el paro de los años en que la brecha de producción es nula. La Comisión estima una NAWRU en torno al 6-7~\% en la zona del euro y al 12~\% en España \parencite{CE2026AMECO}; si es alta, puede haber rigideces (cuña fiscal alta, negociación que protege a los ocupados, prestaciones mal diseñadas, desajuste de cualificaciones).
\item \concepto{Histéresis}: el paro de equilibrio depende del paro pasado, de modo que un paro alto y persistente tiende a quedarse. Tiene 3 causas: la dualidad entre \emph{insiders} y \emph{outsiders}, el paro de larga duración y la pérdida de capital instalado. Si la hay, las políticas de demanda tienen efectos a largo plazo \vertema{3.A.28}.
\item \emph{Ley de \autor{Okun} y curva de \autor{Phillips}} (2021):
  \begin{itemize}
  \item la ley de \autor{Okun} relaciona la variación del paro con el crecimiento por encima del potencial; en EE.~UU., en torno a medio punto menos de paro por cada punto de crecimiento por encima del potencial \parencite{Okun1962}. En el cuadro sirve para ver a partir de qué crecimiento baja el paro (si el PIB crece mucho y el empleo poco, la elasticidad es baja);
  \item la curva de \autor{Phillips} con expectativas, $\pi_t=\pi^e_t-a\,(U_t-U^n)+\varepsilon_t$, relaciona la inflación con la distancia entre la tasa de paro y la NAIRU \parencite{Phillips1958}\footnote{En la curva de \autor{Phillips}, $U_t$ es la \emph{tasa} de paro y $U^n$ la NAIRU, como en \vertema{3.A.28}; en las fórmulas de la EPA, $U$ son los parados y $u=U/A$ la tasa.} \vertema{3.A.28};
  \item son razonables con un mercado de trabajo formal y asalariado y una serie larga y estable; no lo son con mucha informalidad, con esquemas de mantenimiento del empleo, con choques de oferta o cambios de régimen, y el coeficiente de \autor{Okun} varía mucho entre países y en el tiempo \parencite{Ball2017Okun}.
  \end{itemize}
\end{itemize}

\paragraph{Curva de \autor{Beveridge}}
\begin{itemize}
\item \emph{Qué mide}: la relación entre la tasa de vacantes y la de paro.
\item \emph{Cómo se lee}: moverse \emph{a lo largo} de la curva es el ciclo (en las expansiones, más vacantes y menos paro); \emph{desplazarse hacia fuera} indica peor emparejamiento entre parados y puestos (desajuste de cualificaciones o geográfico, paro de larga duración) \parencite{BlanchardDiamond1989}. Si las vacantes suben y el paro no baja, el problema es estructural, no de demanda (2023).
\item \emph{Tema}: \vertema{3.A.26}.
\end{itemize}

\begin{figure}[H]
\centering
\caption{Curva de \autor{Beveridge} de España y de la zona del euro}\label{fig:beveridge}
\grafico[0.7\textwidth]{beveridge-espana-zona-euro}
\fuente{\textcite{Eurostat2026jvs,Eurostat2026une}; medias anuales. Qué hay que ver: la zona del euro se mueve a lo largo de su curva con el ciclo (máximo de vacantes, 3,2~\%, con paro mínimo en 2022); España tiene muy pocas vacantes (en torno al 0,9~\%) con mucho paro, una curva lejos del origen que apunta a problemas de emparejamiento.}
\end{figure}

\paragraph{Demografía, migración y salario mínimo}
\begin{itemize}
\item \emph{Demografía} (2026): el envejecimiento (más población de 65 y más años por cada persona de 15 a 64) reduce el crecimiento del empleo y del PIB potencial si no lo compensan la productividad, la participación o la inmigración; crea escasez de mano de obra en algunos sectores (vacantes altas con paro bajo) y eleva el gasto en pensiones, sanidad y cuidados \parencite{CE2024Envejecimiento}. Medidas de oferta: más actividad de mujeres y mayores, inmigración ordenada e integrada, capital humano y automatización, y pensiones ligadas a la esperanza de vida \vertema{4.A.2}.
\item \emph{Un endurecimiento de los visados de trabajo} (2020): menos oferta de trabajo, más tensión salarial en los sectores con más inmigrantes, más vacantes sin cubrir, menos PIB potencial y, si los inmigrantes eran contribuyentes netos, peor saldo público. Al revés, la inmigración eleva primero el paro y después el crecimiento potencial.
\item \emph{Una subida del salario mínimo} (2019) sube los salarios bajos y comprime la distribución, eleva el consumo (alta propensión a consumir) y los CLU de los sectores intensivos en salarios bajos (algo de inflación y de pérdida de competitividad). Su efecto sobre el empleo depende del poder de mercado de los empleadores (con monopsonio puede no reducirlo) y del tamaño de la subida respecto al salario mediano: la evidencia encuentra efectos pequeños con subidas moderadas \parencite{CardKrueger1994,Manning2021SMI} \vertema{3.A.17}. El saldo público mejora por las cotizaciones de quienes conservan el empleo y empeora por los que lo pierden.
\end{itemize}

\subsection{Mercado de capitales}\label{sec:capitales}

\begin{ideaclave}
\textbf{Qué hay que contestar}: ¿está endeudada la economía, y con quién? ¿En qué punto del ciclo financiero está? ¿Es el sistema bancario solvente, líquido y rentable? ¿Hay burbujas en la vivienda o en la bolsa? ¿Qué debe hacer la política macroprudencial?
\end{ideaclave}

\begin{itemize}
\item Un sistema financiero que funciona bien canaliza el ahorro hacia la inversión productiva; uno que funciona mal frena la productividad por 2 vías: da crédito a empresas \emph{zombis} (con pérdidas, que sobreviven solo gracias al crédito) y deja sin financiar proyectos rentables. El orden de análisis: endeudamiento y ciclo financiero, banca, vivienda, bolsa y política macroprudencial \vertema{3.B.27} \vertema{4.A.19}.
\end{itemize}

\paragraph{Endeudamiento y ciclo financiero}
\begin{itemize}
\item \emph{Qué se mira}: la deuda de los hogares, de las empresas y del Estado (esta, en la pregunta de política fiscal, apartado~\ref{sec:fiscal}) y el flujo de crédito; si gran parte es frente al exterior (posición de inversión internacional muy negativa), la economía es más vulnerable a un encarecimiento de la financiación. El desarrollo del mercado financiero se aproxima por el diferencial entre los tipos de los préstamos y los de los depósitos.
\item \emph{Umbral}: el cuadro de indicadores del procedimiento de desequilibrios alerta con una deuda de los hogares superior al 55~\% del PIB, una de las sociedades no financieras superior al 85~\% y flujos de crédito anuales de más del 14~\% (hogares) y el 13~\% (empresas) de su deuda del año anterior \parencite{CE2025AMR2026,Eurostat2025MIP}.
\item \concepto{Brecha crédito-PIB}: la desviación de la ratio crédito/PIB respecto a su tendencia, la versión financiera de la brecha de producción. Positiva y creciente, auge del crédito (riesgo de sobreendeudamiento y burbujas); negativa, el crédito está por debajo de su tendencia. Por encima de 2 puntos se recomienda activar el colchón de capital anticíclico \parencite{BCBS2010Anticiclico} \vertema{3.A.42}.
\item \emph{Préstamos en moneda extranjera}: si una parte grande del crédito está en divisas y los prestatarios cobran en moneda nacional, una depreciación encarece sus deudas, eleva la morosidad y puede convertir una crisis cambiaria en bancaria \parencite{EHP2005PecadoOriginal}.
\end{itemize}

\paragraph{Sector bancario}
\begin{itemize}
\item \emph{Solvencia} (capacidad de absorber pérdidas con recursos propios): recursos propios sobre activos ponderados por riesgo (APR), con 3 ratios según la calidad del capital: el \concepto{CET1} (capital y reservas), el \concepto{Tier~1} (más instrumentos híbridos) y el \concepto{capital total} (más deuda subordinada). Los mínimos de Basilea~III, con el colchón de conservación, son el 7, el 8,5 y el 10,5~\% de los APR, a los que se suman los colchones anticíclico y sistémicos y el requisito de Pilar~2 de cada banco: en la zona del euro, los requisitos totales rondan el 11~\% de CET1 y el 16~\% de capital total \parencite{BCBS2017Basel3,RegUE2013CRR,BCE2025SREP}. Se completa con la \concepto{ratio de apalancamiento} (Tier~1 sobre la exposición total sin ponderar), con un mínimo del 3~\%.
\item \emph{Calidad de los activos}: \concepto{ratio de dudosos} (préstamos con más de 90 días de impago o de pago improbable sobre el crédito total) y \concepto{ratio de cobertura} (provisiones sobre dudosos). La morosidad va retrasada respecto al ciclo.
\item \emph{Liquidez} (capacidad de atender los pagos sin vender activos con pérdidas; un banco solvente puede caer por falta de liquidez): \concepto{LCR} (activos líquidos sobre las salidas netas en 30 días de tensión) y \concepto{NSFR} (financiación estable disponible sobre la requerida a un año), con mínimos del 100~\% \parencite{BCBS2013LCR,BCBS2014NSFR}, y la ratio créditos/depósitos: por encima del 100~\%, el banco depende de los mercados mayoristas.
\item \emph{Rentabilidad}: \concepto{ROE} (beneficio neto sobre patrimonio neto) y ROA bancario (beneficio neto sobre activo medio, distinto del ROA de la empresa no financiera \vertema{3.B.1}), con $\mathrm{ROE}=\mathrm{ROA}\times\text{activo}/\text{patrimonio}$; y \concepto{ratio de eficiencia} (gastos de explotación sobre margen bruto; cuanto menor, mejor). El ROE es suficiente si supera el \concepto{coste del capital}, la rentabilidad que exigen los accionistas (el tipo sin riesgo más una prima), que los bancos europeos estiman entre el 8 y el 12~\% \parencite{ABE2025RAR}; no basta compararlo con los tipos de mercado.
\item \emph{Referencias orientativas} (no regulatorias) de la Autoridad Bancaria Europea: CET1 holgado por encima del 14~\%; dudosos bajos por debajo del 3~\% y altos por encima del 8~\%; ROE favorable por encima del 10~\%; eficiencia favorable por debajo del 50~\% y desfavorable por encima del 60~\% \parencite{ABE2026RiskDashboard}. A mediados de 2026 la banca española era más rentable que la europea (ROE del 16~\%, frente al 11~\%) pero tenía menos capital de máxima calidad (CET1 del 14~\%, frente al 16~\%) \parencite{BdE2026EstadSupervisoras,ABE2026RiskDashboard}.
\item \emph{Cómo se lee}: las debilidades se refuerzan entre sí: más dudosos, más provisiones, menos beneficio y menos capital. Con los tipos, 2 efectos opuestos: tipos más bajos estrechan el margen de intermediación (efecto precio) pero aumentan el crédito y mejoran la calidad de los préstamos (efecto volumen); con subidas de tipos, al revés.
\end{itemize}

\paragraph{Precio de la vivienda}
\begin{itemize}
\item \emph{Qué mide}: el valor del principal activo de los hogares y de la garantía de buena parte del crédito bancario (se preguntó en 2016 y 2020).
\item \emph{Cómo se lee}: hay que separar fundamentos y burbuja.
  \begin{itemize}
  \item \emph{Fundamentos}: renta disponible, población y formación de hogares, tipos hipotecarios, oferta de vivienda nueva, fiscalidad, demanda extranjera.
  \item \emph{Señales de sobrevaloración}: precio que crece mucho más que la renta o que los alquileres, crédito hipotecario muy dinámico con criterios relajados (préstamo sobre valor o sobre renta altos) y precios que suben más que el VAB de la construcción (más precio que cantidad).
  \end{itemize}
\item \emph{Umbral}: el cuadro de indicadores del procedimiento de desequilibrios alerta con subidas del precio nominal de más del 9~\% anual \parencite{CE2025AMR2026}. Para la accesibilidad, como referencia orientativa, Eurostat habla de sobrecarga cuando los gastos de vivienda superan el 40~\% de la renta disponible \parencite{Eurostat2025Sobrecarga}.
\item \emph{Con qué se cruza}: consumo (efecto riqueza), inversión residencial, crédito y morosidad, y movilidad laboral (vivienda cara, peor emparejamiento) \vertema{4.A.11}.
\end{itemize}

\paragraph{Índice bursátil}
\begin{itemize}
\item \emph{Qué mide}: las expectativas de beneficios empresariales descontadas con una prima de riesgo; es un indicador adelantado del ciclo.
\item \emph{Cómo se lee}: si sube con el PIB y los beneficios, refleja buenas perspectivas; si sube mucho más que los beneficios con tipos muy bajos o con fuertes entradas de inversión de cartera, puede haber sobrevaloración. Se cruza con la inversión (la $q$ de \autor{Tobin}: si el valor de mercado de las empresas supera el coste de reponer su capital, conviene invertir, \vertema{3.A.34}) y con el consumo (efecto riqueza, menor que el de la vivienda).
\end{itemize}

\paragraph{Política macroprudencial}
\begin{itemize}
\item \emph{Objetivo}: limitar el riesgo \emph{sistémico}, el del conjunto del sistema, que la supervisión de cada banco no ve; tiene una dimensión temporal (el ciclo financiero) y otra estructural (bancos grandes e interconectados) \parencite{CERS2014Manual}.
\item \emph{En la fase alcista} (crédito y vivienda muy dinámicos, brecha crédito-PIB positiva): activar el colchón anticíclico (hasta el 2,5~\%) o uno de riesgo sistémico sectorial, y limitar a los prestatarios (préstamo sobre valor, sobre renta y servicio de la deuda sobre renta).
\item \emph{En la fase bajista o en una crisis}: liberar los colchones para que el capital absorba pérdidas sin cortar el crédito, dar flexibilidad supervisora temporal y limitar dividendos para retener capital, como en 2020 \parencite{BCE2020Dividendos}; y vigilar el precio de los activos.
\item \emph{Tipos muy bajos durante mucho tiempo}: pueden ser un riesgo (búsqueda de rentabilidad, proyectos de baja calidad, menor margen de bancos y aseguradoras, activos inflados).
\item \emph{Tema}: \vertema{3.B.27} \vertema{3.B.44} \vertema{4.A.19}.
\end{itemize}

\begin{itemize}
\item \emph{Recapitulación}: el mercado de trabajo dice cuánto empleo crea el crecimiento y si los costes siguen a la productividad (CLU); el de capitales, si el crédito y la banca lo sostienen sin acumular riesgos. Los CLU llevan directamente al bloque siguiente: la competitividad.
\end{itemize}

% ==========================================================================
\section{Competitividad}\label{sec:competitividad}

\begin{itemize}
\item La conclusión por adelantado: la competitividad tiene 2 caras. La de \emph{precio}, que depende de la inflación, los costes y el tipo de cambio, y la de \emph{no precio}, que depende de la calidad, la diferenciación y los mercados. Una economía puede perder la primera y ganar cuota por la segunda.
\end{itemize}

\subsection{Competitividad precio: precios e inflación}\label{sec:precios}

\begin{ideaclave}
\textbf{Qué hay que contestar}: ¿está controlada la inflación? ¿Viene de dentro o de fuera, de la demanda o de los costes? ¿Es transitoria o persistente? ¿Qué política requiere?
\end{ideaclave}

\paragraph{Índices de precios}
\begin{itemize}
\item \emph{Qué miden}: la \concepto{inflación} es el aumento sostenido y generalizado del nivel de precios, no el cambio de un precio relativo; la \concepto{deflación}, su caída. Se mide con la tasa de variación de un índice: $\pi_t=(\mathrm{IPC}_t/\mathrm{IPC}_{t-1}-1)\times100$.\footnote{El IPC y el deflactor son números índice (100 en el año base); si el cuadro da un porcentaje, no es el índice, sino su tasa de variación.}
\item \emph{Los más usados}:
  \begin{itemize}
  \item \concepto{IPC}: los precios de la cesta que compran los hogares, \emph{importaciones incluidas}; en la UE, el armonizado (IAPC) es la referencia del objetivo del BCE;
  \item \concepto{deflactor del PIB}: el cociente entre el PIB nominal y el PIB en volumen; recoge los precios de lo que se produce dentro, \emph{exportaciones, inversión y servicios públicos incluidos}, y excluye las importaciones;
  \item precios industriales, de exportación y de importación (los de importación anticipan la inflación importada), costes laborales (salario medio, CLU) y precios de activos (vivienda, bolsa, materias primas).
  \end{itemize}
\end{itemize}

\begin{anotaciones}
El IPC es un índice de tipo \autor{Laspeyres}: compara el coste de la cesta del periodo base, $\vect{x}^0$, a los precios de hoy y a los del periodo base, $I_L=\frac{\vect{p}^t\cdot\vect{x}^0}{\vect{p}^0\cdot\vect{x}^0}\times100$. Como no recoge la sustitución hacia los bienes que se abaratan, \emph{sobrevalora} la inflación (sesgo de sustitución). El del INE es un \autor{Laspeyres} \emph{encadenado}: cada año se compara con los precios de diciembre del año anterior y las ponderaciones se actualizan cada año (base 2025 desde enero de 2026) \parencite{INE2026IPCMetodologia}. El deflactor es de tipo \autor{Paasche}, con la cesta de hoy, $I_P=\frac{\vect{p}^t\cdot\vect{x}^t}{\vect{p}^0\cdot\vect{x}^t}\times100$, e \emph{infravalora} la inflación; como el PIB en volumen se calcula encadenado, es un \autor{Paasche} encadenado \vertema{3.A.9} \vertema{3.A.30} \vertema{4.A.1}.
\end{anotaciones}

\paragraph{IPC frente a deflactor e inflación importada}
\begin{itemize}
\item Suelen dar resultados parecidos; cuando se separan, dicen algo:
  \begin{itemize}
  \item IPC que crece \emph{más} que el deflactor: \concepto{inflación importada} (energía, alimentos, depreciación), porque las importaciones están en el IPC y no en el deflactor;
  \item IPC que crece \emph{menos}: suben más los precios de lo que el país produce y no consume, como las exportaciones (mejora de la relación real de intercambio, típica de un exportador de materias primas en un auge), los bienes de inversión o los servicios públicos. Ninguno de los 2 índices recoge los precios de los activos financieros.
  \end{itemize}
\item El deflactor aísla el efecto externo, pero refleja peor el poder adquisitivo; por eso el salario real del trabajador se deflacta con el IPC y el coste real para la empresa, con el deflactor.
\item Si los precios de importación crecen más que el IPC, se está importando inflación; si crecen menos, los bienes importados se abaratan frente a los nacionales, lo que resta competitividad a la producción nacional en su propio mercado.
\end{itemize}

\paragraph{Inflación subyacente y efectos de segunda ronda}
\begin{itemize}
\item \emph{Qué mide}: la \concepto{inflación subyacente} quita de la general los componentes más volátiles para ver la tendencia. El INE excluye los alimentos no elaborados y los productos energéticos; Eurostat y el BCE siguen el IAPC sin energía ni alimentos (incluidos alcohol y tabaco). Hay también medias recortadas y medidas que quitan cada mes los componentes más volátiles. El cuadro puede traer cualquiera: hay que decir cuál se usa.\footnote{No hay una base teórica única para medirla; todas las medidas buscan separar los cambios persistentes de los transitorios.}
\item \emph{Cómo se lee}: sirve para ver si un choque de oferta (energía) se está extendiendo al resto de la cesta. Son los \concepto{efectos de segunda ronda}: los salarios y los márgenes suben para recuperar lo perdido y alimentan una espiral de precios y salarios \vertema{3.A.41}.
\end{itemize}

\paragraph{De dónde viene la inflación}
\begin{itemize}
\item \emph{De demanda} (\emph{demand-pull}): la demanda interna crece más que el PIB, la brecha de producción es positiva, el crédito y los agregados monetarios crecen con fuerza, hay estímulo fiscal o las expectativas de inflación suben.
\item \emph{De costes} (\emph{cost-push}): suben los salarios y los CLU (inflación incorporada, \emph{built-in}), la energía y las materias primas, o se deprecia la moneda (más cuanto más pesan las importaciones); los márgenes amplios señalan poca competencia.
\item \concepto{Efecto base}: la tasa interanual compara con el mismo mes del año anterior; si entonces hubo una subida excepcional, la tasa baja cuando sale de la comparación aunque los precios no bajen. En logaritmos, $\pi^{12}_t-\pi^{12}_{t-1}\approx\pi^{m}_t-\pi^{m}_{t-12}$ (tasa mensual que entra menos la que sale). Se preguntó en 2023.
\item \emph{¿Transitoria o persistente?} Transitoria si se concentra en energía y alimentos con la subyacente contenida, salarios moderados, expectativas ancladas y materias primas que ya bajan. Persistente si la subyacente sube y se generaliza (servicios), los salarios y los márgenes se aceleran, la brecha es positiva, el paro está por debajo de la NAIRU y suben las expectativas.
\item \emph{Resistencia inicial y cambio de tendencia} (2025 y 2026): al principio, choques de oferta (energía, cuellos de botella) con demanda fuerte (ahorro acumulado, estímulos) y traspaso de la depreciación; después, la bajada de la energía, el efecto base, el endurecimiento monetario y la normalización de las cadenas de suministro.
\end{itemize}

\begin{figure}[H]
\centering
\caption{Inflación de la zona del euro y contribuciones de sus componentes}\label{fig:iapc-zona-euro}
\grafico[0.85\textwidth]{iapc-zona-euro}
\fuente{\textcite{Eurostat2026hicp}; tasas interanuales y contribuciones en puntos, enero de 2017-septiembre de 2026. Qué hay que ver: el brote de 2021-2022 empieza en la energía y se traslada a alimentos y bienes; la subyacente baja más despacio, por los servicios. En 2026 la energía vuelve a empujar (IAPC del 3,8~\% en septiembre) con la subyacente estable en el 2,5~\%.}
\end{figure}

\paragraph{Qué política requiere}
\begin{itemize}
\item Si la inflación es alta y persistente, lo natural es endurecer la política monetaria, pero depende del origen:
  \begin{itemize}
  \item \emph{de demanda}: subir tipos, y que la política fiscal no empuje en sentido contrario;
  \item \emph{de oferta y transitoria}, con expectativas ancladas: el banco central puede «mirar a través» del choque y vigilar los efectos de segunda ronda; si las expectativas se desanclan o la subyacente se contagia, tiene que subir tipos aunque el choque sea de oferta;
  \item en una unión monetaria el país no tiene política monetaria propia: le quedan la fiscal y las de oferta.
  \end{itemize}
\item \emph{Con deflación}: política monetaria expansiva; si los tipos están en su límite inferior, expectativas, compras de activos y política fiscal expansiva.
\item \emph{Siempre}: políticas que aumenten la competencia y reduzcan los márgenes.
\item El BCE combina el análisis económico con el análisis monetario y financiero, con un objetivo simétrico del 2~\% a medio plazo \parencite{BCE2025Estrategia} \vertema{3.A.36}.
\end{itemize}

\subsection{Tipo de cambio}\label{sec:tipo-cambio}

\begin{itemize}
\item \emph{Qué mide}: el precio de una moneda en otra; el tipo \emph{real} mide la competitividad-precio.
\item \emph{Cómo se calcula} \vertema{3.B.13}:
  \begin{itemize}
  \item \concepto{nominal bilateral} $E$, en convención directa: unidades de moneda nacional por unidad de moneda extranjera; si $E$ sube, la moneda nacional se \emph{deprecia};
  \item \concepto{efectivo}: media geométrica de los bilaterales ponderada por el comercio con cada socio;
  \item \concepto{real}: $q\equiv EP^*/P$; si sube, la producción nacional se abarata frente a la extranjera y la economía gana competitividad-precio. El efectivo real es $q^{ef}=\prod_i(E_iP^*_i/P)^{\omega_i}$; con los CLU en lugar de los precios se obtiene la competitividad en costes.
  \end{itemize}
\item \emph{Cuidado con las unidades}: el BCE y el Banco de España publican los índices al revés (una subida es una apreciación y una pérdida de competitividad) \vertema{4.A.22}, y el cuadernillo puede traer cualquiera de las 2 convenciones: en 2026 el tribunal tuvo que rectificar las unidades del tipo frente al dólar \parencite{ATCEE2026Examenes}.
\item \emph{Umbral}: el cuadro de indicadores del procedimiento de desequilibrios alerta con una apreciación real de más del 3~\% en 3 años en la zona del euro (10~\% fuera) \parencite{CE2025AMR2026}.
\item \emph{Cómo se lee}: la competitividad-precio mejora por depreciación nominal o por una inflación menor que la de los socios (devaluación interna).
\end{itemize}

\paragraph{Traspaso (\emph{pass-through})}
\begin{itemize}
\item \emph{Qué es}: la parte de una depreciación que pasa a los precios interiores, primero a los de importación y después al IPC. Un traspaso de 0,3 significa que una depreciación del 10~\% eleva los precios un 3~\%. Se preguntó en 2026.
\item \emph{Cómo se detecta en el cuadro}: comparando la depreciación con la inflación (sobre todo la de los bienes importados y la subyacente) de los años siguientes; si tras una depreciación fuerte la inflación apenas se mueve, el traspaso es bajo.
\item \emph{De qué depende}: es mayor con inflación alta y poco creíble, con mucho peso de las importaciones en el consumo, con importaciones facturadas en dólares y en los emergentes \parencite{Taylor2000PassThrough,CampaGoldberg2005}; menor si los exportadores extranjeros absorben la depreciación en sus márgenes, si se espera que sea transitoria y si el banco central es creíble. Con el comercio facturado en dólares, la depreciación encarece las importaciones pero apenas abarata las exportaciones a corto plazo \parencite{Gopinath2020DCP}.
\end{itemize}

\paragraph{Condición de \autor{Marshall}-\autor{Lerner} y curva J}
\begin{itemize}
\item Una depreciación real mejora la balanza comercial si, partiendo de comercio equilibrado, la suma de los valores absolutos de las elasticidades-precio de exportaciones e importaciones es mayor que 1 \parencite{Marshall1923,Lerner1944}:
\[
|\varepsilon_{X,q}|+|\varepsilon_{M,q}|>1 .
\]
\item A corto plazo las cantidades reaccionan poco y la balanza puede empeorar antes de mejorar: es la \concepto{curva J} \parencite{Magee1973}.
\item \emph{Cómo se comprueba en el cuadro} (2022): si tras una depreciación real el saldo de bienes y servicios mejora en 1 o 2 años, con más volumen exportado y menos importado, la condición parece cumplirse; si empeora de forma persistente, no. Con prudencia: la renta interior y la exterior también cambian \vertema{3.B.12}.
\end{itemize}

\paragraph{Régimen cambiario e intervención}
\begin{itemize}
\item \emph{Regímenes}: fijo (o caja de conversión), intermedios (bandas, deslizamiento, \emph{flotación sucia}, con intervenciones para suavizar) y flotación libre. Con libertad de capitales, un tipo fijo obliga a renunciar a la política monetaria propia (\concepto{trilema}) \vertema{3.B.15}.
\item \emph{Cómo se detecta la intervención}: por las reservas. Si caen mientras la moneda se deprecia, el banco central vende divisas para frenarla; si suben mientras se aprecia, las compra. Muchos emergentes que dicen flotar intervienen por \concepto{miedo a flotar} \parencite{CalvoReinhart2002}.
\item \emph{Efecto de comprar divisas sobre la inflación} (2021): el banco central las paga con moneda nacional; si no la retira vendiendo bonos (intervención \emph{no esterilizada}), aumentan la base monetaria y la presión inflacionista, y además frena la apreciación (más inflación importada). Si esteriliza, neutraliza el efecto monetario a costa de un coste financiero. El nivel adecuado de reservas se ve en el sector exterior (apartado~\ref{sec:financiacion}).
\end{itemize}

\subsection{Competitividad no precio y política comercial}\label{sec:no-precio}

\begin{itemize}
\item \emph{Qué mide}: la capacidad de vender fuera por razones distintas del precio: calidad, diferenciación e imagen de marca, contenido tecnológico, servicio posventa, diversificación de productos y mercados, número y tamaño de las empresas exportadoras.
\item \emph{Cómo se aproxima} (casi nunca viene en el cuadro): la \concepto{cuota de mercado} de las exportaciones (si crece aunque el tipo de cambio real se aprecie o los CLU relativos suban, la competitividad no precio es fuerte, como en la paradoja de \autor{Kaldor}); exportaciones que crecen más que la demanda de los socios; valores unitarios altos; peso de las manufacturas de tecnología media-alta y de los servicios no turísticos; concentración geográfica y sectorial de las exportaciones.
\item \emph{Umbral}: el cuadro de indicadores del procedimiento de desequilibrios alerta si los resultados de exportación caen más del 3~\% en 3 años frente a las economías avanzadas \parencite{CE2025AMR2026}.
\item \emph{Cómo se lee}: un país que pierde cuota sin perder competitividad-precio tiene un problema de oferta exportadora (especialización poco dinámica, empresas pequeñas, poca innovación) o de mercados de destino.
\item \emph{Recomendaciones}: I+D e innovación, más tamaño empresarial, diversificación de mercados, promoción comercial (red de oficinas comerciales, agencia de promoción de exportaciones, financiación y seguro de crédito a la exportación) y nuevos acuerdos comerciales \vertema{3.B.6} \vertema{4.A.22}.
\end{itemize}

\paragraph{Aranceles y fragmentación del comercio}
\begin{itemize}
\item \emph{Medir la protección}: el \concepto{arancel medio simple} da el mismo peso a cada partida (sencillo y comparable, pero cuenta igual una partida que apenas se importa); el \concepto{medio ponderado por el comercio} pesa por lo que se importa (más representativo, pero infravalora los aranceles muy altos, que reducen precisamente lo que se importa).
\item \emph{Efecto de los aranceles de un tercero} sobre el país (2026, por la política arancelaria de Estados Unidos):
  \begin{itemize}
  \item \emph{directo}: caen las exportaciones a ese destino, según su peso y la elasticidad de cada producto;
  \item \emph{indirecto}: por las cadenas de valor y la menor demanda mundial;
  \item \emph{desviación de comercio}: productos de terceros que ya no entran en el mercado protegido llegan a otros, abaratan las importaciones y presionan a los productores locales;
  \item \emph{incertidumbre}, que retrasa la inversión, y \emph{tipo de cambio}, que cambia la competitividad y la inflación.
  \end{itemize}
\item \emph{Cómo se lee en el cuadro}: peso del destino en las exportaciones, evolución reciente de esas exportaciones, inflación de los bienes importados, confianza y PMI de manufacturas. En la guerra comercial de 2018-2019, los aranceles de EE.~UU. se trasladaron casi por completo a los precios que pagaron sus importadores y consumidores \parencite{Amiti2019Aranceles,Fajgelbaum2020Proteccionismo}.
\item \emph{Integración comercial asimétrica}: con más apertura ganan los sectores con ventaja comparativa y pierden los manufactureros menos productivos; se recomiendan plazos de desarme, ayudas a la reasignación y formación \vertema{3.B.7} \vertema{3.B.33} \vertema{3.B.41}.
\end{itemize}

\begin{itemize}
\item \emph{Recapitulación}: la competitividad precio se lee con la inflación (de dónde viene y si persiste), los CLU y el tipo de cambio real (con cuidado con la convención); la no precio, con las cuotas de exportación. Las 2 se reflejan en el saldo exterior, que es el bloque siguiente.
\end{itemize}

% ==========================================================================
\section{Sector exterior}\label{sec:exterior}

\begin{itemize}
\item La conclusión por adelantado: el sector exterior se analiza en 4 preguntas encadenadas: ¿hay capacidad o necesidad de financiación y por qué?, ¿cómo se financia?, ¿qué posición acumula el país frente al exterior? y ¿es sostenible? Ser deudor no es malo en sí: lo que importa es la sostenibilidad (enfoque intertemporal, \vertema{3.B.12}).
\item Los desequilibrios externos suelen ser síntoma de desequilibrios internos (exceso de gasto, poco ahorro, pérdida de competitividad).
\end{itemize}

\begin{ideaclave}
\textbf{Qué hay que contestar}: ¿hay capacidad o necesidad de financiación y de qué subcuenta viene? ¿Es un problema de ahorro o de inversión, del sector público o del privado? ¿Cómo se financia? ¿Es sostenible la posición exterior?
\end{ideaclave}

\subsection{Capacidad o necesidad de financiación}\label{sec:cuenta-corriente}

\paragraph{La balanza de pagos}
\begin{itemize}
\item \emph{Qué es}: el registro de las transacciones de los residentes con el resto del mundo en un periodo. Es una variable \emph{nominal} (no se compara con variables reales) y está equilibrada en conjunto, aunque sus subcuentas no.
\item \emph{Cómo se escribe} (MBP6, \vertema{3.B.11}):\footnote{Con el sexto manual del FMI \parencite{FMI2009MBP6}, un aumento de los activos netos frente al exterior (salida de capitales) se anota en la cuenta financiera con signo positivo, y una necesidad de financiación sale negativa. También se admite escribir $\mathrm{CC}+\mathrm{CK}+\mathrm{EO}-\mathrm{CF}=0$, que es lo mismo. El FMI publicó en 2025 el séptimo manual, que se aplicará hacia 2029-2030; los datos oficiales siguen en el sexto \parencite{FMI2025MBP7}.}
\begin{equation}\label{eq:identidad-bp}
\underbrace{\mathrm{CC}+\mathrm{CK}}_{\text{capacidad (+) o necesidad ($-$) de financiación}}+\ \mathrm{EO}=\mathrm{CF}=\mathrm{ID}+\mathrm{IC}+\mathrm{D}+\mathrm{OI}+\Delta R ,
\end{equation}
con $\mathrm{CC}=(X-M)+\mathrm{RP}+\mathrm{RS}$ (bienes y servicios, rentas primarias y rentas secundarias), $\mathrm{CK}$ la cuenta de capital, $\mathrm{EO}$ errores y omisiones, y en la cuenta financiera la inversión directa, la de cartera, los derivados, la otra inversión y la variación de las reservas.
\item \emph{Por subcuentas}:
  \begin{itemize}
  \item \emph{bienes y servicios}: competitividad, renta exterior, precios de la energía; un déficit de servicios con mucho turismo puede indicar poca competitividad de los servicios profesionales;
  \item \emph{rentas primarias}: trabajo y, sobre todo, rentas de la inversión (intereses de la deuda externa, beneficios de la inversión extranjera): explican la diferencia entre el PIB y la RNB (apartado~\ref{sec:medidores});
  \item \emph{rentas secundarias}: transferencias sin contrapartida (remesas de emigrantes, muy estables y contracíclicas en los países en desarrollo; cooperación; aportaciones a la UE);
  \item \emph{cuenta de capital}: transferencias de capital (fondos europeos para inversión) y compraventa de activos no producidos; pequeña, pero hay que sumarla.\footnote{Las subvenciones del Mecanismo de Recuperación y Resiliencia se anotan en la renta secundaria o en la cuenta de capital según financien gasto corriente o inversión, y los préstamos, como pasivos en la otra inversión de las administraciones públicas \parencite{BdE2021NGEU}.}
  \end{itemize}
\item \emph{Umbral}: el cuadro de indicadores del procedimiento de desequilibrios alerta con un saldo corriente medio de 3 años inferior al $-4$~\% del PIB o superior al $+6$~\% \parencite{CE2025AMR2026}.
\end{itemize}

\paragraph{Apertura, penetración de importaciones y relación real de intercambio}
\begin{itemize}
\item \emph{Apertura}, $(X+M)/\mathrm{PIB}$ a precios corrientes. Órdenes de magnitud orientativos, sin norma detrás: por debajo del 15-20~\%, economía cerrada (país grande); en torno al 50~\%, normal; por encima del 75~\%, muy abierta (país pequeño). EE.~UU. está en torno al 25~\%, España en torno al 70~\% y Bélgica por encima del 150~\% \parencite{Eurostat2026nama10gdp,BM2026WDI} \vertema{3.B.9}.
\item \emph{Penetración de las importaciones}, $M/\mathrm{PIB}$: un salto brusco apunta a una demanda que la producción interior no cubre (sobrecalentamiento); una caída brusca, a una recesión o a una depreciación.
\item \emph{Relación real de intercambio}, $\mathrm{RRI}=P_X/P_M$ con los deflactores de exportaciones e importaciones: si sube, con las mismas exportaciones se compran más importaciones (ganancia de renta), salvo que suba porque los precios de exportación se encarecen y se pierde competitividad. Para la competitividad es mejor el tipo de cambio real (apartado~\ref{sec:tipo-cambio}) \vertema{3.B.5}.
\end{itemize}

\paragraph{Ahorro, inversión y saldos sectoriales}
\begin{itemize}
\item \emph{Cómo se calcula}: de la identidad de la renta, $\mathrm{CC}=S-I=(S_P-I_P)+(T-G)$: un déficit corriente es un exceso de inversión sobre el ahorro nacional. Por sectores institucionales, la capacidad o necesidad de financiación de la nación es la suma de la de los hogares, las empresas, las instituciones financieras y las administraciones públicas \parencite{SEC2010}.
\item \emph{Cómo se lee}:
  \begin{itemize}
  \item un déficit por \emph{poco ahorro} (consumo excesivo, déficit público) es peor que uno por \emph{mucha inversión} productiva, que eleva la capacidad de pago futura;
  \item si viene del \emph{sector privado} (hogares y empresas que invierten con crédito), el riesgo es una corrección brusca cuando se corte el crédito; si viene del \emph{público}, es un problema de sostenibilidad fiscal;
  \item en las crisis el sector privado suele pasar a ahorrar y el público a endeudarse: los saldos se compensan (se preguntó en 2018);
  \item un \emph{superávit} puede reflejar demanda interna débil (envejecimiento, desapalancamiento, consolidación fiscal), inversión baja (un sistema financiero que no canaliza el ahorro) o ingresos extraordinarios que se ahorran (exportador de materias primas en un auge);
  \item la \emph{volatilidad} de la capacidad de financiación (2023) suele venir de la balanza energética, de los beneficios de las multinacionales o de los fondos europeos.
  \end{itemize}
\item \concepto{Déficits gemelos}: si el saldo privado no cambia, más déficit público es más déficit exterior. La causalidad iría por 2 vías \vertema{3.A.32}: $\downarrow(T-G)\Rightarrow\uparrow i\Rightarrow$ entrada de capitales $\Rightarrow\downarrow E$ (apreciación) $\Rightarrow\downarrow(X-M)$; y $\downarrow(T-G)\Rightarrow\uparrow Y\Rightarrow\uparrow M\Rightarrow\downarrow(X-M)$. En los datos no siempre van juntos: si el sector privado ahorra más cuando el Estado gasta más, el déficit exterior puede corregirse.
\item Si el cuadro no da el ahorro nacional pero sí el de los hogares o el privado, se dice y se razona con él suponiendo que el nacional sigue la misma tónica.
\end{itemize}

\subsection{Cómo se financia}\label{sec:financiacion}

\begin{itemize}
\item \emph{Qué se mira}: la cuenta financiera por instrumentos, porque no es lo mismo financiarse con capital que con deuda:
  \begin{itemize}
  \item \emph{inversión directa}: estable, con transferencia de tecnología y efectos desbordamiento;
  \item \emph{inversión de cartera}: volátil y con efecto rebaño ante un cambio de expectativas;
  \item \emph{derivados}: aún más volátiles;
  \item \emph{otra inversión}: préstamos y depósitos (incluidos los del Eurosistema por TARGET y los de organismos como el FMI o la UE);
  \item \emph{reservas}: su variación es moderada con tipo de cambio flexible y grande si el banco central interviene (apartado~\ref{sec:tipo-cambio}).
  \end{itemize}
\item \emph{Adecuación de las reservas}: los indicadores clásicos son de cobertura \parencite{FMI2011Reservas}:
  \begin{itemize}
  \item al menos 3 meses de importaciones (países con poco acceso a los mercados);
  \item al menos el 100~\% de la deuda externa a corto plazo (regla de \autor{Greenspan}-\autor{Guidotti}, para emergentes con acceso a los mercados; mide la liquidez, no la solvencia);
  \item al menos el 20~\% del dinero en sentido amplio, M2 (riesgo de fuga de depósitos hacia divisas, sobre todo con tipo fijo o caja de conversión);
  \item el FMI combina estos criterios en una métrica compuesta para los emergentes.
  \end{itemize}
  Demasiadas reservas también cuestan (rinden poco y su esterilización tiene un coste). Se preguntó en 2021 y 2024.
\item \emph{Tema}: \vertema{3.B.11} \vertema{3.B.22} \vertema{4.A.21}.
\end{itemize}

\subsection{Posición de inversión internacional y deuda externa}\label{sec:piin}

\begin{itemize}
\item \emph{Qué mide}: los saldos (\emph{stocks}) que acumulan los flujos. La \concepto{posición de inversión internacional neta} (PIIN) es la diferencia entre los activos y los pasivos financieros frente al exterior,
\[
\mathrm{PIIN}\equiv\text{activos exteriores}-\text{pasivos exteriores},
\]
negativa si el país es \emph{deudor neto} \parencite{FMI2009MBP6}. La \concepto{deuda externa} son solo los pasivos exigibles (sin las acciones, la inversión directa en capital ni los derivados; sí los bonos de la inversión de cartera) \vertema{3.B.11}.
\item \emph{Umbral}: el cuadro de indicadores del procedimiento de desequilibrios alerta con una PIIN inferior al $-35$~\% del PIB \parencite{CE2025AMR2026}.
\item \emph{Cómo se lee}: primero, si el país es acreedor o deudor; después, cómo evoluciona y por qué; y la composición de los pasivos, que importa tanto como el nivel:\footnote{Hay debate sobre la mejor métrica de vulnerabilidad: \textcite{Shin2012} defiende los pasivos exteriores \emph{brutos}, mejor indicador de la integración financiera y de la exposición; \textcite{CataoMilesiFerretti2014}, que los \emph{netos} (la PIIN) predicen bien las crisis exteriores. La deuda externa, al recoger solo lo exigible, refleja mejor el riesgo de refinanciación.}
  \begin{itemize}
  \item \emph{instrumentos}: mejor no exigibles (inversión directa antes que cartera y que derivados);
  \item \emph{sectores}: los pasivos del banco central no tienen riesgo de refinanciación;
  \item \emph{plazos}: mejor largos;
  \item \emph{monedas}: mejor en moneda nacional.
  \end{itemize}
\item \emph{Efecto valoración}: la PIIN cambia por las transacciones y por las \emph{otras variaciones} (precios de los activos y tipo de cambio). Un empeoramiento de la PIIN por valoración no siempre es más vulnerabilidad: si suben los precios de las acciones y los bonos nacionales en manos de no residentes porque mejoran las perspectivas, los pasivos valen más (se preguntó en 2016 y 2018).
\end{itemize}

\subsection{¿Es sostenible?}\label{sec:sostenibilidad-exterior}

\begin{itemize}
\item \emph{Dinámica}: en porcentaje del PIB, con $\gamma^N$ el crecimiento del PIB nominal y $\mathrm{OV}$ las otras variaciones \vertema{3.B.12},
\begin{equation}\label{eq:dinamica-piin}
\mathrm{PIIN}_t-\mathrm{PIIN}_{t-1}\approx\mathrm{CF}_t+\mathrm{OV}_t-\frac{\gamma^N_t}{1+\gamma^N_t}\,\mathrm{PIIN}_{t-1} .
\end{equation}
Con PIIN negativa, el crecimiento del PIB nominal la \emph{mejora} aunque el país siga pidiendo prestado: por eso la PIIN puede mejorar con cuenta financiera negativa (2018). Los intereses ya están en la cuenta corriente: no hay que sumarlos aparte.
\item \emph{De qué depende la sostenibilidad}:
  \begin{itemize}
  \item \emph{volumen de la deuda} ($-$): más deuda, más intereses, y el problema se retroalimenta;
  \item \emph{crecimiento} ($+$): de ahí la importancia de en qué se usa la financiación (inversión productiva mejor que gasto corriente);
  \item \emph{tipos de interés} ($-$);
  \item \emph{inflación} (ambigua): solo licúa la deuda emitida en moneda nacional;
  \item \emph{tipo de cambio} ($-$ si la deuda está en divisas): una depreciación la encarece; los emergentes que no pueden endeudarse fuera en su moneda (\concepto{pecado original}) son los más expuestos \parencite{EHP2005PecadoOriginal};
  \item \emph{absorción interna} ($-$): cuanto más se gasta dentro, menos queda para pagar;
  \item \emph{acceso a los mercados} ($+$): sin él, el pago depende de la producción corriente.
  \end{itemize}
\item \emph{Umbrales}: no hay uno único para la deuda externa. La regla práctica de la guía anterior (preocupante por encima del 100~\% del PIB y muy preocupante por encima del 200~\%) es orientativa y no sirve para los países del euro o los centros financieros, con deudas brutas muy altas y en su moneda; en los emergentes con historial de impagos, los problemas aparecen con niveles mucho menores \parencite{Reinhart2003Intolerancia}. Mejor la PIIN y la composición.
\end{itemize}

\begin{itemize}
\item \emph{Recapitulación}: capacidad o necesidad (subcuentas, ahorro e inversión, sectores), financiación (instrumentos y reservas) y posición (PIIN, composición y dinámica). Con el diagnóstico de todos los bloques, pasamos a las políticas.
\end{itemize}

% ==========================================================================
\section{Políticas}\label{sec:politicas}

\begin{itemize}
\item La conclusión por adelantado: la política económica busca un crecimiento alto y sostenible y amortiguar los ciclos. Las políticas de \emph{demanda} (monetaria y fiscal) estabilizan el ciclo; las de \emph{oferta} elevan el crecimiento a largo plazo. En cada una hay que juzgar el \emph{tono}, su relación con el ciclo y si ha funcionado.
\end{itemize}

\subsection{Política monetaria}\label{sec:monetaria}

\begin{ideaclave}
\textbf{Qué hay que contestar}: ¿qué objetivo e instrumentos tiene el banco central? ¿Qué tono ha tenido y por qué? ¿Ha funcionado sobre la inflación y la actividad? ¿Qué debería hacer ahora?
\end{ideaclave}

\begin{itemize}
\item \emph{Diseño}: el objetivo es la estabilidad de precios (a menudo un objetivo de inflación numérico, como el 2~\% simétrico a medio plazo del BCE \parencite{BCE2025Estrategia}; en algunos bancos centrales, también el empleo). A largo plazo la política monetaria no altera el nivel de actividad: a corto plazo, sí \vertema{3.A.36}.
\item \emph{Instrumentos}: tipos oficiales; operaciones de financiación a plazo largo; compras de activos (expansión cuantitativa) y su reversión (endurecimiento cuantitativo); orientación sobre los tipos futuros; en los emergentes, intervención cambiaria, coeficientes de reservas o agregados monetarios. Con exceso de liquidez, como en la zona del euro, el tipo que guía la política es el de la facilidad de depósito, y el del mercado a un día cotiza algo por debajo \parencite{BCE2024Marco}.
\item \emph{Tono}: se juzga por el \emph{nivel} del tipo real a corto, $r_t\equiv i_t-\mathbb{E}_t[\pi_{t+1}]$ (en la práctica, $i_t-\pi_t$), frente al tipo real natural $r^n$: \emph{expansiva} (o \emph{acomodaticia}, que es sinónimo) si $r_t<r^n$; neutral si $r_t\approx r^n$; contractiva si $r_t>r^n$. Como $r^n$ no se observa, se completa con la regla de \autor{Taylor} \parencite{Taylor1993},
\[
i_t=r^*+\pi^*+1{,}5\,(\pi_t-\pi^*)+0{,}5\,\mathrm{OG}_t,\qquad r^*=\pi^*=2~\%,
\]
y, si los tipos están en su límite inferior, con el \emph{balance del banco central}.\footnote{En contextos de mucha inflación hay que matizar: en 2022 muchos bancos centrales subieron con fuerza el tipo nominal, pero la inflación era tan alta que el tipo real caía. Una salvedad útil: «aunque los agentes deciden con el tipo real, la fuerte subida del tipo nominal al final de la serie indica un endurecimiento de la política monetaria».}
\item \emph{Canales de transmisión} (2020, 2024) \vertema{3.A.37}:
  \begin{itemize}
  \item \emph{tipo de interés}: del tipo oficial a la curva de tipos y a los tipos de préstamos y depósitos, y de ahí al consumo y la inversión;
  \item \emph{crédito}: la oferta de crédito de los bancos y el valor de las garantías de los prestatarios;
  \item \emph{precios de los activos}: vivienda y bolsa (efecto riqueza, $q$ de \autor{Tobin});
  \item \emph{tipo de cambio}: un tipo más alto aprecia la moneda y reduce la inflación importada;
  \item \emph{expectativas}: si la política es creíble, ancla las de inflación y modera salarios y precios.
  \end{itemize}
\item \emph{Eficacia sobre variables nominales y reales} (2026): si tras las subidas baja la inflación (con 1 o 2 años de retraso), se frena el crédito y se aprecia la moneda, la transmisión funciona. Es menos eficaz con un sistema financiero poco desarrollado, mucha dolarización o dominancia fiscal.
\item \emph{Balance del banco central} (2021, 2026): crece con las compras de activos, los préstamos a los bancos y la compra de divisas. El del Eurosistema pasó de 2 billones de euros en 2014 a 8,8 a mediados de 2022 y bajó a 5,9 en octubre de 2026 \parencite{BCE2026ILM}.
\item \concepto{Dominancia fiscal}: cuando manda la política fiscal y el banco central acaba financiando el déficit (directamente o manteniendo tipos bajos para que la deuda sea pagable), pierde el control de la inflación \parencite{SargentWallace1981,Leeper1991}. Indicios: déficit y deuda altos, crédito del banco central al Estado, tipos reales muy negativos con inflación alta, objetivos incumplidos \vertema{3.A.39}.
\item \emph{La «última milla» de la desinflación} (2024): tras bajar la energía y los bienes, lo difícil es pasar del 3 al 2~\%, porque la inflación de los servicios depende de salarios que se ajustan con retraso. El riesgo es bajar tipos demasiado pronto (repunte) o demasiado tarde (recesión innecesaria) \parencite{BPI2024Informe}.
\item \emph{Relación con el ciclo}: contracíclica si sube los tipos con la brecha positiva y los baja con la negativa; procíclica si hace lo contrario.
\end{itemize}

\begin{figure}[H]
\centering
\caption{Tipos de interés oficiales del BCE y balance del Eurosistema}\label{fig:tipos-balance-bce}
\grafico[0.9\textwidth]{tipos-balance-bce}
\fuente{\textcite{BCE2026TiposOficiales,BCE2026ILM}; hasta octubre de 2026. Qué hay que ver: desde 2022 el tipo que guía la política es el de la facilidad de depósito; subidas de 450 puntos básicos en 2022-2023, bajadas hasta el 2~\% en 2024-2025 y 2 subidas en 2026 (2,50~\% desde septiembre). El balance crece con las compras de activos y la financiación a los bancos y baja desde 2022 al vencer los préstamos y no reinvertirse los valores.}
\end{figure}

\subsection{Política fiscal y deuda pública}\label{sec:fiscal}

\begin{ideaclave}
\textbf{Qué hay que contestar}: ¿qué tono tiene la política fiscal y es adecuado al ciclo? ¿Es sostenible la deuda? ¿Qué riesgos tiene y cómo se reducen? ¿Qué exigen las reglas fiscales?
\end{ideaclave}

\paragraph{Saldos y tono}
\begin{itemize}
\item \emph{Saldos}, en porcentaje del PIB y con la notación del temario \vertema{4.B.4}: el saldo total $b$ (ingresos menos gastos); el primario $b^p$ (sin intereses); el ajustado por el ciclo $b^{ca}=b-\varepsilon\cdot\mathrm{OG}$, con $\varepsilon$ la semielasticidad del saldo a la brecha (en torno a 0,5-0,6 en la UE); el estructural $b^e$, que además quita las medidas puntuales;\footnote{Una medida puntual (\emph{one-off}) es un ingreso o gasto transitorio que no cambia la posición de las cuentas, como una venta de activos o una ayuda puntual a la banca. Ejemplo: en 2012 el déficit de España fue del 10,6~\% del PIB (111.641 millones de euros) \parencite{Eurostat2013EDP}; de ellos, 38.343 millones (unos 3,7 puntos) fueron ayudas a la banca, y sin ellas el déficit era del 6,98~\% \parencite{SEPG2013Deficit}. Esas ayudas no cuentan en el saldo estructural. Las medidas de emergencia de la COVID-19 se trataron, en cambio, como temporales.} y el estructural primario $b^{pe}$.
\item \emph{Tono}: el \concepto{esfuerzo fiscal} es la variación del saldo estructural primario, $\Delta b^{pe}_t$: positivo, contractivo; negativo, expansivo; en torno a 0, \concepto{neutral} (deja actuar a los estabilizadores automáticos, sea cual sea el nivel del saldo). El FMI llama \emph{impulso fiscal} a $-\Delta b^{pe}_t$. Cuidado con el signo: si el saldo era negativo y se hace más negativo, ha \emph{bajado} y el tono es expansivo.
\item \emph{Con déficits muy altos}, el Consejo Fiscal Europeo propuso mirar también el nivel: un déficit estructural primario alto que se reduce algo sigue siendo una política expansiva \parencite{EFB2021FiscalStance}. Desde 2024 la Comisión mide la orientación con el crecimiento del gasto primario neto frente al potencial.
\item \emph{Relación con el ciclo}: contracíclica si el saldo estructural primario sube con la brecha positiva y baja con la negativa (lo ideal); procíclica si al revés. Con histéresis, una política fiscal expansiva en una recesión puede tener efectos duraderos.
\item \emph{Efectos no keynesianos} (2022): con dudas sobre la sostenibilidad, una consolidación creíble puede elevar la confianza y el consumo hoy (un multiplicador negativo).
\end{itemize}

\paragraph{Las reglas fiscales de la UE desde 2024}
\begin{itemize}
\item \emph{Qué se mantiene}: los valores de referencia del Tratado (déficit del 3~\% y deuda del 60~\% del PIB) y el procedimiento de déficit excesivo \parencite{TFUE2012Protocolo12}.
\item \emph{Qué cambia} \parencite{RegUE2024Preventivo,RegUE2024Correctivo} \vertema{3.B.45}:
  \begin{itemize}
  \item un solo indicador operativo, la senda de \concepto{gasto primario neto} (sin intereses, sin el gasto cíclico en desempleo y sin las medidas discrecionales de ingresos);
  \item un plan fiscal-estructural a 4 años, ampliable a 7 con reformas e inversiones, que debe dejar la deuda en una trayectoria descendente plausible;
  \item salvaguardas: la deuda debe bajar al menos 1 punto al año de media si supera el 90~\% del PIB (0,5 entre el 60 y el 90~\%), y el déficit estructural converger a un margen del 1,5~\% del PIB;
  \item una cuenta de control de las desviaciones y cláusulas de escape; en 2025-2026 se activó la nacional para el gasto en defensa (hasta el 1,5~\% del PIB) en 17 Estados.
  \end{itemize}
\item \emph{En el examen}: para un país con déficit y deuda altos, el marco pide una consolidación gradual y plurianual, con más margen si hace reformas e inversiones (se preguntó en 2023) \vertema{4.B.21}.
\end{itemize}

\paragraph{Dinámica y sostenibilidad de la deuda pública}
\begin{itemize}
\item \emph{Cómo se calcula}: con $d_t$ la deuda sobre el PIB, $g_t-\tau_t$ el déficit primario, $i$ el tipo implícito de la deuda y $\gamma^N$ el crecimiento del PIB nominal \vertema{3.A.39},
\begin{equation}\label{eq:dinamica-deuda}
d_t-d_{t-1}=(g_t-\tau_t)+\frac{i_t-\gamma^N_t}{1+\gamma^N_t}\,d_{t-1}+\mathrm{ADD}_t\;\approx\;(g_t-\tau_t)+\frac{r_t-\gamma_{Y,t}}{1+\gamma_{Y,t}}\,d_{t-1}+\mathrm{ADD}_t ,
\end{equation}
donde el segundo término es el \concepto{efecto bola de nieve} y ADD el ajuste déficit-deuda (privatizaciones, ayudas financieras, depósitos del Tesoro, efecto del tipo de cambio en la deuda en divisas).
\item \emph{Saldo primario que estabiliza la deuda}: $\tau-g=\frac{i-\gamma^N}{1+\gamma^N}\,d$. Si $i<\gamma^N$ (o $r<\gamma_Y$), el crecimiento reduce la ratio por sí solo y basta un déficit primario pequeño \parencite{Blanchard2019Deuda}; si $i>\gamma^N$, hacen falta superávits primarios, mayores cuanto más deuda. Ejemplo: con una deuda del 100~\% del PIB, un tipo implícito del 3~\% y un crecimiento nominal del 4~\%, la deuda se estabiliza con un déficit primario de casi 1 punto.
\item \emph{Riesgos} (2026):
  \begin{itemize}
  \item \emph{macroeconómicos}: menos crecimiento, deflación, subidas de tipos, pasivos contingentes (avales, empresas públicas, banca) y envejecimiento \parencite{CE2024Envejecimiento};
  \item \emph{financieros}: deuda en divisas, mucho vencimiento a corto plazo, tenedores no residentes, concentración en la banca nacional (bucle bancario-soberano, \vertema{3.B.44}) y calificación cerca del grado especulativo (por debajo de BBB$-$/Baa3 la deuda deja de ser admisible como garantía en el Eurosistema \parencite{BCE2026ECAF}).
  \end{itemize}
\item \emph{Vida media de la deuda}: cuanto más larga, más despacio se traslada una subida de tipos al coste medio (solo se paga el tipo nuevo por lo que se refinancia) y menor es el riesgo de refinanciación \parencite{Sahuquillo2023}. También importan la deuda a corto, la regional y la indexada.
\item \emph{La inflación y las cuentas públicas}: a corto plazo la inflación mejora el déficit y la ratio de deuda (los ingresos y el PIB nominal suben antes que el gasto indexado y los intereses); a medio plazo el efecto se invierte \parencite{FMI2023FiscalMonitor}.
\item \emph{Ingresos bajos y tamaño del sector público} (2022, 2025): subir ingresos sin dañar el crecimiento (ampliar bases, menos exenciones, lucha contra el fraude y la economía sumergida, impuestos menos distorsionadores) y comparar gasto e ingreso con países de renta parecida; la pregunta no es el tamaño, sino si los ingresos financian el gasto de forma sostenible y si el gasto es eficiente \vertema{4.B.5} \vertema{4.B.7}.
\item \emph{Estrategias para reducir los riesgos}: consolidación gradual y creíble, anunciada pronto (mejor por el gasto corriente y ampliando bases), alargar la vida media y emitir en moneda nacional, diversificar los inversores, reformas que eleven el crecimiento potencial y, en la UE, un plan fiscal-estructural creíble \vertema{4.B.26}.
\end{itemize}

\subsection{Combinación de políticas (\emph{policy mix})}\label{sec:policy-mix}

\begin{itemize}
\item \emph{Qué es}: la combinación de los tonos de la política monetaria y la fiscal (se preguntó en 2020 y 2021).
\item \emph{Cómo se lee}: el \emph{policy mix} es \emph{coherente} si cada política contribuye a los objetivos de la otra, no si las 2 tienen el mismo signo, y \emph{tensionado} si una deshace lo que hace la otra (gráfico~\ref{fig:policy-mix}) \vertema{3.A.40}:
  \begin{itemize}
  \item \emph{ambas expansivas}: adecuado en una recesión con inflación baja (2020); riesgo de inflación si se mantienen en la recuperación;
  \item \emph{ambas contractivas}: adecuado con sobrecalentamiento e inflación; riesgo de recesión;
  \item \emph{fiscal expansiva y monetaria contractiva}: suben los tipos y se aprecia la moneda; poco coherente salvo que la expansión fiscal sea temporal y focalizada (ayudas por la energía en 2022-2023);
  \item \emph{fiscal contractiva y monetaria expansiva}: la combinación clásica para consolidar con deuda alta y brecha negativa sin hundir la demanda.
  \end{itemize}
\item En una unión monetaria la política monetaria es común: el país solo controla la fiscal, que tiene que compensar si su ciclo va por delante o por detrás del de la zona.
\end{itemize}

\begin{figure}[H]
\centering
\caption{Los 4 cuadrantes del \emph{policy mix}}\label{fig:policy-mix}
\grafico[0.6\textwidth]{policy-mix}
\fuente{Elaboración propia. Qué hay que ver: el tono de cada política se sitúa en un eje; la coherencia depende de la situación (inflación, brecha, deuda), no de que las 2 vayan en el mismo sentido.}
\end{figure}

\subsection{Políticas de oferta}\label{sec:politicas-oferta}

\begin{itemize}
\item \emph{Diagnóstico previo}: crecimiento potencial bajo, productividad estancada, inversión baja, paro estructural alto o tasa de actividad baja. Las recomendaciones se ordenan según cuál sea el cuello de botella (2023, 2024, 2026) \vertema{3.A.45} \vertema{3.A.44}:
  \begin{itemize}
  \item \emph{progreso técnico}: I+D (ayudas, deducciones, colaboración entre universidad y empresa) y capital riesgo;
  \item \emph{capital}: incentivos fiscales a la inversión, estabilidad macroeconómica e infraestructuras que vertebren el mercado;
  \item \emph{capital humano}: menos abandono escolar, formación profesional dual, recualificación de los desplazados por la digitalización;
  \item \emph{mercado de trabajo}: según el tipo de paro (apartado~\ref{sec:por-desequilibrio});
  \item \emph{competencia}: liberalización de sectores, menos barreras al comercio interior y exterior, defensa de la competencia, subvenciones transversales (I+D, internacionalización, medio ambiente) en lugar de sectoriales;
  \item \emph{instituciones}: seguridad jurídica, calidad regulatoria, eficacia de la justicia y de la Administración (en 2022 el tribunal dio un gráfico de calidad institucional).
  \end{itemize}
\end{itemize}

\begin{itemize}
\item \emph{Recapitulación}: el tono se mide por niveles reales frente al neutral (monetaria) y por la variación del saldo estructural primario (fiscal); la sostenibilidad de la deuda, por $i-\gamma^N$ y el saldo primario; y la oferta, por el cuello de botella del crecimiento potencial. Con todo ello, solo queda la última pregunta: qué hacer.
\end{itemize}

% ==========================================================================
\section{Recomendaciones de política económica}\label{sec:recomendaciones}

\subsection{Cómo se plantea la pregunta}\label{sec:plantear}

\begin{itemize}
\item La pregunta de desequilibrios y recomendaciones cayó en 12 de los 14 dictámenes y suele ser la última; desde 2023 a veces va dentro de cada pregunta («proponga medidas…») \parencite{ATCEE2026Examenes}.
\item Se contesta en 3 pasos: \emph{qué va bien} (fortalezas), \emph{qué hay que cambiar} (los 2 o 3 desequilibrios principales, ordenados por importancia) y \emph{cómo} (medida, objetivo y riesgo).
\item Nunca tajante ni demasiado preciso: «parece necesaria una consolidación fiscal, pero sin que resulte excesiva para no lastrar la recuperación».
\end{itemize}

\begin{notaopositor}
Esta pregunta se escribe al final, con la matriz de fortalezas y debilidades que se ha ido anotando durante el examen (apartado~\ref{sec:matriz}). Hay recomendaciones que valen para casi todos los países (estabilidad macroeconómica, reformas que eleven la productividad), pero las que puntúan son las que responden a los datos del cuadro.
\end{notaopositor}

\subsection{Recomendaciones por desequilibrio}\label{sec:por-desequilibrio}

\begin{itemize}
\item La tabla~\ref{tab:recomendaciones} da la recomendación tipo para cada desequilibrio; 3 de ellos piden más detalle.
\end{itemize}

\begin{itemize}
\item \emph{Mercado de trabajo deficiente}: primero, diagnosticar qué tipo de paro hay (apartado~\ref{sec:trabajo}) \vertema{3.A.28} \vertema{4.A.16}:
  \begin{itemize}
  \item \emph{políticas de oferta}, contra el paro clásico y estructural: salarios que reflejen la productividad (negociación más cercana a la empresa, indexación a la productividad y no a la inflación pasada, parte del salario ligada a resultados), menos cuña fiscal sobre los salarios bajos, menos dualidad entre contratos fijos y temporales;
  \item \emph{políticas activas de empleo}, contra el desajuste y el paro de larga duración: orientación e intermediación (servicios de empleo eficaces, colaboración con agencias privadas), formación ligada a la demanda de las empresas e incentivos a contratar a quien tiene más dificultades; funcionan mejor bien focalizadas y evaluadas \parencite{Card2018ALMP};
  \item \emph{políticas pasivas} (prestaciones): protegen la renta y sostienen el consumo en las recesiones, pero con cuantía decreciente y ligadas a la búsqueda activa para no desincentivarla;
  \item \emph{paro keynesiano}: políticas de demanda temporales (menos impuestos sobre la renta, más transferencias);
  \item \emph{paro friccional}: mejor emparejamiento y movilidad geográfica (alquiler, planes de estudio adaptados).
  \end{itemize}
\item \emph{Presiones de gasto a medio plazo}: el envejecimiento (pensiones, sanidad y cuidados), la transición energética y digital y, desde 2025, la defensa piden reducir poco a poco las ratios de deuda altas y proteger la inversión pública \parencite{CE2024Envejecimiento}. El FMI, por ejemplo, recomendaba a España en 2026 un ajuste de unos 0,5 puntos del PIB al año entre 2026 y 2030 \parencite{FMI2026ArticuloIV}.
\item \emph{Riesgos de suministro}: ante la fragmentación de las cadenas globales de valor, una \emph{autonomía estratégica abierta}: seguir con la apertura y el multilateralismo, pero diversificar los suministros críticos (energía, metales, componentes), construir reservas estratégicas y elegir socios más próximos.
\end{itemize}

\begin{table}[H]
\centering\small
\caption{Desequilibrios típicos, indicios y recomendaciones}\label{tab:recomendaciones}
\begin{tabular}{>{\raggedright\arraybackslash}p{2.7cm}>{\raggedright\arraybackslash}p{3.8cm}>{\raggedright\arraybackslash}p{7.8cm}}
\toprule
\rowcolor{tceeazulpalido}
\textbf{Desequilibrio} & \textbf{Indicios en el cuadro} & \textbf{Recomendaciones tipo} \\
\midrule
Crecimiento débil o recesión & $\gamma_Y$ bajo o negativo; brecha negativa & Políticas monetaria y fiscal expansivas a corto plazo si hay margen (inflación baja, deuda sostenible), mientras surten efecto las reformas de oferta \\
Crecimiento potencial bajo & $\gamma_{Y^p}$ bajo; PAT estancada; NAIRU alta; déficit corriente persistente & Reformas estructurales: capital humano, I+D, infraestructuras, competencia, reforma del sistema financiero y del mercado de trabajo; plantear el \emph{policy mix} a largo plazo \\
Demanda interna débil & consumo e inversión estancados; ahorro alto & Apoyo a la renta de los hogares con más propensión a consumir; crédito; certidumbre fiscal; si sobra ahorro empresarial, incentivos a invertir o a repartir dividendos \\
Inflación y sobrecalentamiento & brecha positiva; subyacente al alza; salarios y crédito acelerados & Política monetaria contractiva con un banco central independiente; política fiscal que no empuje (ayudas temporales, focalizadas y que mantengan el incentivo a ahorrar energía); competencia; diversificación energética; dejar apreciar la moneda si hay entradas de capital \\
Paro alto y persistente & paro muy por encima de la NAIRU; larga duración; Beveridge desplazada & Según el tipo de paro (lista anterior) \\
Pérdida de competitividad & CLU y tipo de cambio real al alza; cuota de exportación a la baja & Salarios ligados a la productividad; competencia en bienes y factores; I+D; promoción comercial (red de oficinas, agencia de exportación) \\
Déficit exterior y deuda externa & CC muy negativa; PIIN bajo $-35$~\%; deuda a corto & Más ahorro nacional (consolidación fiscal); competitividad; financiarse con inversión directa antes que con deuda; reservas suficientes; flexibilizar un tipo fijo sobrevalorado \\
Déficit y deuda públicos & $b$ muy negativo; $d$ creciente; $i>\gamma^N$ & Consolidación gradual, creíble y anunciada pronto; techo de gasto; corresponsabilidad fiscal; lucha contra el fraude; quitar gasto ineficaz (subvenciones a los combustibles); reformas de pensiones y sanidad; alargar la vida media de la deuda \\
Fragilidad financiera & morosidad alta; capital justo; ROE bajo el coste del capital; crédito excesivo & Saneamiento y recapitalización; medidas macroprudenciales; supervisión; resolución ordenada de entidades inviables; más profundidad del mercado de capitales y educación financiera \\
Dependencia de un sector & mucho peso en el VAB, las exportaciones o los ingresos públicos & Diversificación productiva; fondo de estabilización o regla fiscal estructural; formación para reasignar el empleo \\
Envejecimiento & dependencia creciente; activos a la baja & Más actividad de mujeres y mayores, inmigración, productividad, reforma de pensiones \\
\bottomrule
\end{tabular}
\fuente{Elaboración propia a partir de la guía anterior y de las fichas de \emph{Claves para preparar el dictamen}.}
\end{table}

\subsection{Un programa del FMI}\label{sec:programa-fmi}

\begin{anotaciones}
\emph{Tipos de programa}: acuerdos de derecho de giro (\emph{Stand-By}, problemas de balanza de pagos a corto plazo), Servicio Ampliado (problemas estructurales, plazos más largos), líneas de crédito flexible y de precaución (países con buenos fundamentos, como seguro) y programas concesionarios para los países de renta baja. El desembolso se hace por tramos, si se cumplen criterios cuantitativos (saldo primario, reservas netas, crédito del banco central al Estado) y referencias estructurales \parencite{FMI2025Condicionalidad} \vertema{3.B.22}.

\emph{Efecto típico de un programa de ajuste} (2021 y 2022):
\begin{itemize}
\item \emph{a corto plazo}: menos crecimiento por el ajuste fiscal y la contracción del crédito; sube la ratio de deuda (cae el denominador); mejora la cuenta corriente al caer la absorción; el tipo de cambio se deprecia hacia su nivel de equilibrio y se estabiliza pronto gracias a la financiación del FMI;
\item \emph{a medio plazo} (de 6 meses a 2 años, según la senda de ajuste): vuelve el crecimiento, se estabiliza la inflación en niveles menores, aumentan las reservas, baja el déficit y se estabiliza la deuda.
\end{itemize}
Si nada de eso se ve en el cuadro, el programa no se habría cumplido o habría sido insuficiente. El marco de análisis del FMI (la programación financiera) se resume en el anexo~\ref{anexo:fmi}.
\end{anotaciones}

% ==========================================================================
\section*{Conclusión}\label{sec:cierre}

\begin{itemize}
\item \emph{Cómo se cierra el dictamen}: con una valoración final, en 1 o 2 frases adecuadas y concisas, que resuma la coyuntura del país y responda a la frase de diagnóstico del principio: dónde está la economía, qué la sostiene, qué la amenaza y qué haría falta. Por ejemplo: «En conclusión, la economía afronta en los próximos años estos retos, y su crecimiento dependerá de cómo los aborde; los fondos europeos son, además, una oportunidad para fundamentar una senda de crecimiento sostenible e inclusivo».
\item \emph{Lo que hay que llevarse de esta guía}:
  \begin{itemize}
  \item el método: leer 10-12 minutos, repartir el tiempo con la regla, empezar cada respuesta con la conclusión y anotar fortalezas y debilidades;
  \item el esquema de cada bloque: fases y motor del crecimiento; empleo, CLU y crédito; inflación y tipo de cambio; capacidad de financiación, financiación y PIIN; tono y sostenibilidad de las políticas;
  \item los umbrales, sabiendo cuáles son oficiales y cuáles orientativos, y las recomendaciones por desequilibrio.
  \end{itemize}
\item \emph{Idea final}: «lo bueno, si breve, dos veces bueno». Un buen dictamen no dice todo lo que se sabe: dice, con pocos datos bien elegidos y con prudencia, qué le pasa a esa economía y qué haría falta, en el tiempo que da el tribunal.
\end{itemize}

\clearpage
\appendix

% ==========================================================================
\section{Indicadores y umbrales de referencia}\label{anexo:indicadores}

\begin{itemize}
\item Cada indicador con su cálculo, su referencia y el tema donde se estudia. Las referencias se separan siempre en 2 columnas:
  \begin{itemize}
  \item \emph{oficial}: fijada por una norma o por una institución (Tratado, reglamentos fiscales, Basilea~III, cuadro de indicadores del procedimiento de desequilibrios macroeconómicos de la Comisión, objetivo del BCE), con su fuente;
  \item \emph{orientativa}: regla práctica o referencia supervisora o de mercado, sin consecuencias legales; sirve para ordenar el análisis, nunca para afirmar nada con rotundidad.
  \end{itemize}
\item Los umbrales del procedimiento de desequilibrios (PDM) son de alerta, no obligaciones, y son los del cuadro revisado que se aplica desde el informe del mecanismo de alerta de 2025 \parencite{CE2025AMR2026,Eurostat2025MIP}.
\end{itemize}

\begin{table}[H]
\centering\footnotesize
\caption{Indicadores del perfil cíclico, del mercado de trabajo y de los precios}\label{tab:indicadores-ciclo}
\begin{tabular}{>{\raggedright\arraybackslash}p{2.6cm}>{\raggedright\arraybackslash}p{2.9cm}>{\raggedright\arraybackslash}p{3.3cm}>{\raggedright\arraybackslash}p{3.6cm}>{\raggedright\arraybackslash}p{1.6cm}}
\toprule
\rowcolor{tceeazulpalido}
\textbf{Indicador} & \textbf{Cálculo} & \textbf{Referencia oficial} & \textbf{Referencia orientativa} & \textbf{Tema} \\
\midrule
\rowcolor{tceesalvia}\multicolumn{5}{l}{\emph{Perfil cíclico}} \\
Crecimiento del PIB real & $\gamma_Y$ & — & $<0$ contracción; 0-1~\% estancamiento; 1-2~\% moderado; 2-3~\% por encima del potencial; $\geq3~\%$ sólido (avanzadas) & \vertema{3.A.30} \\
Brecha de producción & $(Y-Y^p)/Y^p\times100$ & 0: producción en su potencial & Su variación en puntos: $\gamma_Y-\gamma_{Y^p}$ & \vertema{4.B.4} \\
Contribución al crecimiento & $\frac{X_{t-1}}{Y_{t-1}}\gamma_{X,t}$ & — & En pp; interna procíclica, externa contracíclica & \vertema{3.A.30} \\
PIB frente a RNB & RNB $=$ PIB $+$ rentas primarias netas & — & RNB $<$ PIB: deudor neto o mucha IED & \vertema{3.B.11} \\
PMI & índice de difusión & 50 & $>50$: la actividad crece frente al mes anterior & \vertema{3.A.42} \\
Sentimiento económico (ESI) & encuestas & 100 (media de largo plazo) & — & \vertema{3.A.42} \\
Pendiente de la curva de tipos & tipo a largo $-$ tipo a corto & — & Negativa: anticipa recesión & \vertema{3.B.24} \\
\rowcolor{tceesalvia}\multicolumn{5}{l}{\emph{Mercado de trabajo}} \\
Tasa de paro & $u=U/A$ & PDM: 10~\% (15-74 años) & Pleno empleo en avanzadas: 3-6~\%; NAWRU de la Comisión & \vertema{3.A.28} \\
Tasa de actividad & $A/P_{16}$ (INE); 15-64 (Eurostat, OCDE) & PDM: caída de más de 0,2 pp en 3 años & Comparar con la misma edad & \vertema{4.A.16} \\
Productividad aparente & $\gamma_Y-\gamma_L$ & — & Puede subir de forma espuria al destruir empleo & \vertema{3.A.45} \\
CLU nominales & $\gamma_W-\gamma_{\mathrm{PAT}}$ & PDM: $+9~\%$ en 3 años (euro), $+12~\%$ (resto), por hora & En el euro, en torno al 2~\% anual mantiene estables los reales & \vertema{4.A.22} \\
CLU reales y margen & $\gamma_{\mathrm{CLU}}-\gamma_P$; $\gamma_P-\gamma_{\mathrm{CLU}}$ & — & CLU reales al alza: menos margen y empleo & \vertema{3.A.28} \\
\rowcolor{tceesalvia}\multicolumn{5}{l}{\emph{Precios}} \\
Inflación (IPC, IAPC) & $\pi_t$ & BCE: 2~\% simétrico a medio plazo & Mirar subyacente, efecto base y expectativas & \vertema{3.A.41} \\
IPC frente a deflactor & $\pi^{\mathrm{IPC}}-\gamma_P$ & — & $>0$: inflación importada; $<0$: mejora de la RRI & \vertema{3.A.30} \\
\bottomrule
\end{tabular}
\fuente{PDM: \textcite{CE2025AMR2026}; BCE: \textcite{BCE2025Estrategia}; NAWRU: \textcite{CE2026AMECO}; el resto, como en los capítulos~\ref{sec:perfil} y~\ref{sec:factores}.}
\end{table}

\begin{table}[H]
\centering\footnotesize
\caption{Indicadores del sector exterior y de las cuentas públicas}\label{tab:indicadores-exterior}
\begin{tabular}{>{\raggedright\arraybackslash}p{2.6cm}>{\raggedright\arraybackslash}p{2.9cm}>{\raggedright\arraybackslash}p{3.3cm}>{\raggedright\arraybackslash}p{3.6cm}>{\raggedright\arraybackslash}p{1.6cm}}
\toprule
\rowcolor{tceeazulpalido}
\textbf{Indicador} & \textbf{Cálculo} & \textbf{Referencia oficial} & \textbf{Referencia orientativa} & \textbf{Tema} \\
\midrule
\rowcolor{tceesalvia}\multicolumn{5}{l}{\emph{Competitividad y sector exterior}} \\
Tipo de cambio efectivo real & $q^{ef}$, variación en 3 años & PDM: $\pm3~\%$ (euro), $\pm10~\%$ (resto) & Ojo: el BCE lo publica al revés (sube = apreciación) & \vertema{3.B.13} \\
Resultados de exportación & frente a las avanzadas, en 3 años & PDM: $-3~\%$ & Cuota que cae sin perder precio: no precio & \vertema{4.A.22} \\
Apertura & $(X+M)/\mathrm{PIB}$ & — & $<15$-20~\% cerrada; $\approx50~\%$ normal; $>75~\%$ muy abierta & \vertema{3.B.9} \\
Saldo por cuenta corriente & $\mathrm{CC}/\mathrm{PIB}$, media de 3 años & PDM: $-4~\%$ y $+6~\%$ & Mirar ahorro, inversión y sectores & \vertema{3.B.11} \\
PIIN & activos $-$ pasivos exteriores & PDM: $-35~\%$ del PIB & Mirar instrumentos, plazos y monedas & \vertema{3.B.12} \\
Deuda externa bruta & pasivos exigibles/PIB & — & 100~\% y 200~\% (regla práctica; no sirve para el euro) & \vertema{3.B.29} \\
Reservas & /importaciones; /deuda a corto; /M2 & — & 3 meses; 100~\% (\autor{Greenspan}-\autor{Guidotti}); 20~\% de M2 & \vertema{3.B.22} \\
\rowcolor{tceesalvia}\multicolumn{5}{l}{\emph{Cuentas públicas}} \\
Saldo público & $b$ & Tratado: $-3~\%$ del PIB & — & \vertema{4.B.1} \\
Esfuerzo fiscal & $\Delta b^{pe}$ & — & $>0$ contractiva; $\approx0$ neutral; $<0$ expansiva & \vertema{4.B.4} \\
Deuda pública & $d$ & Tratado: 60~\%; baja de 1 pp/año si $>90~\%$ (0,5 entre 60 y 90~\%); estructural hacia $-1{,}5~\%$ & — & \vertema{3.A.39} \\
Saldo primario estabilizador & $\frac{i-\gamma^N}{1+\gamma^N}\,d$ & — & Si $i>\gamma^N$, superávit necesario & \vertema{3.A.39} \\
Calificación & escala de las agencias & Grado de inversión: BBB$-$/Baa3 o mejor (CQS~3 del Eurosistema) & — & \vertema{4.B.26} \\
\bottomrule
\end{tabular}
\fuente{PDM: \textcite{CE2025AMR2026}; Tratado y reglas: \textcite{TFUE2012Protocolo12,RegUE2024Preventivo,RegUE2024Correctivo}; reservas: \textcite{FMI2011Reservas}; calificaciones: \textcite{BCE2026ECAF}.}
\end{table}

\begin{table}[H]
\centering\footnotesize
\caption{Indicadores del sistema financiero}\label{tab:indicadores-financiero}
\begin{tabular}{>{\raggedright\arraybackslash}p{2.6cm}>{\raggedright\arraybackslash}p{2.9cm}>{\raggedright\arraybackslash}p{3.3cm}>{\raggedright\arraybackslash}p{3.6cm}>{\raggedright\arraybackslash}p{1.6cm}}
\toprule
\rowcolor{tceeazulpalido}
\textbf{Indicador} & \textbf{Cálculo} & \textbf{Referencia oficial} & \textbf{Referencia orientativa} & \textbf{Tema} \\
\midrule
CET1 / Tier~1 / total & recursos propios/APR & 4,5/6/8~\% + 2,5~\% de conservación $=$ 7/8,5/10,5~\%; requisitos totales del MUS: 11,2~\% CET1 y 15,6~\% total & ABE: CET1 $>14~\%$ holgado & \vertema{3.B.27} \\
Apalancamiento & Tier~1/exposición total & 3~\% & ABE: $>8~\%$ holgado; $<5~\%$ débil & \vertema{3.B.27} \\
LCR y NSFR & véase el apartado~\ref{sec:capitales} & 100~\% & ABE (LCR): $>140~\%$ holgado & \vertema{3.B.27} \\
Dudosos & dudosos/crédito & Dudoso: más de 90 días de impago & ABE: $<3~\%$ bajo; $>8~\%$ alto & \vertema{4.A.19} \\
ROE & beneficio/patrimonio neto & — & Coste del capital: 8-12~\%; ABE: $>10~\%$ favorable & \vertema{3.B.1} \\
Eficiencia & gastos/margen bruto & — & ABE: $<50~\%$ favorable; $>60~\%$ desfavorable & \vertema{4.A.19} \\
Créditos/depósitos & crédito/depósitos & — & ABE: $<100~\%$; $>150~\%$ desfavorable & \vertema{4.A.19} \\
Brecha crédito-PIB & crédito/PIB $-$ tendencia & — & $>2$ pp: activar el colchón anticíclico & \vertema{3.A.42} \\
Deuda de hogares y de empresas & consolidada/PIB & PDM: 55~\% y 85~\% & — & \vertema{4.A.18} \\
Flujo de crédito & \% de la deuda en $t-1$ & PDM: 14~\% (hogares), 13~\% (empresas) & — & \vertema{4.A.18} \\
Precio de la vivienda & variación nominal anual & PDM: $+9~\%$ & Sobrecarga: gasto en vivienda $>40~\%$ de la renta & \vertema{4.A.11} \\
\bottomrule
\end{tabular}
\fuente{PDM: \textcite{CE2025AMR2026}; Basilea: \textcite{BCBS2017Basel3,RegUE2013CRR,BCE2025SREP}; referencias de la ABE: \textcite{ABE2026RiskDashboard,ABE2025RAR}; brecha crédito-PIB: \textcite{BCBS2010Anticiclico}; vivienda: \textcite{Eurostat2025Sobrecarga}.}
\end{table}

\begin{table}[H]
\centering\footnotesize
\caption{Requisitos de capital y de liquidez de Basilea~III (\% de los activos ponderados por riesgo)}\label{tab:basilea}
\begin{tabular}{>{\raggedright\arraybackslash}p{6.2cm}>{\centering\arraybackslash}p{1.6cm}>{\centering\arraybackslash}p{1.6cm}>{\centering\arraybackslash}p{1.8cm}}
\toprule
\rowcolor{tceeazulpalido}
 & \textbf{CET1} & \textbf{Tier 1} & \textbf{Capital total} \\
\midrule
Mínimo (Pilar 1) & 4,5 & 6,0 & 8,0 \\
+ Colchón de conservación & 2,5 & 2,5 & 2,5 \\
\rowcolor{tceesalvia}
Pilar 1 + conservación & 7,0 & 8,5 & 10,5 \\
+ Colchón anticíclico (España: 1~\% desde octubre de 2026) & \multicolumn{3}{c}{0-2,5 (en CET1)} \\
+ Colchones sistémicos (mundiales, 1-3,5; otros, hasta 3) y de riesgo sistémico & \multicolumn{3}{c}{según la entidad y el país} \\
+ Pilar 2 (P2R) y recomendación (P2G) & \multicolumn{3}{c}{según la entidad (revisión supervisora)} \\
\rowcolor{tceesalvia}
Requisitos y recomendación totales medios del MUS en 2026 & 11,2 & — & 15,6 \\
Ratio de apalancamiento (Tier 1 sobre la exposición total) & \multicolumn{3}{c}{3,0} \\
LCR y NSFR & \multicolumn{3}{c}{100} \\
\bottomrule
\end{tabular}
\fuente{\textcite{BCBS2017Basel3}; en la UE, \textcite{RegUE2013CRR,DirUE2013CRD,RegUE2024CRR3}; colchón anticíclico de España: \textcite{BdE2025CCA}; requisitos del MUS: \textcite{BCE2025SREP}.}
\end{table}

\begin{table}[H]
\centering\footnotesize
\caption{Ratios de la banca que opera en España y de la europea (segundo trimestre de 2026, \%)}\label{tab:ratios-banca}
\begin{tabular}{>{\raggedright\arraybackslash}p{4.2cm}>{\centering\arraybackslash}p{2.2cm}>{\centering\arraybackslash}p{2.4cm}>{\centering\arraybackslash}p{2.4cm}}
\toprule
\rowcolor{tceeazulpalido}
\textbf{Ratio} & \textbf{España} & \textbf{Entidades significativas del MUS} & \textbf{UE (ABE)} \\
\midrule
ROE & 16,05 & 10,73 & 11,3 \\
CET1 & 14,11 & 16,00 & 16,1 \\
Tier 1 & 15,68 & 17,53 & 17,7 \\
Capital total & 18,30 & 20,13 & 20,2 \\
Préstamos dudosos & 2,53 & 2,17 & 1,8 \\
Apalancamiento & 5,60 & — & 5,8 \\
LCR & 171,07 & 154,91 & 158,5 \\
NSFR & — & — & 125,7 \\
Créditos/depósitos & 97,17 & — & 106,3 \\
Eficiencia & — & 53,06 & 51,5 \\
\bottomrule
\end{tabular}
\fuente{\textcite{BdE2026EstadSupervisoras,BCE2026SupervisionSI,ABE2026RiskDashboard}. Qué hay que ver: la banca española es más rentable que la europea, pero tiene menos capital de máxima calidad; hay que mirar las 2 cosas a la vez. La ratio de dudosos del Banco de España no es comparable con la de series antiguas.}
\end{table}

\begin{table}[H]
\centering\footnotesize
\caption{Escala armonizada de calificaciones crediticias del Eurosistema (largo plazo)}\label{tab:calificaciones}
\begin{tabular}{>{\centering\arraybackslash}p{1cm}>{\centering\arraybackslash}p{2.4cm}>{\centering\arraybackslash}p{2.4cm}>{\centering\arraybackslash}p{2.4cm}>{\centering\arraybackslash}p{2.4cm}>{\centering\arraybackslash}p{2.4cm}}
\toprule
\rowcolor{tceeazulpalido}
\textbf{Nivel} & \textbf{DBRS} & \textbf{Fitch} & \textbf{Moody's} & \textbf{S\&P} & \textbf{Scope} \\
\midrule
1 & AAA a AAL & AAA a AA$-$ & Aaa a Aa3 & AAA a AA$-$ & AAA a AA$-$ \\
2 & AH, A, AL & A+, A, A$-$ & A1, A2, A3 & A+, A, A$-$ & A+, A, A$-$ \\
\rowcolor{tceesalvia}
3 & BBBH, BBB, BBBL & BBB+, BBB, BBB$-$ & Baa1, Baa2, Baa3 & BBB+, BBB, BBB$-$ & BBB+, BBB, BBB$-$ \\
4 & BBH & BB+ & Ba1 & BB+ & BB+ \\
5 & BB & BB & Ba2 & BB & BB \\
\bottomrule
\end{tabular}
\fuente{\textcite{BCE2026ECAF}. Qué hay que ver: el grado de inversión llega hasta el nivel 3 (BBB$-$/Baa3); por debajo, la deuda deja de ser admisible como garantía en el Eurosistema.}
\end{table}

\begin{table}[H]
\centering\footnotesize
\caption{Rasgos de una economía avanzada y de una emergente o en desarrollo}\label{tab:avanzada-emergente}
\begin{tabular}{>{\raggedright\arraybackslash}p{2.9cm}>{\raggedright\arraybackslash}p{5.4cm}>{\raggedright\arraybackslash}p{5.4cm}}
\toprule
\rowcolor{tceeazulpalido}
 & \textbf{Economía avanzada} & \textbf{Economía emergente o en desarrollo} \\
\midrule
Crecimiento & potencial del 1-2~\%; menos volátil & potencial mayor (4-6~\%); más volátil \\
Inflación & cerca del objetivo (2~\%) & más alta y volátil; mayor traspaso del tipo de cambio \\
Estructura productiva & servicios en torno a 3 cuartas partes del VAB; primario pequeño & más peso de la agricultura y de las materias primas; dependencia de un sector \\
Política fiscal & contracíclica o acíclica; deuda en moneda propia & a menudo procíclica; ingresos bajos (informalidad); deuda en divisas \\
Política monetaria & tipos de interés y objetivo de inflación creíble & a veces agregados o tipo de cambio; posible dominancia fiscal \\
Sector exterior & déficits financiables con inversión directa y de cartera & riesgo de frenazo súbito de capitales; remesas; reservas como seguro \\
Sistema financiero & profundo, con mercados de capitales & bancarizado y poco profundo; préstamos en divisas \\
Demografía & envejecimiento & transición demográfica inacabada \\
\bottomrule
\end{tabular}
\fuente{Elaboración propia a partir de \textcite{Kaminsky2004Rains}, \textcite{Calvo1998SuddenStops}, \textcite{CalvoReinhart2002} y \textcite{EHP2005PecadoOriginal}.}
\end{table}

% ==========================================================================
\section{Qué se ha preguntado (2011-2026)}\label{anexo:preguntas}

\begin{itemize}
\item De 2011 a 2026 ha habido 14 dictámenes (la numeración es por oferta de empleo público: el de la oferta de 2015 se hizo en enero de 2016, etc.). El sector exterior y la política fiscal han caído en todos; el perfil cíclico, en 13 (no en 2025); las recomendaciones, en 12; la política monetaria y el mercado de trabajo, en 11 \parencite{ATCEE2026Examenes}.
\item Desde 2023 hay menos preguntas (5 o 6), más largas, sin cálculos y sin puntuación, y se pide más el \emph{porqué}: determinantes de la inversión y de la inflación, diferencias entre el PIB y la RNB, traspaso del tipo de cambio, «última milla» de la desinflación, nuevas reglas fiscales, demografía y aranceles.
\item A continuación, las preguntas literales de cada examen (sin los cuadros ni los gráficos).
\end{itemize}

\begin{table}[H]
\centering\footnotesize
\caption{Bloques preguntados en cada dictamen}\label{tab:bloques-anio}
\begin{tabular}{>{\raggedright\arraybackslash}p{2.9cm}*{14}{>{\centering\arraybackslash}p{0.48cm}}}
\toprule
\rowcolor{tceeazulpalido}
\textbf{Bloque} & \rotatebox{90}{2011} & \rotatebox{90}{2013} & \rotatebox{90}{2014} & \rotatebox{90}{2016 ene.} & \rotatebox{90}{2016 nov.} & \rotatebox{90}{2018} & \rotatebox{90}{2019} & \rotatebox{90}{2020} & \rotatebox{90}{2021} & \rotatebox{90}{2022} & \rotatebox{90}{2023} & \rotatebox{90}{2024} & \rotatebox{90}{2025} & \rotatebox{90}{2026} \\
\midrule
Perfil cíclico & • & • & • & • & • & • & • & • & • & • & • & • &  & • \\
Mercado de trabajo & • &  & • & • & • & • & • & • & • &  & • & • &  & • \\
Sistema financiero &  &  &  &  &  &  &  & • &  & • & • &  &  &  \\
Inflación &  &  &  &  & • &  &  &  &  &  & • &  & • & • \\
Sector exterior & • & • & • & • & • & • & • & • & • & • & • & • & • & • \\
Política monetaria & • &  &  & • & • & • &  & • & • & • & • & • & • & • \\
Política fiscal & • & • & • & • & • & • & • & • & • & • & • & • & • & • \\
Recomendaciones &  & • & • & • & • & • & • & • & • & • & • & • &  & • \\
\bottomrule
\end{tabular}
\fuente{Elaboración propia a partir de los enunciados oficiales \parencite{ATCEE2026Examenes}. Una pregunta puede tocar varios bloques.}
\end{table}

\begin{small}

\paragraph{2026}
\textit{Examen del 7 de marzo de 2026.} El tribunal valorará fundamentalmente las capacidades analítica y de síntesis del aspirante, así como su claridad expositiva»; sin calculadora; abreviaturas solo las del enunciado o las convencionales. «En los cuadros y gráficos a continuación de este enunciado se facilita información económica sobre un determinado país referida a los años naturales 2014 a 2025. Atendiendo a esa información, elabore un escrito respondiendo a las siguientes cuestiones:
\begin{enumerate}[leftmargin=*,itemsep=1pt]
\item Describa brevemente el perfil cíclico de la economía desde la óptica de la demanda a lo largo del periodo 2014-2024. ¿Cuáles podrían ser los determinantes del comportamiento de la inversión? ¿A qué podrían deberse las diferencias entre el producto interior bruto (PIB) y la renta nacional bruta (RNB)?
\item Analice la evolución del sector exterior a lo largo de todo el horizonte temporal, identificando sus fortalezas y debilidades. A tenor de los datos disponibles sobre precios internos, ¿puede afirmarse que el tipo de cambio presenta un alto grado de traspaso (pass-through)? ¿Cree que la política arancelaria de la actual administración de Estados Unidos ha tenido algún impacto sobre el país objeto de análisis?
\item Caracterice el mercado laboral del país. ¿Qué restricciones pueden estar imponiendo las dinámicas demográficas sobre este mercado? Teniendo en cuenta estas restricciones, proponga medidas de política económica orientadas a sostener el crecimiento por el lado de la oferta.
\item Describa la situación de las cuentas públicas y los principales rasgos de la política fiscal. ¿Considera que existen riesgos macroeconómicos y financieros apreciables para la sostenibilidad de la deuda pública? En caso positivo, indique posibles estrategias para mitigar tales riesgos.
\item Analice la capacidad aparente de la política monetaria para influir sobre las variables nominales y reales. ¿Qué factores podrían estar detrás tanto de la resistencia al alza inicial de la inflación como de su posterior cambio de tendencia? Valore el comportamiento de la autoridad monetaria en el nuevo contexto inflacionario.
\end{enumerate}

\paragraph{2025}
\textit{Examen del 8 de marzo de 2025.} El Tribunal valorará la capacidad de síntesis y la capacidad de repartir adecuadamente el tiempo proporcionado entre las distintas preguntas»; 1 hora y 45 minutos. «En los cuadros y gráficos adjuntos se facilita información macroeconómica sobre una determinada economía para el período 2017-2026. Atendiendo a dicha información, responda a cada una de las siguientes preguntas:
\begin{enumerate}[leftmargin=*,itemsep=1pt]
\item Analice la importancia del sector petrolero en esta economía, tanto a efectos de su contribución a la variación del PIB y de los precios como su relevancia para el sector exterior y las cuentas públicas.
\item A la luz de los datos facilitados, ¿cuáles son, en su opinión, los factores fundamentales detrás del repunte de la inflación?
\item Analice la evolución de los principales rubros del sector exterior, fortalezas y debilidades.
\item Analice los principales elementos de las cuentas públicas, tanto en términos de su evolución como del tamaño del Sector Público. ¿Cuáles son las claves de la gestión del déficit y de la deuda pública del país para garantizar su sostenibilidad a largo plazo?
\item Analice cómo ha evolucionado la orientación de la política monetaria desde 2020 e indique cuáles son los objetivos e instrumentos principales del Banco Central del país. ¿Cuál cree que debería ser la orientación de la política monetaria en los próximos años?
\end{enumerate}

\paragraph{2024}
\textit{Examen del 5 de junio de 2024.} En el cuadro y gráfico adjuntos se facilita información macroeconómica del país (A) referida al período 2019 – 2026. Dicha economía cuenta con una moneda propia que fluctúa libremente. Conteste a las siguientes preguntas apoyándose únicamente en la información que le ha sido suministrada.
\begin{enumerate}[leftmargin=*,itemsep=1pt]
\item Analice brevemente el perfil cíclico de esta economía. Comente brevemente el comportamiento de los principales componentes del lado de la demanda, así como de la oferta. En función de los indicadores de alta frecuencia de los que dispone, ¿cuál es la situación coyuntural reciente de esta economía?
\item Describa brevemente la evolución y las características del mercado laboral de esta economía. Explique sus principales implicaciones tanto a corto como a largo plazo.
\item Analice la Balanza de Pagos de la presente economía. ¿Qué factores pueden explicar la evolución reciente de la Posición de Inversión Internacional Neta? ¿Qué opinión le merece la evolución de las reservas en esta economía?
\item Describa brevemente cuáles han podido ser los posibles canales de transmisión y efectos sobre los principales agregados macroeconómicos de la política monetaria adoptada durante el periodo 2019-2023. ¿A qué retos se enfrenta la autoridad monetaria para alcanzar la «última milla» del proceso de disminución de la inflación?
\item Evalúe el tono de la política fiscal. ¿Qué recomendación daría a medio y largo plazo en relación con la sostenibilidad y la eficiencia del sistema de finanzas públicas de la economía en cuestión?
\item Detalle los principales impedimentos de esta economía para lograr un crecimiento sostenido. ¿Qué reformas estructurales considera que son clave para mejorar el desempeño económico a largo plazo?
\end{enumerate}

\paragraph{2023}
\textit{Examen del 14 de junio de 2023.} PM para designar política monetaria deberá leerse literalmente»). «En el cuadro y gráfico adjuntos se facilita información macroeconómica referida al periodo 2012-2024 sobre un determinado Estado Miembro de la Eurozona. Atendiendo a dicha información, responda a las siguientes preguntas:
\begin{enumerate}[leftmargin=*,itemsep=1pt]
\item ¿Cuáles podrían ser las explicaciones del desfase entre el Producto Interior Bruto y la Renta Nacional Bruta en el año 2015? Analice el perfil cíclico del consumo privado entre 2019 y 2022 así como sus determinantes. Analice, asimismo, el comportamiento de la inversión privada entre 2019 y 2021, así como sus determinantes.
\item Analice el comportamiento de la inflación desde 2021 y sus principales determinantes. Viendo las previsiones de 2023, ¿podrían estar afectadas éstas por el efecto base? ¿Cree que el brote inflacionista de 2021 es temporal? A la luz de los datos y de su respuesta evalúe la efectividad de la política monetaria a partir de 2022. Valore brevemente la situación del sistema financiero.
\item ¿A qué puede deberse la volatilidad en la Capacidad/Necesidad de Financiación de esta economía entre 2018 y 2022?
\item Defina la tasa de actividad. ¿A qué podría deberse el aumento de la tasa de actividad entre 2020 y 2022? Analice la evolución de la Curva de \autor{Beveridge} a lo largo de la serie.
\item Valore el tono y la orientación de la política fiscal a lo largo de la serie, así como la evolución de la deuda pública. ¿Cómo considera que la entrada en vigor de las nuevas reglas fiscales en 2024 afectará a la política fiscal del país?
\item Dada la situación del país ¿qué políticas recomendaría para estimular el crecimiento en el medio y largo plazo?
\end{enumerate}

\paragraph{2022}
\textit{Examen del 9 de julio de 2022.} En el cuadro y gráficos adjuntos se facilita información macroeconómica, referida al periodo 2015-2021, sobre un determinado país. Atendiendo a dicha información, responda a las siguientes preguntas:
\begin{enumerate}[leftmargin=*,itemsep=1pt]
\item Analice el perfil cíclico de esta economía a lo largo del periodo 2015-2021, así como la composición del crecimiento desde la óptica de la demanda. (1,5 puntos)
\item Realice un análisis en detalle del sector exterior. Entre otros, explique las causas del déficit permanente en la balanza de pagos de esta economía. ¿Se cumple la condición de \autor{Marshall}-\autor{Lerner}? Analice posibles alternativas de política económica en este ámbito. (2,5 puntos)
\item Evalúe el comportamiento del sector público y analice los principales rasgos de la política fiscal. ¿Aprecia usted efectos no keynesianos de la política fiscal? Comente posibles alternativas de política económica en relación con la baja participación de los ingresos públicos en el PIB del país. (1 punto)
\item ¿Qué tipo de política monetaria cree que están aplicando las autoridades? ¿Deberían continuar con esta política en 2022? (1 punto)
\item Realice una valoración de la estabilidad financiera de la economía objeto de análisis, identificando las principales vulnerabilidades y fortalezas del sector bancario. (1 punto)
\item Este país ha recibido financiación del Fondo Monetario Internacional (FMI) en el marco de un programa de ayuda condicionado. ¿Qué impacto ha tenido el programa a la luz de la evolución del saldo público y del sector exterior? (1,5 puntos)
\item Señale cuáles son las principales vulnerabilidades y riesgos de esta economía y comente qué políticas deberían aplicarse en los próximos años para abordarlos. (1,5 puntos)
\end{enumerate}

\paragraph{2021}
\textit{Examen del 10 de julio de 2021.} Los datos, cuadros y gráficos aportados contienen información de un país integrante del G-20. Actualmente el régimen cambiario de esta economía es el de flotación, guardándose la posibilidad de intervenir en los mercados cambiarios (este régimen también se conoce frecuentemente como «flotación sucia»). Sin embargo, en el pasado, esta economía ha mantenido un régimen más cercano al tipo de cambio fijo. Conteste a las siguientes preguntas apoyándose únicamente en la información que le ha sido suministrada.
\begin{enumerate}[leftmargin=*,itemsep=1pt]
\item Analice el perfil cíclico de esta economía a lo largo del período. Comente, BREVEMENTE, el comportamiento de los principales componentes del lado de la demanda. Observe, si es posible, y comente la evolución del PIB potencial. (1,5 puntos)
\item Con los datos aportados, calcule la variación de los costes laborales unitarios reales para 2017, 2018 y 2019. Para este periodo, compare los costes laborales unitarios con la posición cíclica de la economía. ¿Qué factores podrían explicar su evolución? Indique si consideraría oportuno o no utilizar la Ley de \autor{Okun} y la Curva de \autor{Phillips} a efectos de modelizar la economía de este país. (1 punto)
\item Este país sufrió una grave crisis económica en 2018 desencadenada en el sector exterior. Comente con detalle los desequilibrios que se fueron acumulando en el sector exterior y cómo estos pudieron contribuir a que se originara la crisis. (2 puntos)
\item Evalué brevemente el tono de la política monetaria. ¿Considera usted que hay dominancia fiscal? Observando los datos aportados sobre su balance, comente los posibles objetivos que persigue el banco central. El banco central de este país ha anunciado su intención de incrementar sus reservas en moneda extranjera. Explique las posibles consecuencias de esta medida sobre la inflación. (2 puntos)
\item Atendiendo a la política fiscal, ¿considera que el policy mix de esta economía es el adecuado? ¿Cuáles han sido las consecuencias de la política de endeudamiento del sector público? De acuerdo con la posición cíclica de la economía, ¿el saldo primario estructural de 2011 debería ser mayor, menor o igual al 1,6\% del PIB? Razone brevemente la respuesta. (2 puntos)
\item Enumere cuáles son, en su opinión, los desequilibrios económicos más importantes a los que se enfrenta este país. Formule recomendaciones de política económica para corregirlos. ¿Qué implicaciones tendría el acceso a un programa de financiación del Fondo Monetario Internacional para la corrección de estos equilibrios? (1,5 puntos)
\end{enumerate}

\paragraph{2020}
\textit{Examen del 18 de enero de 2020.} Los cuadros y gráficos contienen datos de una economía desarrollada, con una alta apertura al exterior y un peso elevado de su sector agroalimentario. El objetivo prioritario de la política monetaria es mantener la tasa de inflación medida mediante un índice de precios al consumo en una banda de fluctuación en torno al 2\% interanual. El régimen cambiario al que se han comprometido las autoridades es el de flotación libre y ha sido respetado salvo durante la última gran crisis de 2008 y 2009 cuando intervinieron de manera puntual. Entre sus principales mercados de exportación se encuentran la Unión Europea, los Estados Unidos y China, entre otros.
\begin{enumerate}[leftmargin=*,itemsep=1pt]
\item Analice el perfil cíclico de esta economía a lo largo del período relacionando la evolución de las principales macromagnitudes con sus determinantes más significativos. ¿Cuál ha sido la contribución al crecimiento de la demanda interna en los ejercicios 2017 y 2018?
\item Realice un análisis detallado del sector exterior de esta economía identificando sus principales fortalezas y debilidades. Entre otros, evalúe cómo ha evolucionado su competitividad precio así como su patrón de financiación externa. Compare y analice posibles alternativas de política económica en este ámbito. Apoyándose en su conocimiento de la actual coyuntura económica mundial, ¿cuáles son, en su opinión, los riesgos a los que se enfrenta su sector exterior?
\item ¿Cómo valoraría usted la política monetaria aplicada? ¿Cuáles son, en su opinión, sus principales canales de transmisión? Dada la evolución mostrada por el tipo de cambio de la moneda nacional, ¿consideraría usted conveniente que las autoridades abandonasen, siquiera puntualmente, la flotación libre a fin de mejorar el saldo exterior?
\item Realice un análisis detallado del mercado de trabajo de esta economía, evaluando sus principales vulnerabilidades y fortalezas. Compare y analice posibles alternativas de política económica en este ámbito. En 2017 el Gobierno decidió endurecer los requisitos para la concesión de visados a los trabajadores extranjeros. Describa brevemente los posibles mecanismos de transmisión al producto nacional de este shock exógeno.
\item Realice una valoración de la estabilidad financiera de la economía objeto de análisis desde el punto de vista macroprudencial, evaluando los posibles riesgos, tanto internos como externos, a los que se enfrenta. Teniendo en cuenta que su sistema financiero está muy «bancarizado», incluya una valoración de la solvencia y rentabilidad de su sistema bancario. ¿En qué medida los bajos tipos de interés pueden ser considerados un riesgo a la estabilidad financiera? Compare posibles alternativas de política económica a este respecto.
\item Realice una valoración de la participación del sector público analizando tanto el tono de la política fiscal como los riesgos a los que ha de hacer frente. Compare posibles alternativas de política económica en este ámbito indicando sus potenciales interrelaciones con la política monetaria En 2017 se celebraron sus últimas elecciones generales y se produjo un cambio de Gobierno. El nuevo Gobierno estableció como objetivos de política económica reducir el fuerte déficit de infraestructuras del país y solventar la escasez de vivienda para lo que puso en marcha una política fiscal expansiva de inversión en infraestructuras y de construcción de vivienda social. ¿Qué valoración le merece esta medida?
\item Para concluir, tras el análisis que usted acaba de realizar, ¿cuáles son en su opinión los principales desequilibrios macroeconómicos que debía abordar la política económica del nuevo Gobierno que se constituyó en 2017?
\end{enumerate}

\paragraph{2019}
\textit{Examen del 2 de febrero de 2019.} En los siguientes cuadros y gráficos se muestran las principales variables macroeconómicas de un país de la Unión Europea.
\begin{enumerate}[leftmargin=*,itemsep=1pt]
\item Comente brevemente el perfil cíclico de esta economía desde la óptica de la demanda a la largo del período 2012-2017 ¿Cuáles han sido los principales motores de la recuperación económica?
\item ¿Cuáles son los principales determinantes de la evolución de las exportaciones de bienes y servicios? ¿Cómo han influido dichos determinantes en las exportaciones entre 2015 y 2017?
\item ¿Cuáles eran los principales desequilibrios de esta economía antes del inicio de la crisis financiera internacional? ¿Se han corregido algunos de ellos después de una década? ¿Han aparecido otros nuevos?
\item A la luz de la evolución de los indicadores de clima económico en 2018, ¿cómo cree que será el desempeño de esta economía en 2018 y 2019?
\item Analice la evolución de la productividad aparente del factor trabajo durante el período 2012-2017. Calcule la variación de los costes laborales unitarios en 2016 y 2017.
\item Comente las principales transformaciones del mercado de trabajo y la situación del mismo en 2017.
\item El gobierno ha aprobado una subida del salario mínimo interprofesional del 4,1\% a partir del 1 de enero de 2019. Valore los efectos que esta medida podría tener en las principales variables macroeconómicas, incluyendo el saldo público.
\item Calcule la tasa de inversión nacional en 2017. ¿Cuál puede ser el motivo de que, a pesar de la política monetaria expansiva aplicada por el banco central en los últimos años, la tasa de inversión de esta economía sea relativamente baja respecto a la media de la Unión Europea?
\item Analice la política fiscal entre los años 2012 y 2017. ¿Cuál ha sido el «esfuerzo estructural» en los últimos tres años?
\item En un reciente Informe de Perspectivas Económicas Mundiales del Fondo Monetario Internacional, se considera que un agravamiento de las tensiones comerciales a nivel global o una subida abrupta de los tipos de interés, en un entorno de normalización de las políticas monetarias de Estados Unidos, constituyen uno de los riesgos que podrían lastrar el crecimiento económico. ¿Esta economía es vulnerable a alguno de estos riegos? ¿Por qué?
\end{enumerate}

\paragraph{2018}
\textit{Examen del 20 de enero de 2018.} Dictamen económico. En los cuadros y gráficos adjuntos, se facilita información macroeconómica, referida al periodo 2004-2016 (con proyecciones para 2017-2018), sobre un determinado país. Atendiendo a dicha información, responda a las siguientes preguntas:
\begin{enumerate}[leftmargin=*,itemsep=1pt]
\item Evalúe el perfil cíclico de la economía de este país entre 2004 y 2016 valorando la aportación de los diferentes componentes al crecimiento. ¿Cómo ha cambiado la composición de los saldos sectoriales antes y después de 2009? (2 puntos)
\item Calcule, de forma aproximada, las variaciones de la productividad aparente del factor trabajo y de los costes laborales unitarios (CLU) para 2009 y 2010. Analice la serie completa de CLU y la de productividad en relación con la evolución del PIB real. A partir de los datos del mercado de trabajo, ¿detecta alguna posible vulnerabilidad para el crecimiento a medio/largo plazo? (1,5 puntos)
\item ¿Cómo ha evolucionado el sector exterior antes y después de 2009? ¿Qué factores explican la evolución de dicho saldo? Indique si se ha producido una ganancia o pérdida de competitividad después de 2012 y explique qué efecto ha tenido en este sentido la relación de precios con el resto del mundo. (1,5 puntos)
\item A partir de la información de balanza de pagos, ¿cómo se ha financiado el déficit por cuenta corriente desde 2009? ¿Por qué mejora la posición de inversión internacional neta (PIIN) entre 2013 y 2015 si el saldo por cuenta financiera es negativo en esos años? (1 punto)
\item Comente cómo ha evolucionado el saldo público y cuál ha sido el tono de la política fiscal. Calcule el esfuerzo fiscal en el periodo 2015-2016. (1,5 puntos)
\item Analice cuál ha sido la orientación de la política monetaria. (1 punto)
\item A partir de las proyecciones del crecimiento real y potencial, calcule la brecha del producto (output gap) para los años 2017 y 2018. Realice un análisis de la situación económica en 2016 y del crecimiento estimado en 2017 y 2018. Indique cuáles son las principales vulnerabilidades y riesgos de esta economía y qué políticas deberían aplicarse en los próximos años. (1,5 puntos)
\end{enumerate}

\paragraph{2016 (noviembre)}
\textit{Examen del 26 de noviembre de 2016.} Análisis de coyuntura. En el cuadro y gráficos adjuntos se facilita información macroeconómica, referida al periodo 2007-2016, sobre un determinado país. Atendiendo a dicha información, responda a las siguientes preguntas:
\begin{enumerate}[leftmargin=*,itemsep=1pt]
\item Evalúe el perfil cíclico de la economía de este país desde el inicio de la crisis financiera internacional hasta 2015.
\item ¿Cómo ha podido afectar la evolución del precio del petróleo a esta economía desde 2008?
\item Calcule la variación del coste laboral unitario de 2012 respecto al año anterior. Con la información disponible, analice la evolución de los costes laborales unitarios y su relación con la evolución de los salarios, suponiendo que la tasa de asalarización es constante e igual a 1 (2012-2015). ¿Qué puede deducirse en relación con la competitividad de esta economía?
\item ¿Cuáles son los principales desequilibrios de esta economía?
\item Comente la evolución de los principales determinantes de la demanda de exportaciones de bienes y servicios y su contribución al crecimiento de las exportaciones. Enumere otros posibles determinantes de la demanda de exportaciones. ¿Cree que la balanza de servicios es superavitaria?
\item Esta economía ha sufrido un shock en 2016, con efectos sobre el tipo de cambio y los índices de gestores de compras de los sectores de manufacturas y servicios (PMI, siglas en inglés). A la vista de estos gráficos, ¿cuál podría ser el impacto (positivo/negativo) de dicho shock en el crecimiento y en los principales componentes del PIB desde la óptica de la demanda, así como en el empleo y la inflación?
\item Analice el tono de la política económica (política fiscal y política monetaria) en el período 2008-2015.
\item ¿Cuál podría ser el dilema de política monetaria del Banco Central en 2016?
\end{enumerate}

\paragraph{2016 (enero)}
\textit{Examen del 16 de enero de 2016.} Análisis de coyuntura. En el cuadro adjunto se facilita información macroeconómica, referida al periodo 2010-2016, sobre un determinado país importador neto de petróleo. En base a dicha información, responda a las siguientes cuestiones en el orden establecido:
\begin{enumerate}[leftmargin=*,itemsep=1pt]
\item Evalúe el perfil cíclico de la economía real en el período 2010-2015 y la composición del crecimiento desde la óptica de la demanda.
\item ¿Qué razones podrían explicar el repunte de los precios de la vivienda a partir de 2013?
\item Calcule las tasas de variación anuales de la productividad aparente del trabajo y de los costes laborales unitarios nominales de la economía en el período 2013-2015. Analice su evolución.
\item ¿Qué shocks habrían afectado al saldo de la cuenta corriente en 2014 y 2015? Explique su respuesta.
\item Defina el efecto valoración de la Posición Neta de Inversión Internacional (PNII) y calcule su valor en 2015 y 2016. ¿A qué atribuye el empeoramiento de la PNII en estos años?
\item Calcule el esfuerzo estructural fiscal en el período 2015-2016 bajo el supuesto de que no se han adoptado medidas one-off. Analice el tono de la política fiscal en dichos años.
\item ¿Qué tipo de política monetaria cree que están aplicando las autoridades? ¿Deberían continuar con esta política en 2016? Explique la respuesta.
\item Señale los principales riesgos que presenta actualmente esta economía.
\end{enumerate}

\paragraph{2014}
\textit{Examen del 27 de septiembre de 2014.} Primer Ejercicio. Segunda Parte: Análisis de la coyuntura. Duración máxima: 2 horas. En el cuadro adjunto se proporciona información, para el periodo 2000-2013, de los principales agregados macroeconómicos de la economía de un determinado país. En base a dicha información, responda a las siguientes cuestiones:
\begin{enumerate}[leftmargin=*,itemsep=1pt]
\item Analice la evolución cíclica de dicha economía en el periodo especificado en el cuadro macroeconómico, así como la composición del crecimiento desde la óptica de la oferta y de la demanda.
\item Comente los posibles desequilibrios macroeconómicos que presenta esta economía.
\item Evalúe el comportamiento del mercado de trabajo, señalando las posibles causas de su evolución.
\item Analice la evolución del sector exterior y sus principales determinantes.
\item Evalúe el comportamiento del Sector Público y analice los principales rasgos de la política fiscal.
\item Exponga sus recomendaciones de política económica orientadas al crecimiento económico sostenido a medio y largo plazo.
\end{enumerate}

\paragraph{2013}
\textit{Examen del 30 de noviembre de 2013.} Primer Ejercicio. 2ª parte. Duración máxima: 2 horas. Se permite el uso de calculadora. Se permiten referencias a datos que conozca el opositor, aunque no estén en las tablas adjuntas y siempre y cuando el grueso del análisis se centre en los datos que se proporcionan.
\begin{enumerate}[leftmargin=*,itemsep=1pt]
\item \emph{Cuestión 1}: En el cuadro 1 que se adjunta, se facilita información, referente al periodo 2011-2013, de la contabilidad trimestral de España para ese periodo, en tasas de variación intertrimestral, de acuerdo con los datos ofrecidos por el INE. Sobre la base de esta información, responda a las siguientes cuestiones:
  \par I. Determine las tasas medias intertrimestrales para el PIB para los periodos anuales de 2011 y 2012. Calcule las tasas equivalentes anuales para los mismos periodos.
  \par II. Comente la evolución cíclica de la economía española durante el periodo cubierto en el cuadro adjunto de los componentes de la demanda.
  \par III. Lo mismo para los componentes de la oferta.
  \par IV. Realice un análisis específico de la evolución del sector exterior.
\item \emph{Cuestión 2}: Sobre la base de los datos que se proporcionan en el cuadro 2 relativos a deuda y saldo presupuestario de las Administraciones Públicas españolas, analice la evolución del saldo presupuestario y la sostenibilidad de la deuda pública.
\item \emph{Cuestión 3}: De acuerdo con su diagnóstico sobre esta economía exponga sus principales recomendaciones de política económica orientadas al crecimiento a medio y largo plazo.
\end{enumerate}

\paragraph{2011}
\textit{Examen del 22 de octubre de 2011.} Primer ejercicio. 2ª parte. En los cuadros adjuntos se muestra la evolución de los principales agregados macroeconómicos para España y Alemania, así como las perspectivas de la Comisión Europea para 2011 y 2012. Responda, por favor, a las siguientes cuestiones en un tiempo máximo de 2 horas:
\begin{enumerate}[leftmargin=*,itemsep=1pt]
\item a) Compare el perfil cíclico de la economía española con el de la alemana y las perspectivas para 2011 y 2012, así como la composición del crecimiento del PIB desde la óptica de la demanda. b) Calcule la contribución de la demanda externa neta al crecimiento del PIB en España en 2009.
\item Analice y compare la evolución en ambas economías del empleo en el período 2007-2009.
\item a) Calcule los crecimientos en media anual de la productividad real por ocupado y del coste laboral unitario nominal (CLU) en España y Alemania en 2010. b) Comente la evolución de la competitividad en ambos países utilizando el tipo de cambio efectivo real.
\item Compare la evolución de los tipos de interés reales en España y Alemania y su posible impacto en el crecimiento.
\item a) Según la Comisión Europea, la elasticidad de las exportaciones españolas de bienes y servicios respecto al comercio mundial es de 1,1. ¿Cuál sería el crecimiento de las exportaciones españolas en 2011 si el comercio mundial aumentara un 7,3\% y la competitividad se mantuviera estable? b) Calcule el saldo exterior por cuenta corriente, la capacidad o necesidad de financiación y la tasa de inversión de España en 2010. ¿Qué efectos ha tenido la crisis financiera internacional en el déficit exterior de la economía española?
\item a) Explique cuál fue la orientación de la política fiscal en España y Alemania en el período 2007-2009 y si fue adecuada o no a la situación cíclica dichos países. b) Si el artículo 135 de la Constitución Española fijara un límite para el déficit público estructural (déficit primario cíclicamente ajustado) del 0,4\% del PIB a partir de 2020 y se cumplieran las previsiones de la Comisión Europea hasta 2012, ¿a qué ritmo debería reducirse el déficit estructural cada año?
\end{enumerate}

\end{small}

% ==========================================================================
\section{Expresiones útiles para redactar}\label{anexo:expresiones}

\begin{itemize}
\item Plantillas para no quedarse en blanco; los puntos suspensivos se rellenan con los datos del cuadro. Ninguna sustituye al análisis.
\item \emph{Diagnóstico}:
  \begin{itemize}
  \item «La economía muestra un crecimiento sólido pero desequilibrado, apoyado en la demanda interna y con una necesidad de financiación exterior creciente».
  \item «Parece tratarse de una economía avanzada en recuperación, con la brecha de producción aún negativa, la inflación contenida y un margen fiscal reducido».
  \end{itemize}
\item \emph{Ciclo}:
  \begin{itemize}
  \item «El cuadro recoge más de un ciclo: …».
  \item «La brecha de producción se cierra en … y se abre de nuevo en …» (la brecha \emph{se cierra} o \emph{se abre}; el PIB \emph{se acelera} o \emph{se desacelera}).
  \item «la recuperación tiene forma de uve: …».
  \end{itemize}
\item \emph{Contribuciones}:
  \begin{itemize}
  \item «De los … puntos que crece el PIB, … los aporta la demanda interna y … la externa».
  \item «La demanda externa amortigua la caída del PIB, más por el desplome de las importaciones que por el empuje de las exportaciones».
  \item «La demanda interna es procíclica y la externa, contracíclica».
  \end{itemize}
\item \emph{Prudencia}:
  \begin{itemize}
  \item «Con la información disponible, …»; «esta evolución sería coherente con …, aunque no puede descartarse …»; «a falta de datos sobre …, cabe suponer que …, lo que habría que confirmar».
  \item «Aunque los agentes deciden con el tipo real, la fuerte subida del tipo nominal al final de la serie indica un endurecimiento de la política monetaria».
  \end{itemize}
\item \emph{Interpretación}: «más que …, lo que llama la atención es …»; «este dato es especialmente relevante porque …».
\item \emph{Sector exterior}:
  \begin{itemize}
  \item «Los desequilibrios externos suelen ser síntoma de desequilibrios internos».
  \item «El déficit exterior refleja un exceso de inversión sobre el ahorro, cuyo origen parece estar en el sector …».
  \item «La mejora de la posición de inversión internacional, pese a la necesidad de financiación, se debería a efectos valoración y al crecimiento del PIB nominal».
  \end{itemize}
\item \emph{Políticas}:
  \begin{itemize}
  \item «El tono de la política fiscal, medido por la variación del saldo estructural primario, habría sido expansivo en …».
  \item «Aunque el tipo nominal sube, el tipo real sigue siendo negativo: la política monetaria seguiría siendo acomodaticia».
  \end{itemize}
\item \emph{Recomendaciones}: «parece necesaria una consolidación fiscal, pero gradual, para no lastrar la recuperación»; «a corto plazo …; a medio y largo plazo …».
\item \emph{Cierre}: «en conclusión, la economía afronta … retos; cómo los aborde determinará su crecimiento en los próximos años».
\end{itemize}

\begin{anotaciones}
\emph{Ejemplo de respuesta: un sistema bancario que entra en una crisis bien capitalizado.}

«El sistema financiero parece llegar a la crisis en una posición sólida. En los años previos mejora su \emph{solvencia}: la ratio de capital de nivel 1 sube hasta niveles holgadamente superiores a los mínimos de Basilea~III, y la morosidad baja hasta niveles cercanos al 3~\%, aunque con margen de mejora.

En la crisis, la solvencia resiste. Es probable que hayan ayudado las moratorias de préstamos y las garantías públicas, que retrasan el reconocimiento de la morosidad, y la suspensión del reparto de dividendos, que retiene beneficios como capital \parencite{BCE2020Dividendos,ABE2020Moratorias}; habrá que vigilar la calidad de los préstamos cuando se retiren esas medidas.

La \emph{liquidez} no parece un problema: la política monetaria aporta financiación abundante y barata a los bancos (operaciones a plazo largo en condiciones favorables, compras de activos) y los depósitos crecen.

La \emph{rentabilidad} es el punto débil. Los depósitos aumentan por el ahorro forzoso (no se puede gastar) y precautorio (incertidumbre), pero con tipos muy bajos el margen de intermediación se estrecha, y las provisiones por el deterioro esperado del crédito reducen el beneficio: el ROE cae por debajo del coste del capital \parencite{BCE2020IEF}.»
\end{anotaciones}

\begin{anotaciones}
\emph{Ejemplo de respuesta: estabilidad financiera con tipos de interés muy bajos durante años.}

«Al inicio de la serie la economía necesita reducir su endeudamiento y su riesgo cambiario, y una política de demanda acomodaticia parece haber mitigado el coste del ajuste. Al final de la serie parecen adecuadas 3 líneas de actuación macroprudencial:
\begin{itemize}
\item si la brecha crédito-PIB es negativa, el crédito está por debajo de su tendencia: no procede activar el colchón anticíclico y, si estuviera activado, convendría liberarlo;
\item unos tipos de interés muy bajos durante años, negativos en términos reales, pueden ser un riesgo para la estabilidad financiera (financiación de proyectos de poca calidad, menor rentabilidad de bancos y aseguradoras, precios de los activos inflados): podría justificarse un colchón de riesgo sistémico;
\item si el precio de la vivienda crece a 2 dígitos y más que el VAB de la construcción (más precio que cantidad), herramientas sectoriales que limiten el crédito hipotecario, porque la vivienda es a la vez el principal activo de los hogares, su principal deuda y una gran exposición de los bancos.»
\end{itemize}
\end{anotaciones}

\begin{anotaciones}
\emph{Los fondos europeos en la balanza de pagos.} Las subvenciones del Mecanismo de Recuperación y Resiliencia se anotan en 2 momentos: cuando llegan los fondos, en la cuenta financiera (otra inversión del banco central, frente a un pasivo con la Comisión), y cuando se gastan, como transferencia de la Comisión, en la renta secundaria si financian gasto corriente o en la cuenta de capital si financian inversión. Los préstamos del Mecanismo son deuda frente a las instituciones de la UE (otra inversión de las administraciones públicas) y sus intereses van a la renta primaria \parencite{BdE2021NGEU}. Por eso los fondos pueden explicar saltos de la cuenta de capital y de la capacidad de financiación.
\end{anotaciones}

% ==========================================================================
\section{Lectura de series y hechos estilizados del ciclo}\label{anexo:series}

\paragraph{Conceptos}
\begin{itemize}
\item \emph{Tendencia}: el movimiento persistente a largo plazo de una serie (lineal, si crece a una tasa constante en logaritmos; determinista o estocástica, si la mueven choques aleatorios).
\item \emph{Ciclo}: las oscilaciones alrededor de la tendencia, irregulares y recurrentes, con 2 puntos de giro (pico y valle) y 2 fases (contracción y expansión). Una serie puede alejarse de su tendencia por un \emph{cambio estructural} y no por el ciclo (gráfico~\ref{fig:ciclo-tendencia}) \vertema{3.A.42}.
\item \emph{Estacionariedad}: una serie estacionaria vuelve a su media tras un choque; una no estacionaria, no \vertema{5.B.5}.
\item \emph{Comovimiento con el PIB}, en 5 rasgos:
  \begin{itemize}
  \item \emph{dirección}: procíclica (se mueve con el PIB), contracíclica (al revés) o acíclica (correlación cercana a 0);
  \item \emph{volatilidad}: si oscila más o menos que el PIB;
  \item \emph{adelanto o retraso}: si gira antes o después que el PIB;
  \item \emph{persistencia}: cuánto dura una desviación;
  \item \emph{conformidad}: si estas relaciones se mantienen en toda la muestra (depende de la serie que dé el tribunal).
  \end{itemize}
\end{itemize}

\begin{figure}[H]
\centering
\caption{El ciclo económico: tendencia, fases y brecha de producción}\label{fig:ciclo-tendencia}
\grafico[0.7\textwidth]{ciclo-tendencia}
\fuente{Elaboración propia. Qué hay que ver: el ciclo es la oscilación del PIB alrededor de su tendencia; de pico a valle, contracción (recesión si es intensa); de valle a pico, expansión. La distancia vertical a la tendencia es la brecha de producción.}
\end{figure}

\begin{figure}[H]
\centering
\caption{Comovimientos con el PIB: procíclica, contracíclica y adelantada}\label{fig:comovimientos}
\grafico[0.9\textwidth]{comovimientos}
\fuente{Elaboración propia. Qué hay que ver: si la otra variable se mueve en el mismo sentido que el PIB o en el contrario, si oscila más o menos y si gira antes o después. Los indicadores adelantados sirven para anticipar el ciclo \parencite{ConferenceBoard2026LEI}.}
\end{figure}

\begin{table}[H]
\centering\footnotesize
\caption{Hechos estilizados del ciclo en las economías avanzadas}\label{tab:hechos-ciclo}
\begin{tabular}{>{\raggedright\arraybackslash}p{3.2cm}>{\raggedright\arraybackslash}p{3.6cm}>{\raggedright\arraybackslash}p{3.2cm}>{\raggedright\arraybackslash}p{3.6cm}}
\toprule
\rowcolor{tceeazulpalido}
\textbf{Variable} & \textbf{Dirección} & \textbf{Volatilidad frente al PIB} & \textbf{Adelanto o retraso} \\
\midrule
Consumo privado & procíclico & menor (duraderos, mayor) & contemporáneo \\
Inversión en equipo & procíclica & varias veces mayor & contemporánea o algo retrasada \\
Inversión en vivienda & procíclica & varias veces mayor & adelantada \\
Variación de existencias & procíclica & muy alta & contemporánea \\
Consumo público & acíclico (correlación cercana a 0); procíclico en muchos emergentes & depende del país & sin patrón claro \\
Importaciones & procíclicas & mayor (el comercio, 2-3 veces) & contemporáneas \\
Demanda externa neta & contracíclica & — & — \\
Empleo & procíclico & similar o algo menor & contemporáneo o algo retrasado \\
Tasa de paro & contracíclica & — & retrasada \\
Productividad del trabajo & procíclica (en España, a menudo contracíclica) & menor & algo adelantada \\
Salario real & débilmente procíclico o acíclico & mucho menor & sin patrón claro \\
Inflación & procíclica (contracíclica con choques de oferta) & — & retrasada \\
Tipos de interés a corto & procíclicos & — & retrasados \\
Pendiente de la curva de tipos & — & — & adelantada \\
Bolsa & procíclica & mucho mayor & adelantada \\
\bottomrule
\end{tabular}
\fuente{Elaboración propia a partir de \textcite{StockWatson1999Ciclo}, \textcite{BlinderMaccini1991Existencias}, \textcite{Kaminsky2004Rains} y \textcite{EstrellaMishkin1998}, coherente con \vertema{3.A.42}.}
\end{table}

\paragraph{Gráficos de lectura}
\begin{itemize}
\item Ejemplos de lectura con datos recientes, para practicar el «qué hay que ver».
\end{itemize}

\begin{figure}[H]
\centering
\caption{España: PIB por persona en edad de trabajar, por hora y horas por persona (variación acumulada desde el cuarto trimestre de 2019)}\label{fig:pib-por-pet}
\grafico[0.8\textwidth]{pib-por-pet-espana}
\fuente{\textcite{Eurostat2026namq10gdp,Eurostat2026namq10a10e,Eurostat2026lfsq}; cálculo propio. Qué hay que ver: desde 2019 el PIB por persona en edad de trabajar crece sobre todo porque se trabajan más horas por persona ($+3{,}9~\%$) y poco por la productividad por hora ($+1{,}7~\%$).}
\end{figure}

\begin{figure}[H]
\centering
\caption{Zona del euro: empleo, horas trabajadas, tasa de actividad y tasa de paro}\label{fig:empleo-horas}
\grafico[0.85\textwidth]{empleo-horas-actividad-zona-euro}
\fuente{\textcite{Eurostat2026namq10a10e,Eurostat2026lfsq,Eurostat2026une}. Qué hay que ver: el empleo supera en un 6~\% el nivel de 2019 y las horas solo en un 5~\% (baja la jornada media); el paro, del 6,4~\% a mediados de 2026, está cerca de su mínimo.}
\end{figure}

\begin{figure}[H]
\centering
\caption{Precio del petróleo Brent en dólares y en euros}\label{fig:brent}
\grafico[0.8\textwidth]{brent}
\fuente{\textcite{EIA2026Brent}; en euros con el tipo de referencia del BCE \parencite{BCE2026EXR}. Qué hay que ver: los grandes picos, 2008 (133 dólares), 2011-2012 (125 dólares en marzo de 2012), 2022 (123 dólares; en euros, 116, por la depreciación del euro) y 2026 (117 dólares en abril y 114 en septiembre, por el conflicto de Oriente Medio).}
\end{figure}

\begin{figure}[H]
\centering
\caption{Impulso fiscal en algunas economías avanzadas}\label{fig:impulso-fiscal}
\grafico[0.8\textwidth]{impulso-fiscal}
\fuente{\textcite{FMI2026FMdatos}; impulso $=-\Delta$ saldo primario ajustado por el ciclo; 2026 y 2027, previsiones. Qué hay que ver: positivo es expansivo; 2020 lo es en todos, en 2021-2022 se retira y EE.~UU. vuelve a expandir en 2023.}
\end{figure}

\begin{figure}[H]
\centering
\caption{Efectos inmediatos de la inflación en las cuentas públicas}\label{fig:inflacion-cuentas}
\grafico[0.75\textwidth]{inflacion-cuentas-publicas}
\fuente{Elaboración propia a partir de \textcite{FMI2023FiscalMonitor}. Qué hay que ver: los ingresos y el PIB nominal suben enseguida y el gasto indexado y los intereses, con retraso; a corto plazo mejoran el déficit y la ratio de deuda, y a medio plazo el efecto se invierte.}
\end{figure}

% ==========================================================================
\section{Cálculos tipo}\label{anexo:calculos}

\begin{itemize}
\item Desde 2022 no se piden cálculos, pero el análisis los usa. Siempre: escribir la fórmula, los datos y el resultado con su unidad.
\item \emph{Contribución al crecimiento} de un componente $X$: $\frac{X_{t-1}}{Y_{t-1}}\,\gamma_{X,t}$, en puntos porcentuales (las importaciones, con signo menos). Ejemplo: si el consumo pesa el 55~\% del PIB y crece un 2~\%, contribuye 1,1 puntos.
\item \emph{Tasa media anual a partir de datos trimestrales}: la media de los 4 niveles del año frente a la del año anterior (no la suma de las tasas intertrimestrales). \emph{Tasa anualizada} de una tasa trimestral $g$: $(1+g)^4-1\approx4g$; con $g=0{,}5~\%$, en torno al 2~\% (se preguntó en 2013).
\item \emph{Productividad, CLU y margen}: $\gamma_{\mathrm{PAT}}=\gamma_Y-\gamma_L$; $\gamma_{\mathrm{CLU}}=\gamma_W-\gamma_{\mathrm{PAT}}$; CLU real: $\gamma_{\mathrm{CLU}}-\gamma_P$; margen: $\gamma_{1+\mu}\approx\gamma_P-\gamma_{\mathrm{CLU}}$, con $P$ el deflactor del PIB. Ejemplo: PIB $+2~\%$, empleo $+1{,}5~\%$, remuneración por asalariado $+3~\%$ y deflactor $+2~\%$ dan $\gamma_{\mathrm{PAT}}=0{,}5$, $\gamma_{\mathrm{CLU}}=2{,}5$, CLU reales $+0{,}5$ y margen $-0{,}5$. Si se da la remuneración total y la tasa de asalarización es 1 (2016), $W$ es el salario medio.
\item \emph{Salario real}: $\gamma_W-\pi$, con el IPC si interesa el poder adquisitivo del trabajador y con el deflactor si interesa el coste para la empresa.
\item \emph{Deflactor}: $\gamma_P\approx\gamma^N-\gamma_Y$ (crecimiento nominal menos real).
\item \emph{Tipo de interés real}: $r\approx i-\pi^e$; a falta de expectativas, con la inflación del periodo (diciéndolo); a largo plazo, mejor con el deflactor.
\item \emph{PIB potencial y brecha}: $\gamma_{Y^p}\approx\gamma_Y-\Delta\mathrm{OG}$; brecha prevista: $\mathrm{OG}_t\approx\mathrm{OG}_{t-1}+\gamma_{Y,t}-\gamma_{Y^p,t}$ (2018 y 2021).
\item \emph{Esfuerzo fiscal}: $\Delta b^{pe}$, con $b^{pe}=b-\varepsilon\cdot\mathrm{OG}-$ medidas puntuales $+$ intereses (2011, 2016, 2018, 2019).
\item \emph{Saldo primario estabilizador}: $\frac{i-\gamma^N}{1+\gamma^N}\,d$; con $d=100~\%$, $i=3~\%$ y $\gamma^N=4~\%$, un déficit primario de casi 1 punto ($-0{,}96$).
\item \emph{Efecto valoración de la PIIN}: $\mathrm{OV}_t=\Delta\mathrm{PIIN}_t-\mathrm{CF}_t$ (2016).
\item \emph{Tasas de inversión y de ahorro nacional}: $I/\mathrm{PIB}$ y $S/\mathrm{PIB}=I/\mathrm{PIB}+\mathrm{CC}/\mathrm{PIB}$ (2011, 2019).
\item \emph{Exportaciones con su elasticidad-renta}: $\gamma_X\approx\varepsilon_{X,Y^*}\,\gamma_{Y^*}$ con competitividad constante; en 2011, $1{,}1\times7{,}3~\%\approx8~\%$.
\item \emph{Pendiente de la curva de tipos}: tipo a largo menos tipo a corto (un diferencial, no una fórmula).
\item \emph{Índices}: índice simple $=$ valor/valor del año base $\times100$; cambio de base con una regla de 3; en los índices encadenados la base es móvil. Si el dato viene en porcentaje, no es el índice sino su tasa de variación \vertema{4.A.1}.
\end{itemize}

% ==========================================================================
\section{La programación financiera del FMI}\label{anexo:fmi}

\begin{itemize}
\item \emph{Qué es}: un \emph{marco} (no un modelo) para formular políticas fiscales, monetarias y estructurales que mantengan o recuperen la estabilidad macroeconómica: crecimiento duradero, inflación bajo control, sistema financiero sólido, finanzas públicas sin grandes desajustes y resistencia a los choques \parencite{FMI2022edXFPP}.
\item \emph{De qué depende la estabilidad}: de factores externos (crisis de los socios, catástrofes), de factores propios difíciles de cambiar (desarrollo, demografía) y de las políticas.
\item \emph{Los 2 objetivos}:
  \begin{itemize}
  \item \emph{equilibrio interno}: producción en su potencial con una inflación baja y estable;
  \item \emph{equilibrio externo}: una cuenta corriente que se pueda financiar de forma sostenible, con reservas suficientes.
  \end{itemize}
\item \emph{Qué los rompe}: demanda excesiva o insuficiente; incertidumbre (sistema financiero frágil, finanzas públicas débiles); una combinación de inflación y tipo de cambio que hace caro producir dentro; burbujas en los precios de los activos. Si se agravan, desembocan en una crisis.
\item \emph{Las 4 cuentas y sus vínculos}: la real (producción y gasto), la fiscal (ingresos, gastos y financiación del Estado), la monetaria (balance del banco central y de los bancos) y la exterior (balanza de pagos). Se relacionan por identidades: la necesidad de financiación de la nación es la suma de la de sus sectores (capítulo~\ref{sec:exterior}); el déficit público se financia con deuda interna, externa o con el banco central; la variación de las reservas cambia la base monetaria. Un programa fija metas coherentes para las 4 a la vez \vertema{3.B.22}.
\end{itemize}

% ==========================================================================
\section{Un caso: la economía española}\label{anexo:espana}

\begin{itemize}
\item \emph{Para qué}: aplicar el método a un país conocido, con datos recientes, como ejercicio. La situación y las políticas de España se estudian en \vertema{4.A.29}.
\item \emph{Diagnóstico}: el FMI resumía así su evaluación de 2026: la economía española «siguió creciendo más que la zona del euro pese al impacto adverso de la guerra en Oriente Medio» \parencite{FMI2026ArticuloIV}.
\item \emph{Fortalezas}:
  \begin{itemize}
  \item crecimiento alto para una economía avanzada: 3,5~\% en 2024 y 2,8~\% en 2025 \parencite{FMI2026ArticuloIV}, tras recuperar el nivel de PIB previo a la pandemia en el segundo trimestre de 2022 \parencite{Eurostat2026namq10gdp};
  \item más empleo: la tasa de actividad de 15 a 64 años (74,9~\% en 2025) ya está por encima de la media de la OCDE \parencite{Eurostat2026lfsa,OCDE2026LFS};
  \item una economía abierta (exportaciones más importaciones en torno al 70~\% del PIB) \parencite{Eurostat2026nama10gdp};
  \item una banca rentable (ROE del 16~\% a mediados de 2026) \parencite{BdE2026EstadSupervisoras}.
  \end{itemize}
\item \emph{Debilidades}:
  \begin{itemize}
  \item el paro sigue siendo alto (10,5~\% en 2025) con muy pocas vacantes, una curva de \autor{Beveridge} lejos del origen \parencite{Eurostat2026une,Eurostat2026jvs};
  \item productividad que crece poco: desde 2019 el PIB por persona en edad de trabajar ha subido sobre todo por más horas por persona y poco por la productividad por hora (gráfico~\ref{fig:pib-por-pet});
  \item cuentas públicas con déficit (2,5~\% del PIB en 2025) y deuda alta, ante el gasto del envejecimiento \parencite{FMI2026ArticuloIV,CE2024Envejecimiento};
  \item menos capital de máxima calidad en la banca que la media europea (CET1 del 14~\%, frente al 16~\%) \parencite{BdE2026EstadSupervisoras,ABE2026RiskDashboard}.
  \end{itemize}
\item \emph{Recomendaciones} (tabla~\ref{tab:recomendaciones}): un ajuste fiscal gradual (el FMI proponía unos 0,5 puntos del PIB al año entre 2026 y 2030) \parencite{FMI2026ArticuloIV}; políticas activas de empleo y de vivienda que mejoren el emparejamiento; y reformas que eleven la productividad (capital humano, tamaño empresarial, inversión en I+D).
\end{itemize}

% ==========================================================================
\section{Dónde se estudia cada bloque}\label{anexo:temas}

\begin{itemize}
\item La guía remite en cada indicador al tema donde se estudia a fondo: primero el del tercer ejercicio (la teoría) y después el del cuarto (la aplicación a España o al sector público).
\end{itemize}

\begin{table}[H]
\centering\footnotesize
\caption{Temas del temario que sostienen cada bloque del dictamen}\label{tab:donde-se-estudia}
\begin{tabular}{>{\raggedright\arraybackslash}p{2.8cm}>{\raggedright\arraybackslash}p{6.6cm}>{\raggedright\arraybackslash}p{4.8cm}}
\toprule
\rowcolor{tceeazulpalido}
\textbf{Bloque} & \textbf{Tercer ejercicio} & \textbf{Cuarto ejercicio}\\
\midrule
Perfil cíclico & \vertemacorto{3.A.30}, \vertemacorto{3.A.31}, \vertemacorto{3.A.33}, \vertemacorto{3.A.34}, \vertemacorto{3.A.42}, \vertemacorto{3.A.43}, \vertemacorto{3.A.44}, \vertemacorto{3.A.45}, \vertemacorto{3.B.28}, \vertemacorto{3.B.32} & \vertemacorto{4.A.1}, \vertemacorto{4.A.4}, \vertemacorto{4.A.5}, \vertemacorto{4.A.9}, \vertemacorto{4.A.11}, \vertemacorto{4.A.13}, \vertemacorto{4.A.29}, \vertemacorto{4.B.4}\\
Mercados de factores & \vertemacorto{3.A.13}, \vertemacorto{3.A.17}, \vertemacorto{3.A.25}, \vertemacorto{3.A.26}, \vertemacorto{3.A.27}, \vertemacorto{3.A.28}, \vertemacorto{3.B.18}, \vertemacorto{3.B.23}, \vertemacorto{3.B.24}, \vertemacorto{3.B.26}, \vertemacorto{3.B.27}, \vertemacorto{3.B.44} & \vertemacorto{4.A.2}, \vertemacorto{4.A.11}, \vertemacorto{4.A.16}, \vertemacorto{4.A.18}, \vertemacorto{4.A.19}, \vertemacorto{4.A.20}\\
Competitividad & \vertemacorto{3.A.9}, \vertemacorto{3.A.41}, \vertemacorto{3.B.5}, \vertemacorto{3.B.6}, \vertemacorto{3.B.7}, \vertemacorto{3.B.12}, \vertemacorto{3.B.13}, \vertemacorto{3.B.14}, \vertemacorto{3.B.15}, \vertemacorto{3.B.33}, \vertemacorto{3.B.41} & \vertemacorto{4.A.1}, \vertemacorto{4.A.22}, \vertemacorto{4.A.25}\\
Sector exterior & \vertemacorto{3.B.11}, \vertemacorto{3.B.12}, \vertemacorto{3.B.15}, \vertemacorto{3.B.17}, \vertemacorto{3.B.18}, \vertemacorto{3.B.22}, \vertemacorto{3.B.29}, \vertemacorto{3.B.45} & \vertemacorto{4.A.21}, \vertemacorto{4.A.22}, \vertemacorto{4.A.23}, \vertemacorto{4.A.29}\\
Políticas & \vertemacorto{3.A.32}, \vertemacorto{3.A.35}, \vertemacorto{3.A.36}, \vertemacorto{3.A.37}, \vertemacorto{3.A.38}, \vertemacorto{3.A.39}, \vertemacorto{3.A.40}, \vertemacorto{3.A.41}, \vertemacorto{3.B.19}, \vertemacorto{3.B.43}, \vertemacorto{3.B.45} & \vertemacorto{4.A.15}, \vertemacorto{4.A.16}, \vertemacorto{4.B.1}, \vertemacorto{4.B.4}, \vertemacorto{4.B.5}, \vertemacorto{4.B.7}, \vertemacorto{4.B.21}, \vertemacorto{4.B.23}, \vertemacorto{4.B.26}\\
Recomendaciones & \vertemacorto{3.A.15}, \vertemacorto{3.A.20}, \vertemacorto{3.A.23}, \vertemacorto{3.B.22} & \vertemacorto{4.A.7}, \vertemacorto{4.A.15}, \vertemacorto{4.A.16}, \vertemacorto{4.A.24}, \vertemacorto{4.A.25}\\
Lectura de series & \vertemacorto{3.A.42}, \vertemacorto{5.B.5} & \vertemacorto{4.A.1}\\
\bottomrule
\end{tabular}
\fuente{Elaboración propia a partir de los títulos del temario vigente.}
\end{table}

\bibliografia

\end{document}
